\documentclass[UTF8,12pt]{ctexart}
\usepackage[a4paper,margin=2.3cm]{geometry}
\usepackage{amsmath,amssymb,amsthm,mathtools}
\usepackage{enumitem}
\usepackage{hyperref}

\usepackage{newunicodechar}
\newcommand{\umath}[1]{\ensuremath{#1}}
\newunicodechar{∶}{\umath{:}}
\newunicodechar{∓}{\umath{\mp}}
\newunicodechar{⧵}{\umath{\backslash}}
\newunicodechar{⃗}{}
\newunicodechar{̸}{\umath{/}}
\newunicodechar{∈}{\umath{\in}}
\newunicodechar{∉}{\umath{\notin}}
\newunicodechar{∋}{\umath{\ni}}
\newunicodechar{∀}{\umath{\forall}}
\newunicodechar{∃}{\umath{\exists}}
\newunicodechar{∅}{\umath{\varnothing}}
\newunicodechar{≤}{\umath{\le}}
\newunicodechar{≥}{\umath{\ge}}
\newunicodechar{⩽}{\umath{\le}}
\newunicodechar{⩾}{\umath{\ge}}
\newunicodechar{≠}{\umath{\ne}}
\newunicodechar{≡}{\umath{\equiv}}
\newunicodechar{≈}{\umath{\approx}}
\newunicodechar{≃}{\umath{\simeq}}
\newunicodechar{≅}{\umath{\cong}}
\newunicodechar{∼}{\umath{\sim}}
\newunicodechar{⊂}{\umath{\subset}}
\newunicodechar{⊆}{\umath{\subseteq}}
\newunicodechar{⊕}{\umath{\oplus}}
\newunicodechar{⊗}{\umath{\otimes}}
\newunicodechar{⊥}{\umath{\perp}}
\newunicodechar{∩}{\umath{\cap}}
\newunicodechar{∪}{\umath{\cup}}
\newunicodechar{∧}{\umath{\wedge}}
\newunicodechar{∣}{\umath{\mid}}
\newunicodechar{∤}{\umath{\nmid}}
\newunicodechar{∥}{\umath{\parallel}}
\newunicodechar{⋅}{\umath{\cdot}}
\newunicodechar{⋯}{\umath{\cdots}}
\newunicodechar{∑}{\umath{\sum}}
\newunicodechar{∏}{\umath{\prod}}
\newunicodechar{∫}{\umath{\int}}
\newunicodechar{∂}{\umath{\partial}}
\newunicodechar{∇}{\umath{\nabla}}
\newunicodechar{∆}{\umath{\Delta}}
\newunicodechar{⇒}{\umath{\Rightarrow}}
\newunicodechar{⇐}{\umath{\Leftarrow}}
\newunicodechar{↔}{\umath{\leftrightarrow}}
\newunicodechar{⟹}{\umath{\Longrightarrow}}
\newunicodechar{⟶}{\umath{\longrightarrow}}
\newunicodechar{⟨}{\umath{\langle}}
\newunicodechar{⟩}{\umath{\rangle}}
\newunicodechar{⌊}{\umath{\lfloor}}
\newunicodechar{⌋}{\umath{\rfloor}}
\newunicodechar{′}{\umath{'}}
\newunicodechar{″}{\umath{''}}
\newunicodechar{‴}{\umath{'''}}
\newunicodechar{ℎ}{\umath{h}}
\newunicodechar{ℓ}{\umath{\ell}}
\newunicodechar{ℙ}{\umath{\mathbb P}}
\newunicodechar{ℬ}{\umath{\mathcal B}}
\newunicodechar{ℱ}{\umath{\mathcal F}}
\newunicodechar{α}{\umath{\alpha}}
\newunicodechar{β}{\umath{\beta}}
\newunicodechar{γ}{\umath{\gamma}}
\newunicodechar{δ}{\umath{\delta}}
\newunicodechar{ε}{\umath{\varepsilon}}
\newunicodechar{ϵ}{\umath{\epsilon}}
\newunicodechar{ζ}{\umath{\zeta}}
\newunicodechar{η}{\umath{\eta}}
\newunicodechar{θ}{\umath{\theta}}
\newunicodechar{ι}{\umath{\iota}}
\newunicodechar{κ}{\umath{\kappa}}
\newunicodechar{λ}{\umath{\lambda}}
\newunicodechar{ν}{\umath{\nu}}
\newunicodechar{ξ}{\umath{\xi}}
\newunicodechar{π}{\umath{\pi}}
\newunicodechar{ρ}{\umath{\rho}}
\newunicodechar{σ}{\umath{\sigma}}
\newunicodechar{τ}{\umath{\tau}}
\newunicodechar{φ}{\umath{\varphi}}
\newunicodechar{ϕ}{\umath{\phi}}
\newunicodechar{χ}{\umath{\chi}}
\newunicodechar{ψ}{\umath{\psi}}
\newunicodechar{ω}{\umath{\omega}}
\newunicodechar{𝛼}{\umath{\alpha}}
\newunicodechar{𝛽}{\umath{\beta}}
\newunicodechar{𝛾}{\umath{\gamma}}
\newunicodechar{𝛿}{\umath{\delta}}
\newunicodechar{𝜁}{\umath{\zeta}}
\newunicodechar{𝜃}{\umath{\theta}}
\newunicodechar{𝜆}{\umath{\lambda}}
\newunicodechar{𝜇}{\umath{\mu}}
\newunicodechar{𝜈}{\umath{\nu}}
\newunicodechar{𝜋}{\umath{\pi}}
\newunicodechar{𝜍}{\umath{\varsigma}}
\newunicodechar{𝜎}{\umath{\sigma}}
\newunicodechar{𝜏}{\umath{\tau}}
\newunicodechar{𝜑}{\umath{\varphi}}
\newunicodechar{𝜒}{\umath{\chi}}
\newunicodechar{𝜓}{\umath{\psi}}
\newunicodechar{𝜔}{\umath{\omega}}
\newunicodechar{𝜕}{\umath{\partial}}
\newunicodechar{𝜽}{\umath{\boldsymbol\theta}}
\newunicodechar{𝐀}{\umath{\mathbf A}}
\newunicodechar{𝐁}{\umath{\mathbf B}}
\newunicodechar{𝐂}{\umath{\mathbb C}}
\newunicodechar{𝐄}{\umath{\mathbf E}}
\newunicodechar{𝐅}{\umath{\mathbf F}}
\newunicodechar{𝐇}{\umath{\mathbf H}}
\newunicodechar{𝐈}{\umath{\mathbf I}}
\newunicodechar{𝐍}{\umath{\mathbb N}}
\newunicodechar{𝐐}{\umath{\mathbb Q}}
\newunicodechar{𝐑}{\umath{\mathbb R}}
\newunicodechar{𝐗}{\umath{\mathbf X}}
\newunicodechar{𝐘}{\umath{\mathbf Y}}
\newunicodechar{𝐙}{\umath{\mathbb Z}}
\newunicodechar{𝐝}{\umath{\mathbf d}}
\newunicodechar{𝐞}{\umath{\mathbf e}}
\newunicodechar{𝐦}{\umath{\mathbf m}}
\newunicodechar{𝐱}{\umath{\mathbf x}}
\newunicodechar{𝐴}{\umath{A}}
\newunicodechar{𝐵}{\umath{B}}
\newunicodechar{𝐶}{\umath{C}}
\newunicodechar{𝐸}{\umath{E}}
\newunicodechar{𝐹}{\umath{F}}
\newunicodechar{𝐻}{\umath{H}}
\newunicodechar{𝐼}{\umath{I}}
\newunicodechar{𝐾}{\umath{K}}
\newunicodechar{𝐿}{\umath{L}}
\newunicodechar{𝑀}{\umath{M}}
\newunicodechar{𝑁}{\umath{N}}
\newunicodechar{𝑂}{\umath{O}}
\newunicodechar{𝑃}{\umath{P}}
\newunicodechar{𝑄}{\umath{Q}}
\newunicodechar{𝑅}{\umath{R}}
\newunicodechar{𝑆}{\umath{S}}
\newunicodechar{𝑇}{\umath{T}}
\newunicodechar{𝑈}{\umath{U}}
\newunicodechar{𝑉}{\umath{V}}
\newunicodechar{𝑊}{\umath{W}}
\newunicodechar{𝑋}{\umath{X}}
\newunicodechar{𝑌}{\umath{Y}}
\newunicodechar{𝑎}{\umath{a}}
\newunicodechar{𝑏}{\umath{b}}
\newunicodechar{𝑐}{\umath{c}}
\newunicodechar{𝑑}{\umath{d}}
\newunicodechar{𝑒}{\umath{e}}
\newunicodechar{𝑓}{\umath{f}}
\newunicodechar{𝑔}{\umath{g}}
\newunicodechar{𝑖}{\umath{i}}
\newunicodechar{𝑗}{\umath{j}}
\newunicodechar{𝑘}{\umath{k}}
\newunicodechar{𝑙}{\umath{l}}
\newunicodechar{𝑚}{\umath{m}}
\newunicodechar{𝑛}{\umath{n}}
\newunicodechar{𝑝}{\umath{p}}
\newunicodechar{𝑞}{\umath{q}}
\newunicodechar{𝑟}{\umath{r}}
\newunicodechar{𝑠}{\umath{s}}
\newunicodechar{𝑡}{\umath{t}}
\newunicodechar{𝑢}{\umath{u}}
\newunicodechar{𝑣}{\umath{v}}
\newunicodechar{𝑤}{\umath{w}}
\newunicodechar{𝑥}{\umath{x}}
\newunicodechar{𝑦}{\umath{y}}
\newunicodechar{𝑧}{\umath{z}}
\newunicodechar{𝒞}{\umath{\mathcal C}}
\newunicodechar{𝒩}{\umath{\mathcal N}}
\newunicodechar{𝒪}{\umath{\mathcal O}}
\newunicodechar{𝔞}{\umath{\mathfrak a}}
\newunicodechar{𝔤}{\umath{\mathfrak g}}
\newunicodechar{𝔩}{\umath{\mathfrak l}}
\newunicodechar{𝔬}{\umath{\mathfrak o}}
\newunicodechar{𝔭}{\umath{\mathfrak p}}
\newunicodechar{𝔮}{\umath{\mathfrak q}}
\newunicodechar{𝔰}{\umath{\mathfrak s}}
\newunicodechar{𝔼}{\umath{\mathbb E}}
\newunicodechar{}{}
\newunicodechar{}{}
\newunicodechar{}{}
\newunicodechar{}{}
\newunicodechar{}{}
\newunicodechar{}{}
\newunicodechar{}{}
\newunicodechar{}{}
\newunicodechar{}{}
\newunicodechar{}{}
\newunicodechar{}{}
\newunicodechar{}{}
\newunicodechar{}{}
\newunicodechar{}{}
\newunicodechar{}{}
\newunicodechar{}{}

\hypersetup{colorlinks=true,linkcolor=blue,urlcolor=blue}

\setlength{\parindent}{2em}
\setlength{\parskip}{0.35em}
\setlist[enumerate]{leftmargin=2.2em,itemsep=0.35em,topsep=0.35em}

\newcommand{\C}{\mathbb C}
\newcommand{\R}{\mathbb R}
\newcommand{\Q}{\mathbb Q}
\newcommand{\N}{\mathbb N}
\newcommand{\Z}{\mathbb Z}
\newcommand{\D}{\mathbb D}
\newcommand{\supp}{\operatorname{supp}}
\newcommand{\dist}{\operatorname{dist}}
\newcommand{\Arg}{\operatorname{Arg}}
\newcommand{\esssup}{\operatorname*{ess\,sup}}

\newcounter{problem}
\newenvironment{problem}[1][]{%
  \refstepcounter{problem}
  \par\bigskip
  \noindent{\large\bfseries 题目 \theproblem\if\relax\detokenize{#1}\relax\else\quad #1\fi}\par
  \smallskip
  \noindent\ignorespaces
}{\par}
\newenvironment{solution}{%
  \par\medskip
  \noindent{\bfseries 解答}\quad\ignorespaces
}{\hfill$\square$\par\medskip}

\title{丘成桐数学竞赛几何与拓扑真题}
\author{}
\date{}

\begin{document}
\maketitle
\tableofcontents
\newpage

\section{个人赛题}

\subsection{2010 年个人赛：几何与拓扑}

原题要求：从以下 6 道题中选择 5 道作答。

\begin{problem}
设
\[
D^*=\{(x,y)\in\R^2:0<x^2+y^2<1\}
\]
是 Euclidean 平面中的去心单位圆盘。设 \(g\) 是 \(D^*\) 上完备且常曲率为 \(-1\) 的 Riemann 度量。求点
\[
(e^{-2\pi},0),\qquad (-e^{-\pi},0)
\]
在该度量下的距离。
\end{problem}

\begin{solution}
去心圆盘上的完备常曲率 \(-1\) 度量是 cusp 双曲度量。在坐标 \(z=re^{i\theta}\) 中，它可写为
\[
ds^2=\frac{|dz|^2}{|z|^2(\log |z|)^2}.
\]
令
\[
t=-\log r>0.
\]
则
\[
\frac{dr^2+r^2d\theta^2}{r^2(\log r)^2}
=\frac{dt^2+d\theta^2}{t^2}.
\]
这就是上半平面模型中坐标 \(\tau=\theta+it\) 的双曲度量。

两点对应为
\[
(e^{-2\pi},0):\quad \theta=0,\ t=2\pi,
\]
\[
(-e^{-\pi},0):\quad \theta=\pi,\ t=\pi.
\]
上半平面中双曲距离公式为
\[
\cosh d
=1+\frac{(\Delta\theta)^2+(\Delta t)^2}{2t_1t_2}.
\]
代入 \(\Delta\theta=\pi\)、\(\Delta t=\pi\)、\(t_1=2\pi\)、\(t_2=\pi\)，得
\[
\cosh d
=1+\frac{\pi^2+\pi^2}{2(2\pi)(\pi)}
=1+\frac12=\frac32.
\]
因此
\[
d=\operatorname{arcosh}\frac32
=\log\frac{3+\sqrt5}{2}.
\]
\end{solution}

\begin{problem}
证明：\(\R^n\) 中每个闭超曲面都有一点，其第二基本形式正定。
\end{problem}

\begin{solution}
设闭超曲面为 \(M\subset\R^n\)。取包含 \(M\) 的最小闭球 \(\overline{B_R(a)}\)。由于 \(M\) 紧，存在 \(p\in M\cap\partial B_R(a)\)。在该点，函数
\[
\rho(x)=|x-a|^2
\]
限制到 \(M\) 上取得最大值。于是对任意切向量 \(v\in T_pM\)，有
\[
\operatorname{Hess}_M\rho(v,v)\le0.
\]
另一方面，Euclidean Hessian 满足
\[
\operatorname{Hess}_{\R^n}\rho(v,v)=2|v|^2.
\]
设外法向取为
\[
\nu=\frac{p-a}{|p-a|}.
\]
Hessian 分解给出
\[
\operatorname{Hess}_M\rho(v,v)
=2|v|^2+2(p-a)\cdot II(v,v).
\]
在支撑球外法向约定下，第二基本形式满足
\[
II_\nu(v,v)\ge \frac1R|v|^2
\]
或按相反法向为负定。选择内法向或外法向的符号后，即可得到某点第二基本形式正定。几何直观是：在最小包围球的接触点，超曲面位于球的一侧，因此其法曲率至少为球的法曲率。
\end{solution}

\begin{problem}
证明实射影空间 \(\mathbb{RP}^n\) 可定向当且仅当 \(n\) 为奇数。
\end{problem}

\begin{solution}
\(\mathbb{RP}^n\) 是 \(S^n\) 对 antipodal 映射
\[
A:S^n\to S^n,\qquad A(x)=-x
\]
的商。覆盖 \(S^n\to\mathbb{RP}^n\) 的 deck 变换为 \(A\)。

商空间可定向当且仅当 deck 变换保持 \(S^n\) 的定向。映射 \(A\) 是 \(\R^{n+1}\) 上线性映射 \(-I\) 对单位球面的限制。其在球面上的 degree 为
\[
\deg A=(-1)^{n+1}.
\]
因此 \(A\) 保定向当且仅当
\[
(-1)^{n+1}=1,
\]
即 \(n\) 为奇数。故 \(\mathbb{RP}^n\) 可定向当且仅当 \(n\) 为奇数。
\end{solution}

\begin{problem}
设 \(\pi:M_1\to M_2\) 是从一个连通可微流形到另一个可微流形的 \(C^\infty\) 映射。假设对每个 \(p\in M_1\)，微分
\[
\pi_*:T_pM_1\to T_{\pi(p)}M_2
\]
都是向量空间同构。
\begin{enumerate}[label=(\alph*)]
\item 题面要求证明若 \(M_1\) 连通，则 \(\pi\) 是覆盖映射；
\item 给出 \(M_2\) 紧但 \(\pi:M_1\to M_2\) 不是覆盖映射、同时满足 \(\pi_*\) 同构性质的例子。
\end{enumerate}
\end{problem}

\begin{solution}
按题面原样，(a) 不成立。微分处处同构只说明 \(\pi\) 是局部微分同胚。局部微分同胚若再满足适当性、或 \(M_1\) 紧且 \(\pi\) 满射到连通 \(M_2\)，才可推出覆盖映射。

反例也回答(b)。取
\[
M_1=(0,1),\qquad M_2=S^1,
\]
定义
\[
\pi(t)=e^{2\pi i t}.
\]
则
\[
\pi'(t)=2\pi i e^{2\pi it}\ne0,
\]
所以微分处处同构。但 \(\pi\) 的像是 \(S^1\setminus\{1\}\)，不是满射，因此不是覆盖映射；而 \(M_2=S^1\) 是紧流形。

若补充假设 \(\pi\) 是 proper 且 \(M_2\) 连通，则证明如下：局部微分同胚说明每个点有均匀小邻域被各个原像邻域微分同胚覆盖；proper 保证纤维离散且紧，从而有限。于是每个点都有 evenly covered 邻域，\(\pi\) 是覆盖映射。
\end{solution}

\begin{problem}
设 \(\Sigma_g\) 为 genus \(g\) 的闭可定向曲面。证明若 \(g>1\)，则 \(\Sigma_g\) 是 \(\Sigma_2\) 的覆盖空间。
\end{problem}

\begin{solution}
闭可定向曲面的 Euler 示性数为
\[
\chi(\Sigma_g)=2-2g.
\]
若 \(\Sigma_g\to\Sigma_2\) 是 \(d\)-重覆盖，则
\[
\chi(\Sigma_g)=d\,\chi(\Sigma_2).
\]
由于
\[
\chi(\Sigma_2)=2-4=-2,
\]
需要
\[
2-2g=-2d,
\qquad d=g-1.
\]
对每个 \(d\ge1\)，genus \(2\) 曲面的基本群有指数 \(d\) 子群；对应的连通覆盖曲面 Euler 示性数为
\[
d(-2)=-2d.
\]
该覆盖曲面为闭可定向曲面，其 genus \(h\) 满足
\[
2-2h=-2d,
\]
所以
\[
h=d+1=g.
\]
因此当 \(g>1\) 时，取 \(d=g-1\) 的连通覆盖即可得到 \(\Sigma_g\to\Sigma_2\)。
\end{solution}

\begin{problem}
设 \(M\) 是光滑四维流形。辛形式是闭 \(2\)-形式 \(\omega\)，满足 \(\omega\wedge\omega\) 是处处非零的 \(4\)-形式。
\begin{enumerate}[label=(\alph*)]
\item 在 \(\R^4\) 上构造一个辛形式；
\item 证明 \(S^4\) 上不存在辛形式。
\end{enumerate}
\end{problem}

\begin{solution}
\begin{enumerate}[label=(\alph*)]
\item 在坐标 \((x_1,y_1,x_2,y_2)\) 上取
\[
\omega=dx_1\wedge dy_1+dx_2\wedge dy_2.
\]
它显然闭，并且
\[
\omega\wedge\omega
=2\,dx_1\wedge dy_1\wedge dx_2\wedge dy_2
\]
处处非零。

\item 若 \(S^4\) 上存在辛形式 \(\omega\)，则 \(\omega\) 闭，且 \(\omega\wedge\omega\) 是体积形式。因此
\[
\int_{S^4}\omega\wedge\omega\ne0.
\]
但
\[
H^2(S^4;\R)=0,
\]
所以闭 \(2\)-形式 \(\omega\) 必为恰当形式：\(\omega=d\alpha\)。于是
\[
\omega\wedge\omega=d(\alpha\wedge\omega).
\]
由 Stokes 定理，
\[
\int_{S^4}\omega\wedge\omega=0,
\]
矛盾。因此 \(S^4\) 上无辛形式。
\end{enumerate}
\end{solution}

\subsection{2011 年个人赛：几何与拓扑}

原题要求：从以下 6 道题中选择 5 道作答。

\begin{problem}
设 \(M\) 是闭光滑 \(n\) 维流形。
\begin{enumerate}[label=(\alph*)]
\item 是否总存在光滑映射 \(f:M\to S^n\)，使得 \(f\) 是本质的，即 \(f\) 不同伦于常值映射？请说明理由。
\item 若把 \(S^n\) 换成 \(n\) 维环面 \(T^n\)，同样的问题是否成立？
\end{enumerate}
\end{problem}

\begin{solution}
\begin{enumerate}[label=(\alph*)]
\item 答案是否定的。取 \(M=\mathbb{RP}^2\)，这是闭光滑 \(2\) 维流形。因为 \(S^2\) 是单连通且 \(2\) 维 CW 复形上的映射到 \(S^2\) 的同伦类由顶维上同调类决定，更具体地说
\[
[\mathbb{RP}^2,S^2]\cong H^2(\mathbb{RP}^2;\mathbb Z).
\]
而
\[
H^2(\mathbb{RP}^2;\mathbb Z)=0.
\]
因此每个连续映射 \(\mathbb{RP}^2\to S^2\) 都同伦于常值映射；光滑映射当然也是如此。故并非每个闭光滑 \(n\) 维流形都存在到 \(S^n\) 的本质光滑映射。

为说明上述判别的来源，可把 \(\mathbb{RP}^2\) 看成 \(2\) 维 CW 复形。由于 \(S^2\) 是 \(1\)-连通的，映射在 \(1\)-骨架上没有本质障碍；唯一可能的障碍和分类数据在 \(H^2(\mathbb{RP}^2;\pi_2(S^2))=H^2(\mathbb{RP}^2;\mathbb Z)\) 中。该群为零，所以不存在非平凡同伦类。

\item 答案也是否定的。取 \(M=S^n\)，其中 \(n\ge2\)。环面 \(T^n\) 是 Eilenberg--MacLane 空间 \(K(\mathbb Z^n,1)\)，其高阶同伦群为零：
\[
\pi_k(T^n)=0,\qquad k\ge2.
\]
任意映射 \(f:S^n\to T^n\) 的同伦类由 \(\pi_n(T^n)\) 中的元素给出，而 \(\pi_n(T^n)=0\)。故每个 \(S^n\to T^n\) 的映射都同伦于常值映射。于是也不能总保证存在到 \(T^n\) 的本质映射。
\end{enumerate}
\end{solution}

\begin{problem}
设 \((X,d)\) 是紧度量空间，\(f:X\to X\) 满足
\[
d(f(x),f(y))=d(x,y),\qquad x,y\in X.
\]
证明 \(f\) 是满射。
\end{problem}

\begin{solution}
由距离保持性可知，若 \(f(x)=f(y)\)，则
\[
d(x,y)=d(f(x),f(y))=0,
\]
所以 \(x=y\)。因此 \(f\) 是单射，并且 \(f\) 连续。

反设 \(f\) 不是满射。由于 \(X\) 紧，\(f(X)\) 也是紧集，从而是闭集。可取
\[
y\in X\setminus f(X).
\]
于是点 \(y\) 到闭集 \(f(X)\) 的距离为正：
\[
\delta=\dist(y,f(X))>0.
\]
考虑轨道
\[
y,\ f(y),\ f^2(y),\ldots .
\]
对任意整数 \(j>i\ge0\)，由距离保持性反复使用得
\[
d(f^i(y),f^j(y))
=d(y,f^{j-i}(y)).
\]
而 \(f^{j-i}(y)\in f(X)\)，所以
\[
d(y,f^{j-i}(y))\ge \delta.
\]
因此序列 \(\{f^m(y)\}_{m\ge0}\) 中任意两个不同项的距离都至少为 \(\delta\)，它不可能有 Cauchy 子列。

但紧度量空间中的任意序列都有收敛子列，收敛子列必为 Cauchy 子列，矛盾。因此 \(f\) 必为满射。
\end{solution}

\begin{problem}
设 \(C_1,C_2\) 是 \(\R^3\) 中两个相互链环的圆。证明 \(C_1\) 在 \(\R^3\setminus C_2\) 中不能同伦于一点。
\end{problem}

\begin{solution}
把 \(C_2\) 看成标准圆
\[
C_2=\{(x,y,0):x^2+y^2=1\}.
\]
它的补空间 \(\R^3\setminus C_2\) 的基本群中有一个由穿过 \(C_2\) 的小圆给出的生成元。等价地，可以用链环数刻画：若闭曲线 \(\gamma\subset \R^3\setminus C_2\)，则 \(\operatorname{lk}(\gamma,C_2)\) 在同伦中保持不变。

具体说明如下。取以 \(C_2\) 为边界的圆盘
\[
D=\{(x,y,0):x^2+y^2\le1\}.
\]
若闭曲线 \(\gamma\) 与 \(D\) 横截，则代数交数
\[
\gamma\cdot D
\]
等于 \(\operatorname{lk}(\gamma,C_2)\)。当 \(\gamma\) 在 \(\R^3\setminus C_2\) 中作同伦时，同伦轨迹不会穿过 \(C_2=\partial D\)，所以该代数交数只能在成对产生或抵消的横截交点中改变，总数保持不变。因此链环数是 \(\R^3\setminus C_2\) 中闭曲线同伦类的不变量。

因为 \(C_1\) 与 \(C_2\) 是两个相互链环的圆，按定义
\[
\operatorname{lk}(C_1,C_2)=\pm1.
\]
而常值曲线与 \(C_2\) 的链环数为 \(0\)。若 \(C_1\) 在 \(\R^3\setminus C_2\) 中能同伦于一点，则链环数必须从 \(\pm1\) 变成 \(0\)，这与同伦不变性矛盾。因此 \(C_1\) 不能在 \(\R^3\setminus C_2\) 中缩为一点。
\end{solution}

\begin{problem}
设
\[
M=\R^2/\Z^2
\]
为二维环面，\(L\) 是 \(\R^2\) 中直线 \(3x=7y\)，并设
\[
S=\pi(L)\subset M,
\]
其中 \(\pi:\R^2\to M\) 是自然投影。求 \(M\) 上一个代表 \(S\) 的 Poincare 对偶的微分形式。
\end{problem}

\begin{solution}
直线 \(3x=7y\) 的方向向量可取为
\[
(7,3).
\]
因为 \(\gcd(7,3)=1\)，它在环面中的像 \(S\) 是一条嵌入闭曲线，其同调类为
\[
[S]=7[a]+3[b]\in H_1(T^2;\mathbb Z),
\]
其中 \([a]\) 表示 \(x\)-方向基本圆，\([b]\) 表示 \(y\)-方向基本圆。

在取定环面方向 \(dx\wedge dy\) 后，若闭曲线 \(C\) 的同调类为
\[
[C]=m[a]+n[b],
\]
则 \(C\) 与 \(S\) 的代数交数为
\[
C\cdot S=\det
\begin{pmatrix}
m&7\\
n&3
\end{pmatrix}
=3m-7n.
\]
因此 Poincare 对偶 \(\operatorname{PD}[S]\in H^1(T^2;\mathbb R)\) 应满足
\[
\int_C \operatorname{PD}[S]=3m-7n.
\]
常系数闭 \(1\)-形式
\[
\alpha=3\,dx-7\,dy
\]
正满足这一点，因为
\[
\int_C \alpha=3m-7n.
\]
所以
\[
\boxed{\alpha=3\,dx-7\,dy}
\]
代表 \(S\) 的 Poincare 对偶。注意 \(\alpha(7\partial_x+3\partial_y)=21-21=0\)，这也与它作为该闭曲线对偶形式的方向关系一致。
\end{solution}

\begin{problem}
\(\R^3\) 中正则曲线 \(C\) 称为 Bertrand 曲线，如果存在一个从 \(C\) 到另一条不同正则曲线 \(D\) 的微分同胚 \(f:C\to D\)，使得对任意 \(x\in C\)，曲线 \(C\) 在 \(x\) 处的主法线与曲线 \(D\) 在 \(f(x)\) 处的主法线相同。这里 \(N_xC\) 表示曲线 \(C\) 过 \(x\) 的主法线，\(T_xC\) 表示 \(C\) 在 \(x\) 处的切线。证明：
\begin{enumerate}[label=(\alph*)]
\item 距离 \(|x-f(x)|\) 为常数，并且两条切线 \(T_xC\) 与 \(T_{f(x)}D\) 的夹角也为常数；
\item 若 \(C\) 的曲率 \(k\) 与挠率 \(\tau\) 处处非零，则存在常数 \(\lambda,\mu\)，使得
\[
\lambda k+\mu\tau=1.
\]
\end{enumerate}
\end{problem}

\begin{solution}
用弧长 \(s\) 参数化 \(C\)，记
\[
\alpha(s)\in C,\qquad T=\alpha'(s),\qquad N,\qquad B
\]
为 Frenet 标架，并记曲率、挠率为 \(k,\tau\)。由于两条曲线在对应点有同一条主法线，存在函数 \(\lambda(s)\)，使得
\[
\beta(s):=f(\alpha(s))=\alpha(s)+\lambda(s)N(s).
\]
这里 \(\beta\) 未必以自身弧长为参数。

对上式求导，并用 Frenet 公式
\[
T'=kN,\qquad N'=-kT+\tau B,\qquad B'=-\tau N,
\]
得
\[
\beta'
=(1-\lambda k)T+\lambda' N+\lambda\tau B.
\]
曲线 \(D\) 在 \(\beta(s)\) 处的切向量 \(\beta'(s)\) 必与它的主法线 \(N\) 垂直；而 \(T,N,B\) 是正交标架，所以 \(\beta'\) 的 \(N\)-分量必须为零，即
\[
\lambda'(s)=0.
\]
因此 \(\lambda\) 是常数，距离
\[
|\beta(s)-\alpha(s)|=|\lambda|
\]
也是常数。

现在有
\[
\beta'=(1-\lambda k)T+\lambda\tau B.
\]
设 \(\theta(s)\) 为 \(C\) 与 \(D\) 对应切线的有向夹角，则可写成
\[
\frac{\beta'}{|\beta'|}
=\cos\theta\,T+\sin\theta\,B.
\]
对右端关于 \(s\) 求导：
\[
\frac{d}{ds}(\cos\theta\,T+\sin\theta\,B)
=-\theta'\sin\theta\,T+(k\cos\theta-\tau\sin\theta)N+\theta'\cos\theta\,B.
\]
另一方面，曲线 \(D\) 的单位切向量沿自身弧长的导数必须沿其主法线方向，而该主法线仍为 \(N\)。因此上式中的 \(T\)-分量和 \(B\)-分量都必须为零。于是
\[
\theta'\sin\theta=0,\qquad \theta'\cos\theta=0,
\]
从而 \(\theta'=0\)。故两条对应切线的夹角为常数。

最后比较
\[
\beta'=(1-\lambda k)T+\lambda\tau B
\]
与
\[
\beta'=|\beta'|(\cos\theta\,T+\sin\theta\,B).
\]
因为 \(C\) 与 \(D\) 是不同曲线，\(\lambda\ne0\)。由两边 \(T\)-分量和 \(B\)-分量得
\[
\tan\theta=\frac{\lambda\tau}{1-\lambda k}.
\]
夹角 \(\theta\) 是常数。若 \(\cos\theta\ne0\)，则
\[
1-\lambda k=\lambda\tau\cot\theta,
\]
即
\[
\lambda k+\lambda\cot\theta\,\tau=1.
\]
令 \(\mu=\lambda\cot\theta\)，便得
\[
\lambda k+\mu\tau=1.
\]
若 \(\cos\theta=0\)，则 \(1-\lambda k=0\)，即 \(\lambda k=1\)，这也属于上式中取 \(\mu=0\) 的情形。命题得证。
\end{solution}

\begin{problem}
设 \(M\) 是如下闭曲面：把半径为 \(r\) 的小圆沿 \(\R^3\) 中一条嵌入闭曲线 \(C\) 搬运一周，使圆心沿 \(C\) 运动，且每一时刻小圆都位于 \(C\) 的法平面中。证明
\[
\int_M H^2\,d\sigma\ge 2\pi^2,
\]
且等号成立当且仅当 \(C\) 是半径为 \(\sqrt2\,r\) 的圆。这里 \(H\) 是 \(M\) 的平均曲率，\(d\sigma\) 是 \(M\) 的面积元。
\end{problem}

\begin{solution}
设 \(C\) 以弧长 \(s\) 参数化为 \(\gamma(s)\)，曲率为 \(\kappa(s)\)。取沿 \(C\) 的 Bishop 正交法标架 \(e_1(s),e_2(s)\)，使
\[
e_1'(s)=-\kappa_1(s)T(s),\qquad e_2'(s)=-\kappa_2(s)T(s),
\]
其中 \(T=\gamma'\)，且 \(\kappa(s)^2=\kappa_1(s)^2+\kappa_2(s)^2\)。管面的参数化为
\[
X(s,\theta)=\gamma(s)+r(\cos\theta\,e_1(s)+\sin\theta\,e_2(s)).
\]
记
\[
q(s,\theta)=\kappa_1(s)\cos\theta+\kappa_2(s)\sin\theta.
\]
直接求导得
\[
X_s=(1-rq)T,\qquad
X_\theta=r(-\sin\theta\,e_1+\cos\theta\,e_2),
\]
所以
\[
d\sigma=r(1-rq)\,ds\,d\theta,
\]
其中管面正则且取 \(r\) 足够小时有 \(1-rq>0\)。

管面的两个主曲率分别为
\[
\frac1r,\qquad -\frac{q}{1-rq}
\]
（整体符号取决于法向选择，但 \(H^2\) 不受影响）。按平均曲率为两个主曲率平均值的约定，
\[
H=\frac12\left(\frac1r-\frac{q}{1-rq}\right)
=\frac{1-2rq}{2r(1-rq)}.
\]
于是
\[
\int_M H^2\,d\sigma
=\int \int_0^{2\pi}
\frac{(1-2rq)^2}{4r^2(1-rq)^2}\,r(1-rq)\,d\theta\,ds.
\]
对固定的 \(s\)，可把法平面旋转，使
\[
q=\kappa(s)\cos\theta.
\]
于是需要计算
\[
I(a)=\int_0^{2\pi}\frac{(1-2a\cos\theta)^2}{1-a\cos\theta}\,d\theta,
\qquad a=r\kappa(s).
\]
由恒等变形
\[
\frac{(1-2a\cos\theta)^2}{1-a\cos\theta}
=-4a\cos\theta+\frac{1}{1-a\cos\theta}
\]
以及
\[
\int_0^{2\pi}\frac{d\theta}{1-a\cos\theta}
=\frac{2\pi}{\sqrt{1-a^2}}\qquad (|a|<1),
\]
得到
\[
I(a)=\frac{2\pi}{\sqrt{1-a^2}}.
\]
代回即得管面的 Willmore 积分公式
\[
\int_M H^2\,d\sigma
=\frac{\pi}{2r}\int_C\frac{ds}{\sqrt{1-r^2\kappa(s)^2}}.
\]

因为函数
\[
\frac1{\sqrt{1-u^2}}\ge 2u,\qquad 0\le u<1,
\]
其等号当且仅当 \(u=1/\sqrt2\)，所以
\[
\int_M H^2\,d\sigma
\ge \frac{\pi}{2r}\int_C 2r\kappa(s)\,ds
=\pi\int_C\kappa(s)\,ds.
\]
闭空间曲线的 Fenchel 定理给出总曲率估计
\[
\int_C\kappa(s)\,ds\ge 2\pi,
\]
且等号当且仅当 \(C\) 是平面凸曲线；在这里进一步结合前一个等号条件可知
\[
r\kappa(s)=\frac1{\sqrt2}
\]
处处成立。因此
\[
\int_M H^2\,d\sigma\ge 2\pi^2.
\]

若等号成立，则 Fenchel 定理取等号，\(C\) 为平面凸闭曲线；同时 \(\kappa\) 为常数 \(1/(\sqrt2\,r)\)。平面中曲率为正常数的闭曲线只能是圆，半径为
\[
\frac1\kappa=\sqrt2\,r.
\]
反过来，若 \(C\) 是半径 \(\sqrt2\,r\) 的圆，则 \(r\kappa=1/\sqrt2\)，标准管面公式中每一步均取等号，故
\[
\int_M H^2\,d\sigma=2\pi^2.
\]
命题得证。
\end{solution}

\subsection{2012 年个人赛：几何与拓扑}

原题要求：从以下 6 道题中选择 5 道作答；若作答超过 5 题，则取其中最高的 5 题计分。

\begin{problem}
证明
\[
\pi_3(S^2)\ne0.
\]
\end{problem}

\begin{solution}
考虑 Hopf 映射
\[
h:S^3\to S^2.
\]
把
\[
S^3=\{(z_1,z_2)\in \C^2:|z_1|^2+|z_2|^2=1\}
\]
看成单位球面，定义
\[
h(z_1,z_2)=\bigl(2z_1\overline z_2,\ |z_1|^2-|z_2|^2\bigr)\in \C\times\R\cong\R^3.
\]
因为
\[
|2z_1\overline z_2|^2+\bigl(|z_1|^2-|z_2|^2\bigr)^2
=(|z_1|^2+|z_2|^2)^2=1,
\]
所以 \(h\) 的像在 \(S^2\) 中。

若 \(\pi_3(S^2)=0\)，则 \(h\) 同伦于常值映射。取 \(S^2\) 上的面积形式 \(\omega\)，归一化为
\[
\int_{S^2}\omega=1.
\]
因为 \(H^2(S^3;\R)=0\)，存在 \(1\)-形式 \(\alpha\) 使
\[
d\alpha=h^*\omega.
\]
Hopf 不变量定义为
\[
H(h)=\int_{S^3}\alpha\wedge d\alpha.
\]
若换取另一个原始形式 \(\alpha+\eta\)，其中 \(d\eta=0\)，则 \(S^3\) 上闭 \(1\)-形式恰当，\(\eta=d\varphi\)，于是
\[
\int_{S^3}d\varphi\wedge d\alpha
=\int_{S^3}d(\varphi\,d\alpha)=0,
\]
所以 \(H(h)\) 定义良好。同伦不改变该积分，因此同伦于常值映射时 Hopf 不变量应为 \(0\)。

另一方面，Hopf 映射的两个正则值的原像是 \(S^3\) 中相互链环一次的两条圆，故
\[
H(h)=\operatorname{lk}(h^{-1}(q_1),h^{-1}(q_2))=1.
\]
这与 \(H(h)=0\) 矛盾。因此 \(h\) 给出 \(\pi_3(S^2)\) 中的非零元素，故
\[
\pi_3(S^2)\ne0.
\]
\end{solution}

\begin{problem}
设 \(M\) 是 \(n\) 维光滑流形，\(X_1,\ldots,X_k\) 是开集 \(U\subset M\) 上处处线性无关的 \(k\) 个光滑向量场，并且
\[
[X_i,X_j]=0,\qquad 1\le i,j\le k.
\]
证明：对任意 \(p\in U\)，存在坐标图 \((V,y^i)\)，其中 \(p\in V\subset U\)，坐标为 \(y^1,\ldots,y^n\)，使得在 \(V\) 上
\[
X_i=\frac{\partial}{\partial y^i},\qquad 1\le i\le k.
\]
\end{problem}

\begin{solution}
先在 \(p\) 附近选取一个 \((n-k)\) 维子流形 \(S\)，使得
\[
T_pM=\operatorname{span}\{X_1(p),\ldots,X_k(p)\}\oplus T_pS.
\]
缩小 \(S\) 后，上述直和关系在 \(S\) 附近仍成立。设 \(\Phi_i^t\) 为 \(X_i\) 的局部流。因为
\[
[X_i,X_j]=0,
\]
这些局部流两两交换：
\[
\Phi_i^s\circ\Phi_j^t=\Phi_j^t\circ\Phi_i^s
\]
在共同定义域上成立。

取 \(S\) 上局部坐标 \(z^{k+1},\ldots,z^n\)。定义映射
\[
F(t^1,\ldots,t^k,z)
=\Phi_1^{t^1}\circ\cdots\circ\Phi_k^{t^k}(z),
\]
其中 \(z\in S\) 且所有参数足够小。在 \((0,\ldots,0,p)\) 处，\(F\) 的微分把 \(\partial/\partial t^i\) 送到 \(X_i(p)\)，把 \(T_pS\) 同胚地送入补空间。故 \(dF\) 是同构。由反函数定理，缩小定义域后，\(F\) 是到某邻域 \(V\) 的微分同胚。

令 \(y^i=t^i\)（\(1\le i\le k\)），并令 \(y^{k+1},\ldots,y^n\) 为 \(S\) 上的坐标 \(z^{k+1},\ldots,z^n\) 经 \(F\) 推到 \(V\) 上所得坐标。由于流两两交换，
\[
F(t^1,\ldots,t^i+s,\ldots,t^k,z)
=\Phi_i^s(F(t^1,\ldots,t^k,z)).
\]
对 \(s\) 在 \(0\) 处求导，得到
\[
\frac{\partial F}{\partial t^i}=X_i\circ F.
\]
因此在坐标 \(y\) 中，
\[
X_i=\frac{\partial}{\partial y^i},\qquad 1\le i\le k.
\]
\end{solution}

\begin{problem}
证明：\(\mathbb{CP}^2\) 的任意自同胚都是保定向的。
\end{problem}

\begin{solution}
设 \(f:\mathbb{CP}^2\to\mathbb{CP}^2\) 是同胚。复射影平面的上同调环为
\[
H^*(\mathbb{CP}^2;\mathbb Z)\cong \mathbb Z[a]/(a^3),
\qquad \deg a=2,
\]
其中 \(a^2\) 是 \(H^4(\mathbb{CP}^2;\mathbb Z)\cong\mathbb Z\) 的正生成元，对应复流形的自然定向。

因为
\[
H^2(\mathbb{CP}^2;\mathbb Z)\cong\mathbb Z,
\]
同胚诱导的同构必满足
\[
f^*a=\varepsilon a,\qquad \varepsilon=\pm1.
\]
于是
\[
f^*(a^2)=(f^*a)^2=(\varepsilon a)^2=a^2.
\]
这说明 \(f\) 在顶维上同调 \(H^4\) 上诱导恒等映射，所以 \(f\) 的度为 \(+1\)。度为 \(+1\) 等价于保定向，因此 \(f\) 必保定向。
\end{solution}

\begin{problem}
证明如下等周不等式：设 \(C\) 是 Euclidean 平面上一条简单、光滑、闭曲线，长度为 \(L\)。证明 \(C\) 围成的面积 \(A\) 满足
\[
A\le \frac{L^2}{4\pi},
\]
且等号成立当且仅当 \(C\) 是圆。
\end{problem}

\begin{solution}
用参数 \(t\in[0,2\pi]\) 按常速参数化闭曲线：
\[
\gamma(t)=(x(t),y(t)),\qquad |\gamma'(t)|=\frac{L}{2\pi}.
\]
平移坐标原点不改变面积和长度，故可假设
\[
\int_0^{2\pi}\gamma(t)\,dt=0.
\]
由 Green 公式，带符号面积满足
\[
A=\frac12\int_0^{2\pi}\bigl(xy'-yx'\bigr)\,dt
=\frac12\int_0^{2\pi}\langle J\gamma,\gamma'\rangle\,dt,
\]
其中 \(J(x,y)=(-y,x)\)。取绝对值并用 Cauchy--Schwarz 不等式，
\[
A\le \frac12
\left(\int_0^{2\pi}|\gamma|^2\,dt\right)^{1/2}
\left(\int_0^{2\pi}|\gamma'|^2\,dt\right)^{1/2}.
\]
由于 \(\gamma\) 的平均值为 \(0\)，Wirtinger 不等式给出
\[
\int_0^{2\pi}|\gamma|^2\,dt
\le
\int_0^{2\pi}|\gamma'|^2\,dt.
\]
因此
\[
A\le \frac12\int_0^{2\pi}|\gamma'|^2\,dt
=\frac12\cdot 2\pi\left(\frac{L}{2\pi}\right)^2
=\frac{L^2}{4\pi}.
\]

若等号成立，则 Cauchy--Schwarz 与 Wirtinger 不等式都取等号。Wirtinger 取等号说明 \(\gamma\) 只有一阶 Fourier 模式：
\[
\gamma(t)=u\cos t+v\sin t
\]
其中 \(u,v\in\R^2\)。Cauchy--Schwarz 取等号说明 \(\gamma'\) 与 \(J\gamma\) 处处成比例；结合常速参数化可得该比例为常数，因而 \(\gamma\) 是圆的参数方程。反过来，半径为 \(R\) 的圆有
\[
A=\pi R^2,\qquad L=2\pi R,
\]
故
\[
A=\frac{L^2}{4\pi}.
\]
所以等号当且仅当 \(C\) 是圆。
\end{solution}

\begin{problem}
设 \(x:M\to\R^3\) 是 Euclidean 三维空间中的闭曲面，其 Gaussian 曲率和平均曲率分别记为 \(K\) 和 \(H\)。证明
\[
\iint_M H\,dA+\iint_M pK\,dA=0,
\qquad
\iint_M pH\,dA+\iint_M dA=0,
\]
其中
\[
p=\vec x\cdot \vec n
\]
是 \(M\) 的支撑函数，\(\vec x\) 是位置向量，\(\vec n\) 是 \(M\) 的单位法向量，\(dA\) 是 \(M\) 的面积元。
\end{problem}

\begin{solution}
采用与题面公式一致的符号约定：形算子 \(S\) 定义为
\[
S(v)=-\nabla_v n,
\]
其主曲率为 \(k_1,k_2\)，并取
\[
H=\frac{k_1+k_2}{2},\qquad K=k_1k_2.
\]
在局部取主方向正交标架 \(e_1,e_2\)。把位置向量分解为切向和法向部分：
\[
x=x^T+p n.
\]
注意 \(x^T=x-pn\)。计算切向散度：
\[
\operatorname{div}_M x^T
=\sum_{i=1}^2\left\langle \nabla_{e_i}(x-pn),e_i\right\rangle.
\]
因为 \(\nabla_{e_i}x=e_i\)，且 \(\nabla_{e_i}n=-k_i e_i\)，得到
\[
\left\langle \nabla_{e_i}(x-pn),e_i\right\rangle
=1+p k_i.
\]
于是
\[
\operatorname{div}_M x^T=2+2pH.
\]
闭曲面无边界，散度积分为零，所以
\[
0=\int_M(2+2pH)\,dA,
\]
即
\[
\int_M pH\,dA+\int_M dA=0.
\]

再设第一 Newton 变换为
\[
P_1=2H\,I-S.
\]
在 Euclidean 三维空间的曲面上，Codazzi 方程蕴含
\[
\operatorname{div}_M P_1=0.
\]
因此
\[
\operatorname{div}_M(P_1x^T)
=\operatorname{tr}\bigl(P_1\circ\nabla x^T\bigr).
\]
在主方向 \(e_i\) 上，\(P_1\) 的特征值分别为 \(k_2,k_1\)，而上面已算出
\[
\left\langle \nabla_{e_i}x^T,e_i\right\rangle=1+p k_i.
\]
故
\[
\operatorname{div}_M(P_1x^T)
=k_2(1+pk_1)+k_1(1+pk_2)
=2H+2pK.
\]
再次对闭曲面积分并用散度定理，得
\[
0=\int_M(2H+2pK)\,dA,
\]
即
\[
\int_M H\,dA+\int_M pK\,dA=0.
\]
两式均得证。
\end{solution}

\begin{problem}
写出 Riemann 流形上正交标架的结构方程。并用结构方程证明下列 Riemann 度量 \(g\) 有常截面曲率 \(c\)：
\[
g=
\frac{\sum_{i=1}^n (dx^i)^2}
{\left[1+\frac c4\sum_{i=1}^n (x^i)^2\right]^2},
\]
其中 \((x^1,\ldots,x^n)\) 是局部坐标系。
\end{problem}

\begin{solution}
设 \(\{\theta_i\}\) 是正交余标架，\(\omega_{ij}\) 是 Levi-Civita 联络形式，\(\Omega_{ij}\) 是曲率形式。结构方程为
\[
d\theta_i=-\sum_j\omega_{ij}\wedge\theta_j,\qquad
\omega_{ij}+\omega_{ji}=0,
\]
以及
\[
\Omega_{ij}=d\omega_{ij}+\sum_k\omega_{ik}\wedge\omega_{kj}
=\frac12\sum_{k,l}R_{ijkl}\theta_k\wedge\theta_l.
\]

对题中度量，记
\[
\rho=1+\frac c4\sum_{i=1}^n(x^i)^2.
\]
取正交余标架
\[
\theta_i=\frac{dx^i}{\rho}.
\]
于是 \(g=\sum_i\theta_i^2\)。直接计算有
\[
d\rho=\frac c2\sum_m x^m dx^m
=\frac c2\,\rho\sum_m x^m\theta_m,
\]
从而
\[
d\theta_i
=-\,\frac{d\rho}{\rho}\wedge\theta_i
=-\frac c2\sum_m x^m\theta_m\wedge\theta_i.
\]
令
\[
\omega_{ij}=\frac c2(x^i\theta_j-x^j\theta_i).
\]
显然 \(\omega_{ij}+\omega_{ji}=0\)，并且
\[
-\sum_j\omega_{ij}\wedge\theta_j
=-\sum_j\frac c2(x^i\theta_j-x^j\theta_i)\wedge\theta_j
=-\frac c2\sum_j x^j\theta_j\wedge\theta_i
=d\theta_i.
\]
故这些 \(\omega_{ij}\) 正是 Levi-Civita 联络形式。

下面计算曲率形式。由
\[
d(x^i\theta_j)=dx^i\wedge\theta_j+x^i\,d\theta_j
=\rho\,\theta_i\wedge\theta_j
-\frac c2x^i\sum_m x^m\theta_m\wedge\theta_j
\]
以及对应的 \(i,j\) 互换式，代入
\[
d\omega_{ij}+\sum_k\omega_{ik}\wedge\omega_{kj}
\]
并整理含 \(x^m\) 的项，可得它们相互抵消，剩下
\[
\Omega_{ij}=c\,\theta_i\wedge\theta_j.
\]
因此
\[
R_{ijij}=c,\qquad i\ne j,
\]
并且曲率张量具有常曲率张量的形式
\[
R_{ijkl}=c(\delta_{ik}\delta_{jl}-\delta_{il}\delta_{jk}).
\]
所以任意二维截面的截面曲率都等于 \(c\)，该度量有常截面曲率 \(c\)。
\end{solution}

\subsection{2013 年个人赛：几何与拓扑}

原题要求：从以下 6 道题中选择 5 道作答。

\begin{problem}
设 \(X=S^1\times S^1/\{p,q\}\) 是把环面上两个不同点 \(p,q\) 识别为一个点所得的空间。求 \(X\) 的同调群和基本群。
\end{problem}

\begin{solution}
把环面记为 \(T^2\)。在 \(T^2\) 中取一条嵌入弧 \(\alpha\) 连接 \(p\) 与 \(q\)，并且弧内部不经过这两个端点。识别 \(p\sim q\) 后，\(\alpha\) 的像成为一条圆。直观上，识别两个点相当于在原空间上额外楔上一条圆；严格地说，商映射
\[
T^2\vee S^1\longrightarrow X
\]
可构造为同伦等价。因此
\[
X\simeq T^2\vee S^1.
\]

于是约化同调满足楔和公式
\[
\widetilde H_i(X)\cong \widetilde H_i(T^2)\oplus \widetilde H_i(S^1).
\]
由
\[
H_0(T^2)=\mathbb Z,\quad H_1(T^2)=\mathbb Z^2,\quad H_2(T^2)=\mathbb Z
\]
以及
\[
H_0(S^1)=\mathbb Z,\quad H_1(S^1)=\mathbb Z,
\]
得到
\[
H_0(X)\cong\mathbb Z,\qquad
H_1(X)\cong\mathbb Z^3,\qquad
H_2(X)\cong\mathbb Z,
\]
并且
\[
H_i(X)=0,\qquad i\ge3.
\]

基本群由 Seifert--van Kampen 定理给出：
\[
\pi_1(X)\cong \pi_1(T^2)*\pi_1(S^1)
\cong \mathbb Z^2 * \mathbb Z.
\]
若写成生成元与关系，则
\[
\pi_1(X)\cong
\langle a,b,c\mid aba^{-1}b^{-1}=1\rangle,
\]
其中 \(a,b\) 来自环面两个基本方向，\(c\) 来自识别两点后新增的圆。
\end{solution}

\begin{problem}
设 \((X,d)\) 是紧度量空间，\(f:X\to X\) 满足
\[
d(f(x),f(y))=d(x,y),\qquad x,y\in X.
\]
证明 \(f\) 是满射。
\end{problem}

\begin{solution}
这与 2011 年个人赛第 2 题相同。为完整起见重写证明。

距离保持映射必连续且单射。若 \(f\) 不满射，则因 \(X\) 紧，\(f(X)\) 紧且闭。取
\[
y\in X\setminus f(X),
\qquad
\delta=\dist(y,f(X))>0.
\]
对任意 \(j>i\ge0\)，利用 \(f\) 保距离得
\[
d(f^i(y),f^j(y))
=d(y,f^{j-i}(y)).
\]
由于 \(j-i\ge1\)，有 \(f^{j-i}(y)\in f(X)\)，所以
\[
d(f^i(y),f^j(y))\ge\delta.
\]
因此序列
\[
y,\ f(y),\ f^2(y),\ldots
\]
任意两项之间距离至少为 \(\delta\)，不可能存在收敛子列。但紧度量空间中任意序列都有收敛子列，矛盾。故 \(f\) 必为满射。
\end{solution}

\begin{problem}
设 \(M^2\) 是完备正则曲面，\(K\) 为 Gaussian 曲率。设 \(\sigma:[0,\infty)\to M\) 是测地线，并满足
\[
K(\sigma(t))\le f(t),
\]
其中 \(f\) 是 \([0,\infty)\) 上可微函数。证明：方程
\[
u''(t)+f(t)u(t)=0
\]
的任一非零解在 \([0,t_0]\) 上有零点，其中 \(\sigma(t_0)\) 是沿 \(\sigma\) 从 \(\sigma(0)\) 出发的第一个共轭点。
\end{problem}

\begin{solution}
沿曲面测地线 \(\sigma\)，垂直 Jacobi 场可写为
\[
J(t)=j(t)E(t),
\]
其中 \(E(t)\) 是沿 \(\sigma\) 的平行单位法向量，标量函数 \(j\) 满足
\[
j''(t)+K(\sigma(t))j(t)=0.
\]
\(\sigma(t_0)\) 是第一个共轭点，意味着存在非零 Jacobi 场 \(J\)，使
\[
j(0)=0,\qquad j(t_0)=0,
\]
并且可取 \(j(t)>0\) 对 \(0<t<t_0\) 成立。

设 \(u\) 是
\[
u''+fu=0
\]
的非零解。反设 \(u\) 在 \([0,t_0]\) 上没有零点。必要时把 \(u\) 换成 \(-u\)，可设
\[
u(t)>0,\qquad 0\le t\le t_0.
\]
考虑 Wronskian
\[
W(t)=u(t)j'(t)-u'(t)j(t).
\]
直接计算得
\[
W'(t)=u j''-u''j
=u[-K(\sigma(t))j]-[-f(t)u]j
=\bigl(f(t)-K(\sigma(t))\bigr)u(t)j(t)\ge0
\]
在 \(0<t<t_0\) 上成立。因此 \(W\) 单调不减。

又因为 \(j(0)=j(t_0)=0\)，且 \(j>0\) 在中间成立，所以
\[
j'(0)>0,\qquad j'(t_0)<0.
\]
于是
\[
W(0)=u(0)j'(0)>0,
\qquad
W(t_0)=u(t_0)j'(t_0)<0.
\]
这与 \(W\) 单调不减矛盾。因此 \(u\) 在 \([0,t_0]\) 上必有零点。
\end{solution}

\begin{problem}
设 \(g_1,g_2\) 是可微流形 \(M\) 上的 Riemann 度量，\(R_1,R_2\) 分别为它们的 Riemann 曲率张量。假设对任意 \(X,Y\in T_pM\)，都有
\[
R_1(X,Y,Y,X)=R_2(X,Y,Y,X).
\]
证明对任意 \(X,Y,Z,W\in T_pM\)，都有
\[
R_1(X,Y,Z,W)=R_2(X,Y,Z,W).
\]
\end{problem}

\begin{solution}
令
\[
A=R_1-R_2.
\]
则 \(A\) 是一个 algebraic curvature tensor，满足曲率张量的全部代数对称性：
\[
A(X,Y,Z,W)=-A(Y,X,Z,W)=-A(X,Y,W,Z),
\]
\[
A(X,Y,Z,W)=A(Z,W,X,Y),
\]
以及第一 Bianchi 恒等式。题设等价于
\[
A(X,Y,Y,X)=0,\qquad X,Y\in T_pM.
\]

先对 \(Y+Z\) 极化。由
\[
0=A(X,Y+Z,Y+Z,X)
\]
并用 \(A(X,Y,Y,X)=A(X,Z,Z,X)=0\)，得
\[
A(X,Y,Z,X)+A(X,Z,Y,X)=0.
\]
利用最后两个变量反对称与成对对称，可把它写成
\[
A(X,Y,X,Z)=A(X,Z,X,Y).
\]
再对 \(X+W\) 极化，并反复使用上述对称性和 Bianchi 恒等式，得到标准极化公式
\[
6A(X,Y,Z,W)
=Q(X+Z,Y+W)-Q(X+Z,Y-W)
-Q(X-Z,Y+W)+Q(X-Z,Y-W),
\]
其中
\[
Q(U,V)=A(U,V,V,U).
\]
而题设给出 \(Q(U,V)=0\) 对所有 \(U,V\) 成立，因此
\[
A(X,Y,Z,W)=0
\]
对所有 \(X,Y,Z,W\) 成立。故 \(R_1=R_2\) 在点 \(p\) 上作为四线性张量完全相同。
\end{solution}

\begin{problem}
设 \(M^n\) 是偶数维、可定向、具有正截面曲率的 Riemann 流形。设 \(\sigma:[0,l]\to M\) 是闭测地线，即
\[
\sigma(0)=\sigma(l),\qquad \sigma'(0)=\sigma'(l).
\]
证明存在 \(\varepsilon>0\) 和光滑映射
\[
F:[0,l]\times(-\varepsilon,\varepsilon)\to M,
\]
使得 \(F(t,0)=\sigma(t)\)，且对任意固定的 \(s\ne0\)，曲线
\[
\sigma_s(t)=F(t,s)
\]
是长度严格小于 \(\sigma\) 的闭光滑曲线。
\end{problem}

\begin{solution}
沿闭测地线 \(\sigma\) 考虑法向量丛
\[
E(t)=\{\sigma'(t)\}^{\perp}\subset T_{\sigma(t)}M.
\]
由于 \(M\) 可定向且 \(\sigma'(0)=\sigma'(l)\)，沿 \(\sigma\) 的平行移动在法空间 \(E(0)\) 上给出一个保定向正交变换
\[
P:E(0)\to E(0).
\]
因为 \(n\) 为偶数，所以
\[
\dim E(0)=n-1
\]
为奇数。任意奇数维实正交空间上的保定向正交变换必有特征值 \(1\)。故存在非零 \(v\in E(0)\)，使
\[
P(v)=v.
\]
令 \(V(t)\) 为沿 \(\sigma\) 的平行法向量场，满足 \(V(0)=v\)。由 \(P(v)=v\) 得
\[
V(l)=V(0),
\]
所以 \(V\) 是沿闭曲线的闭变分向量场。

定义变分
\[
F(t,s)=\exp_{\sigma(t)}(sV(t)).
\]
当 \(|s|\) 足够小时，\(F\) 光滑，且每条 \(\sigma_s(t)=F(t,s)\) 都是闭曲线。

设 \(L(s)\) 为 \(\sigma_s\) 的长度。因为 \(\sigma\) 是测地线，第一变分给出
\[
L'(0)=0.
\]
闭测地线的第二变分公式为
\[
L''(0)=\frac1l\int_0^l
\left(|\nabla_tV|^2-\langle R(V,\sigma')\sigma',V\rangle\right)\,dt
\]
在 \(|\sigma'|=1\) 且 \(V\perp\sigma'\) 的情形下成立。这里 \(V\) 是平行场，所以 \(\nabla_tV=0\)。又因为截面曲率处处为正，
\[
\langle R(V,\sigma')\sigma',V\rangle
=K(V,\sigma')|V|^2|\sigma'|^2>0.
\]
因此
\[
L''(0)<0.
\]
于是 \(s=0\) 是长度函数的严格局部最大方向，故存在 \(\varepsilon>0\)，使得当 \(0<|s|<\varepsilon\) 时
\[
L(s)<L(0).
\]
这正是所需结论。
\end{solution}

\begin{problem}
设 \((M^2,ds^2)\) 是 \(\R^3\) 中的极小曲面，其中 \(ds^2\) 为 Euclidean 度量限制到 \(M\) 上的第一基本形式。假设其 Gaussian 曲率 \(K\) 为负。记
\[
d\widetilde s^{\,2}=-K\,ds^2
\]
的 Gaussian 曲率为 \(\widetilde K\)。证明
\[
\widetilde K=1.
\]
\end{problem}

\begin{solution}
二维度量在共形变换下的曲率公式如下：若
\[
\widetilde g=e^{2u}g,
\]
则
\[
\widetilde K=e^{-2u}(K-\Delta_g u),
\]
其中 \(\Delta_g\) 取正号约定
\[
\Delta_g=\operatorname{div}\nabla.
\]
本题中
\[
\widetilde g=-K\,g,
\qquad
e^{2u}=-K,
\qquad
u=\frac12\log(-K).
\]

对 \(\R^3\) 中 Gaussian 曲率为负的极小曲面，有 Ricci 方程
\[
\Delta_g\log(-K)=4K.
\]
于是
\[
\Delta_g u=2K.
\]
代入共形曲率公式得
\[
\widetilde K
=(-K)^{-1}(K-2K)
=1.
\]

为说明 Ricci 方程的来源，取局部共形坐标使
\[
ds^2=e^{2\lambda}|dz|^2.
\]
极小曲面的 Hopf 微分
\[
Q\,dz^2
\]
全纯，并且
\[
K=-4|Q|^2e^{-4\lambda}
\]
在相应规范下成立。因为 \(Q\ne0\) 的区域中 \(\log|Q|\) 调和，直接计算
\[
\Delta_g\log(-K)
=e^{-2\lambda}\Delta_0(\log 4+2\log|Q|-4\lambda)
=-4e^{-2\lambda}\Delta_0\lambda.
\]
又共形度量的 Gaussian 曲率满足
\[
K=-e^{-2\lambda}\Delta_0\lambda,
\]
故
\[
\Delta_g\log(-K)=4K.
\]
由于题设 \(K<0\)，无零点问题，Ricci 方程在全局成立。
\end{solution}

\subsection{2014 年个人赛：几何与拓扑}

原题要求：从以下 6 道题中选择 5 道作答。

\begin{problem}
设 \(X\) 是把 \(S^2\) 的赤道 \(S^1\) 上每一点 \(x\) 与其对径点 \(-x\) 识别后所得的商空间。计算同调群 \(H_n(X)\)。再对 \(S^3\) 作同样计算，即把赤道 \(S^2\subset S^3\) 上的对径点识别。
\end{problem}

\begin{solution}
先看 \(S^2\) 的情形。上下两个闭半球在商空间中分别变成一个 \(\mathbb{RP}^2\)，它们沿公共的赤道商空间
\[
S^1/(x\sim -x)\cong \mathbb{RP}^1\cong S^1
\]
粘合。因此
\[
X=\mathbb{RP}^2\cup_{\mathbb{RP}^1}\mathbb{RP}^2.
\]
取 CW 结构：一个 \(0\)-胞腔、一个 \(1\)-胞腔来自 \(\mathbb{RP}^1\)，再有两个 \(2\)-胞腔分别来自上下半球。每个 \(2\)-胞腔的边界映到 \(1\)-骨架的次数为 \(2\)，方向相差一个符号。因此链复形可写为
\[
0\longrightarrow \mathbb Z^2
\xrightarrow{(2,-2)}
\mathbb Z
\longrightarrow \mathbb Z
\longrightarrow0.
\]
于是
\[
H_0(X)\cong\mathbb Z,\qquad
H_1(X)\cong\mathbb Z/2\mathbb Z,\qquad
H_2(X)\cong\mathbb Z,
\]
并且 \(H_n(X)=0\) 对 \(n\ge3\) 成立。

再看 \(S^3\) 的情形。上下两个闭半球 \(D^3\) 的边界按对径点识别后分别成为 \(\mathbb{RP}^3\)，公共部分是
\[
S^2/(x\sim -x)=\mathbb{RP}^2.
\]
因此商空间为
\[
Y=\mathbb{RP}^3\cup_{\mathbb{RP}^2}\mathbb{RP}^3.
\]
用 Mayer--Vietoris 序列，记 \(A=B=\mathbb{RP}^3\)，\(A\cap B=\mathbb{RP}^2\)。已知
\[
H_*(\mathbb{RP}^3):\quad H_0=\mathbb Z,\ H_1=\mathbb Z/2,\ H_2=0,\ H_3=\mathbb Z,
\]
\[
H_*(\mathbb{RP}^2):\quad H_0=\mathbb Z,\ H_1=\mathbb Z/2,\ H_i=0\ (i\ge2).
\]
在 \(H_1\) 上，包含映射 \(\mathbb{RP}^2\hookrightarrow\mathbb{RP}^3\) 诱导同构 \(\mathbb Z/2\to\mathbb Z/2\)。因此
\[
H_3(Y)\cong H_3(A)\oplus H_3(B)\cong\mathbb Z^2,
\]
\[
H_2(Y)=0,
\]
而
\[
H_1(Y)\cong
(\mathbb Z/2\oplus\mathbb Z/2)/\langle(1,1)\rangle
\cong\mathbb Z/2.
\]
再加上连通性，
\[
H_0(Y)\cong\mathbb Z.
\]
其余同调群均为 \(0\)。
\end{solution}

\begin{problem}
设 \(M\to\R^3\) 是图面
\[
z=f(u,v),
\]
其中 \(\{u,v,z\}\) 是 \(\R^3\) 中的直角坐标系。假设 \(M\) 是极小曲面。证明：
\begin{enumerate}[label=(\alph*)]
\item \(M\) 的 Gaussian 曲率 \(K\) 可表示为
\[
K=\Delta\log\left(1+\frac1W\right),
\qquad
W=\sqrt{1+f_u^2+f_v^2},
\]
其中 \(\Delta\) 是关于 \(M\) 上诱导度量的 Laplace 算子；
\item 若 \(f\) 定义在整个 \(uv\)-平面上，则 \(f\) 是线性函数。
\end{enumerate}
\end{problem}

\begin{solution}
\begin{enumerate}[label=(\alph*)]
\item 取图面的向上单位法向量
\[
N=\frac{(-f_u,-f_v,1)}{W}.
\]
因此
\[
N_3=\langle N,e_3\rangle=\frac1W.
\]
对 \(\R^3\) 中的极小曲面，有 Chern 公式
\[
\Delta\log(1+\langle N,a\rangle)=K
\]
对任意常单位向量 \(a\) 成立，只要 \(1+\langle N,a\rangle>0\)。这里取 \(a=e_3\)，由于图面的向上法向满足 \(N_3>0\)，便得
\[
K=\Delta\log\left(1+\frac1W\right).
\]

下面说明该公式。极小曲面的 Gauss 映射
\[
N:M\to S^2
\]
在局部共形坐标下是反全纯或全纯映射，并且第三坐标函数 \(N_3\) 满足
\[
\Delta N_3+|A|^2N_3=0,
\]
其中 \(A\) 是第二基本形式。极小曲面有
\[
|A|^2=-2K.
\]
又有恒等式
\[
|\nabla N_3|^2=-K(1-N_3^2).
\]
于是
\[
\Delta\log(1+N_3)
=\frac{\Delta N_3}{1+N_3}
-\frac{|\nabla N_3|^2}{(1+N_3)^2}
=\frac{2KN_3}{1+N_3}
\frac{K(1-N_3^2)}{(1+N_3)^2}
=K.
\]
代入 \(N_3=1/W\) 即为所求。

\item 因为 \(M\) 是整图，Gauss 映射 \(N\) 的像完全落在开上半球
\[
\{q\in S^2:q_3>0\}.
\]
用从南极的立体投影把该开半球映到单位圆盘，得到一个有界的全纯函数 \(g\)，它就是极小曲面的 Weierstrass Gauss 映射。

整极小图的共形参数域是全平面 \(\C\)。于是 \(g\) 是定义在 \(\C\) 上的有界全纯函数。由 Liouville 定理，\(g\) 为常数。因此 Gauss 映射 \(N\) 为常值，图面的单位法向量处处相同。法向量常值说明切平面处处相同，故 \(M\) 是一个平面。于是
\[
f(u,v)=au+bv+d
\]
为线性函数。
\end{enumerate}
\end{solution}

\begin{problem}
设 \(M=\R^2/\Z^2\) 是二维环面，\(L\) 是 \(\R^2\) 中直线 \(3x=7y\)，并设 \(S=\pi(L)\subset M\)，其中 \(\pi:\R^2\to M\) 是自然投影。求 \(M\) 上一个代表 \(S\) 的 Poincare 对偶的微分形式。
\end{problem}

\begin{solution}
这与 2011 年个人赛第 4 题相同。直线方向向量为 \((7,3)\)，故
\[
[S]=7[a]+3[b]\in H_1(T^2;\mathbb Z).
\]
若闭曲线 \(C\) 的同调类为 \(m[a]+n[b]\)，则
\[
C\cdot S=
\det\begin{pmatrix}m&7\\ n&3\end{pmatrix}
=3m-7n.
\]
因此代表 Poincare 对偶的 \(1\)-形式应在该曲线上积分为 \(3m-7n\)。取
\[
\alpha=3\,dx-7\,dy.
\]
则
\[
\int_C\alpha=3m-7n,
\]
所以
\[
\boxed{3\,dx-7\,dy}
\]
代表 \(\operatorname{PD}[S]\)。
\end{solution}

\begin{problem}
设
\[
p:(\widetilde M,\widetilde g)\to(M,g)
\]
是 Riemannian submersion，即 \(p\) 是浸没，且对每个 \(x\in\widetilde M\)，微分
\[
Dp:\ker(Dp)^\perp\to T_{p(x)}M
\]
是线性等距同构。
\begin{enumerate}[label=(\alph*)]
\item 证明 \(p\) 缩短距离；
\item 若 \((\widetilde M,\widetilde g)\) 完备，则 \((M,g)\) 完备；
\item 举例说明即使 \((M,g)\) 完备，\((\widetilde M,\widetilde g)\) 也未必完备。
\end{enumerate}
\end{problem}

\begin{solution}
\begin{enumerate}[label=(\alph*)]
\item 设 \(\gamma:[a,b]\to\widetilde M\) 是分段光滑曲线。把 \(\gamma'\) 分解为水平和垂直部分：
\[
\gamma'=(\gamma')^H+(\gamma')^V.
\]
由于 \(Dp\) 在水平空间上等距、在垂直空间上为零，
\[
|(p\circ\gamma)'|
=|Dp((\gamma')^H)|
=|(\gamma')^H|
\le|\gamma'|.
\]
积分得
\[
L(p\circ\gamma)\le L(\gamma).
\]
对所有连接两点的曲线取下确界，得到
\[
d_M(p(\tilde x),p(\tilde y))\le d_{\widetilde M}(\tilde x,\tilde y),
\]
故 \(p\) 缩短距离。

\item 证明 \(M\) 测地完备。取任意 \(q\in M\) 和 \(v\in T_qM\)。选 \(\tilde q\in p^{-1}(q)\)，并取唯一水平向量 \(\tilde v\in\ker(Dp)^\perp\)，使
\[
Dp_{\tilde q}(\tilde v)=v.
\]
因为 \(\widetilde M\) 完备，存在定义在整个 \(\R\) 上的测地线 \(\tilde\gamma\)，满足
\[
\tilde\gamma(0)=\tilde q,\qquad \tilde\gamma'(0)=\tilde v.
\]
Riemannian submersion 的基本性质给出：水平测地线的投影是 \(M\) 中的测地线，并且 \(\tilde\gamma\) 保持水平。因此
\[
\gamma=p\circ\tilde\gamma
\]
是 \(M\) 上以 \(q\) 为初点、初速度为 \(v\) 的测地线，且定义在整个 \(\R\) 上。由 Hopf--Rinow 定理，\(M\) 完备。

\item 令
\[
\widetilde M=(0,1)\times S^1,\qquad
\widetilde g=dt^2+d\theta^2,
\]
并令
\[
M=S^1,\qquad g=d\theta^2.
\]
投影
\[
p(t,\theta)=\theta
\]
是 Riemannian submersion，因为水平空间由 \(\partial_\theta\) 张成且 \(Dp\) 在该方向上等距。但 \((0,1)\times S^1\) 在 product metric 下不完备，例如 \(t\to0\) 的曲线有限长却无极限点；而 \(S^1\) 是完备的。
\end{enumerate}
\end{solution}

\begin{problem}
设 \(\Psi:M\to\R^3\) 是紧曲面 \(M\) 到 \(\R^3\) 的等距浸入。证明
\[
\int_M H^2\,d\sigma\ge4\pi,
\]
其中 \(H\) 是平均曲率，\(d\sigma\) 是 \(M\) 的面积元。
\end{problem}

\begin{solution}
这是曲面的 Willmore 不等式。给出标准证明。对任意 \(a\in\R^3\setminus\Psi(M)\)，作反演
\[
I_a(x)=\frac{x-a}{|x-a|^2}.
\]
曲面的 Willmore 能量
\[
W(M)=\int_M H^2\,d\sigma
\]
在 \(\R^3\) 的共形变换下不变。因此反演后的曲面 \(\widehat M=I_a(\Psi(M))\) 满足
\[
\int_{\widehat M}\widehat H^2\,d\widehat\sigma
=\int_M H^2\,d\sigma.
\]

选取 \(a\) 为 \(\Psi(M)\) 围成区域中的一点，或一般地取 \(a\) 趋近于曲面上一点并考虑极限；反演后得到一个完备非紧曲面 \(\widehat M\)，它有一个平面型端。对非紧曲面使用点态不等式
\[
\widehat H^2-\widehat K
=\frac14(\widehat k_1-\widehat k_2)^2\ge0,
\]
并对截去大圆盘后的紧部分使用 Gauss--Bonnet 定理。平面型端的边界项在极限中贡献 \(2\pi\)，而端本身的总转角再贡献 \(2\pi\)。于是得到
\[
\int_{\widehat M}\widehat H^2\,d\widehat\sigma\ge4\pi.
\]
由共形不变性，
\[
\int_M H^2\,d\sigma\ge4\pi.
\]

等价地，这一证明说明了 \(4\pi\) 来自反演后一个平面端的 Gauss--Bonnet 边界贡献；球面取等号，因为半径为 \(R\) 的圆球上
\[
H=\frac1R,\qquad d\sigma=R^2\,d\sigma_{S^2},
\]
从而
\[
\int_{S_R^2}H^2\,d\sigma
=\frac1{R^2}\cdot4\pi R^2=4\pi.
\]
\end{solution}

\begin{problem}
\(S^2\) 的单位切丛是
\[
T^1(S^2)=\{(p,v)\in\R^3:\|p\|=1,\ (p,v)=0,\ \|v\|=1\}.
\]
证明它是 \(S^2\) 的切丛 \(T(S^2)\) 中的光滑子流形，并且
\[
T^1(S^2)\cong\mathbb{RP}^3.
\]
\end{problem}

\begin{solution}
先证明光滑子流形。定义
\[
F:\R^3\times\R^3\to\R^3,\qquad
F(p,v)=\bigl(|p|^2,\ p\cdot v,\ |v|^2\bigr).
\]
则
\[
T^1(S^2)=F^{-1}(1,0,1).
\]
在满足 \(|p|=|v|=1\)、\(p\perp v\) 的点，三个函数的微分分别为
\[
2p\cdot dp,\qquad v\cdot dp+p\cdot dv,\qquad 2v\cdot dv.
\]
它们线性无关，因此 \((1,0,1)\) 是正则值，\(T^1(S^2)\) 是 \(\R^6\) 中的光滑子流形。由于它还满足 \(p\in S^2\)、\(v\in T_pS^2\)，它自然是 \(T(S^2)\) 中的光滑子流形。

构造到 \(SO(3)\) 的微分同胚。给定 \((p,v)\in T^1(S^2)\)，令
\[
w=p\times v.
\]
则 \((p,v,w)\) 是 \(\R^3\) 的一个正向标准正交标架。令矩阵 \(A\) 的三列分别为 \(p,v,w\)，则 \(A\in SO(3)\)。反过来，任意 \(A\in SO(3)\) 的第一列 \(p\) 与第二列 \(v\) 满足
\[
|p|=|v|=1,\qquad p\perp v,
\]
从而给出 \(T^1(S^2)\) 中一点。这两个构造互逆且光滑，所以
\[
T^1(S^2)\cong SO(3).
\]

最后说明
\[
SO(3)\cong\mathbb{RP}^3.
\]
单位四元数群
\[
S^3=\{q\in\mathbb H:|q|=1\}
\]
通过共轭作用在纯虚四元数空间 \(\operatorname{Im}\mathbb H\cong\R^3\) 上：
\[
x\mapsto qxq^{-1}.
\]
这给出满同态 \(S^3\to SO(3)\)，其核为 \(\{\pm1\}\)。因此
\[
SO(3)\cong S^3/\{\pm1\}=\mathbb{RP}^3.
\]
综上，
\[
T^1(S^2)\cong\mathbb{RP}^3.
\]
\end{solution}

\subsection{2015 年个人赛：几何与拓扑}

原题要求：从以下 6 道题中选择 5 道作答。

\begin{problem}
设 \(n,m\) 为正整数。证明球面乘积
\[
S^n\times S^m
\]
的切丛平凡，当且仅当 \(n\) 或 \(m\) 为奇数。
\end{problem}

\begin{solution}
把 \(S^n\subset\R^{n+1}\)、\(S^m\subset\R^{m+1}\) 看作单位球面。球面的法线丛由位置向量给出，是平凡线丛。因此
\[
TS^n\oplus\varepsilon^1\cong\varepsilon^{n+1},
\qquad
TS^m\oplus\varepsilon^1\cong\varepsilon^{m+1}.
\]
在乘积上，
\[
T(S^n\times S^m)
\cong \pi_1^*TS^n\oplus\pi_2^*TS^m.
\]
再加上两个平凡线丛得
\[
T(S^n\times S^m)\oplus\varepsilon^2
\cong\varepsilon^{n+m+2}.
\]
所以该切丛是稳定平凡的。

若 \(n\) 为奇数，则 \(S^n\) 存在一个处处非零切向量场 \(V\)。利用上面的稳定平凡化，可以把两条法向平凡线中的一条与 \(V\) 交换，得到 \(n+m\) 个处处线性无关的切向量场，从而
\[
T(S^n\times S^m)
\]
平凡。若 \(m\) 为奇数同理。

反过来，若 \(T(S^n\times S^m)\) 平凡，则 Euler 类为零，因而 Euler 示性数为零：
\[
\chi(S^n\times S^m)=0.
\]
但
\[
\chi(S^k)=1+(-1)^k,
\]
所以
\[
\chi(S^n\times S^m)
=(1+(-1)^n)(1+(-1)^m).
\]
该乘积为零当且仅当 \(n\) 或 \(m\) 为奇数。故命题成立。
\end{solution}

\begin{problem}
证明不存在紧三维流形 \(M\)，使其边界为实射影平面 \(\mathbb{RP}^2\)。
\end{problem}

\begin{solution}
反设存在紧三维流形 \(M\)，满足
\[
\partial M=\mathbb{RP}^2.
\]
取 \(\mathbb Z/2\) 系数。任意紧奇数维流形与其边界满足
\[
\chi(\partial M)=2\chi(M)
\]
在 \(\mathbb Z/2\) 系数 Betti 数定义的 Euler 示性数下仍成立。这可由 Poincare--Lefschetz 对偶和长正合列推出。

因此 \(\chi(\partial M)\) 必为偶数。但
\[
\chi(\mathbb{RP}^2)=1.
\]
矛盾。所以不存在这样的紧三维流形。

也可从可定向双覆盖看出矛盾：若三维紧流形以 \(\mathbb{RP}^2\) 为唯一边界分量，则其边界 Euler 示性数必须为偶数；而 \(\mathbb{RP}^2\) 的 Euler 示性数为奇数。
\end{solution}

\begin{problem}
设 \(M^n\) 是无边界光滑流形，\(X\) 是 \(M\) 上光滑向量场。若 \(X(p)\ne0\)，证明存在以 \(p\) 为中心的局部坐标图
\[
(U;x^1,\ldots,x^n),
\]
使得在 \(U\) 中
\[
X=\frac{\partial}{\partial x^1}.
\]
\end{problem}

\begin{solution}
这是流盒定理。取过 \(p\) 的 \((n-1)\) 维局部截面 \(S\)，使
\[
T_pM=\R X(p)\oplus T_pS.
\]
令 \(\Phi_t\) 为 \(X\) 的局部流。在 \(S\) 上取局部坐标
\[
(y^2,\ldots,y^n),
\]
并定义
\[
F(t,y^2,\ldots,y^n)=\Phi_t(q(y^2,\ldots,y^n)),
\]
其中 \(q(y)\in S\)。在 \((0,0,\ldots,0)\) 处，\(dF\) 把 \(\partial/\partial t\) 送到 \(X(p)\)，把其余坐标方向送到 \(T_pS\)。这些向量构成 \(T_pM\) 的基，因此 \(dF\) 是同构。

由反函数定理，缩小邻域后，\(F\) 是到某开集 \(U\) 的微分同胚。令
\[
x^1=t,\qquad x^i=y^i\ (2\le i\le n).
\]
由于
\[
F(t+s,y)=\Phi_s(F(t,y)),
\]
对 \(s\) 在 \(0\) 处求导得到
\[
\frac{\partial F}{\partial t}=X\circ F.
\]
因此在坐标 \(x\) 中
\[
X=\frac{\partial}{\partial x^1}.
\]
\end{solution}

\begin{problem}
设 \(M\to\R^3\) 是紧、单连通、闭曲面。证明：若 \(M\) 有常平均曲率，则 \(M\) 是标准球面。
\end{problem}

\begin{solution}
由于 \(M\) 是紧且单连通的闭曲面，由闭曲面分类定理，\(M\) 同胚于 \(S^2\)。取局部共形坐标 \(z\)，第一基本形式写成
\[
I=e^{2\lambda}|dz|^2.
\]
第二基本形式可写为
\[
II=Q\,dz^2+H e^{2\lambda}dz\,d\bar z+\overline Q\,d\bar z^2,
\]
其中
\[
Q\,dz^2
\]
称为 Hopf 二次微分。

Codazzi 方程在共形坐标中给出
\[
Q_{\bar z}=\frac12 e^{2\lambda}H_z.
\]
因为题设平均曲率 \(H\) 为常数，所以
\[
Q_{\bar z}=0.
\]
因此 \(Q\,dz^2\) 是 \(M\) 上的全纯二次微分。

但 \(M\cong S^2\)，而 \(S^2\) 上不存在非零全纯二次微分。事实上全纯二次微分是典范线丛平方 \(K_{S^2}^{\otimes2}\) 的全纯截面，而
\[
\deg K_{S^2}^{\otimes2}=-4<0,
\]
故其全纯截面只能为零。于是
\[
Q\equiv0.
\]
这说明第二基本形式的无迹部分为零，曲面处处全脐，即两个主曲率处处相等。

若公共主曲率为 \(\kappa\)，Codazzi 方程又推出 \(\kappa\) 为常数。若 \(\kappa=0\)，曲面是平面的一部分，不可能紧闭；所以 \(\kappa\ne0\)。设 \(N\) 为单位法向，则
\[
d\left(x+\frac1\kappa N\right)=dx+\frac1\kappa dN=0
\]
在适当法向约定下成立。因此
\[
x+\frac1\kappa N
\]
为常向量，\(M\) 的像是以该常向量为球心、半径 \(|1/\kappa|\) 的圆球。故 \(M\) 是标准球面。
\end{solution}

\begin{problem}
设 \(M\) 是 \(n\) 维紧 Riemann 流形，直径为 \(\pi/c\)，并满足 Ricci 曲率
\[
\operatorname{Ric}\ge (n-1)c^2>0.
\]
证明 \(M\) 等距于 \(\R^{n+1}\) 中半径为 \(1/c\) 的标准 \(n\) 维球面。
\end{problem}

\begin{solution}
由 Bonnet--Myers 定理，若
\[
\operatorname{Ric}\ge(n-1)c^2,
\]
则
\[
\operatorname{diam}(M)\le\frac{\pi}{c}.
\]
题设假设直径正好达到该上界。下面用最大直径刚性定理证明。

取 \(p,q\in M\)，使
\[
d(p,q)=\frac{\pi}{c}.
\]
对任意 \(x\in M\)，三角不等式给出
\[
d(p,x)+d(x,q)\ge d(p,q)=\frac{\pi}{c}.
\]
另一方面，Laplacian 比较定理给出距离函数 \(r_p=d(p,\cdot)\) 在割迹外满足
\[
\Delta r_p\le (n-1)c\cot(cr_p),
\]
类似地
\[
\Delta r_q\le (n-1)c\cot(cr_q).
\]
令
\[
u=r_p+r_q.
\]
在 \(0<r_p,r_q<\pi/c\) 的区域中，利用 \(u\ge\pi/c\) 和 cot 函数的单调性，可由最大值原理推出
\[
r_p(x)+r_q(x)=\frac{\pi}{c}
\]
对所有 \(x\) 成立。于是每个点都位于从 \(p\) 到 \(q\) 的某条最短测地线上。

进一步，Laplacian 比较中的等号处处成立。这迫使沿所有从 \(p\) 出发的径向测地线，所有径向截面曲率都等于 \(c^2\)，并且测地极坐标中的度量为
\[
dr^2+\frac1{c^2}\sin^2(cr)\,g_{S^{n-1}}.
\]
这正是半径 \(1/c\) 的标准球面度量。故
\[
M\cong S^n_{1/c}
\]
等距。
\end{solution}

\begin{problem}
设 \((M,g)\) 是 Riemann 流形，\(p\in M\)。证明在以 \(p\) 为中心的正规坐标中，度量的二阶 Taylor 展开为
\[
g_{ij}(x)=\delta_{ij}-\frac13\sum_{k,l}R_{iklj}x^kx^l+O(|x|^3).
\]
\end{problem}

\begin{solution}
取以 \(p\) 为中心的正规坐标 \(x^1,\ldots,x^n\)。在 \(p\) 点有
\[
g_{ij}(0)=\delta_{ij},\qquad \partial_k g_{ij}(0)=0,
\]
因为 Christoffel 符号在正规坐标中心为零。故二阶展开形如
\[
g_{ij}(x)=\delta_{ij}
\frac12\sum_{k,l}\partial_k\partial_l g_{ij}(0)x^kx^l
O(|x|^3).
\]
需要计算二阶导数。

在正规坐标中，从径向测地线性质可得
\[
\sum_{j,k}\Gamma^i_{jk}(x)x^jx^k=0.
\]
对该恒等式在 \(0\) 处求导，得到
\[
\partial_l\Gamma^i_{jk}(0)+\partial_j\Gamma^i_{kl}(0)+\partial_k\Gamma^i_{lj}(0)=0.
\]
曲率在 \(p\) 点为
\[
R^i{}_{klj}(0)
=\partial_l\Gamma^i_{jk}(0)-\partial_j\Gamma^i_{lk}(0),
\]
因为 \(\Gamma(0)=0\)。结合上面的循环关系，解得
\[
\partial_l\Gamma^i_{jk}(0)
=-\frac13\bigl(R^i{}_{jkl}+R^i{}_{kjl}\bigr)(0).
\]
另一方面，
\[
\Gamma_{ijk}
=\frac12(\partial_jg_{ik}+\partial_kg_{ij}-\partial_i g_{jk}),
\]
且在 \(0\) 点一阶导数为零。由此反解得到
\[
\partial_k\partial_l g_{ij}(0)
=-\frac13\bigl(R_{iklj}+R_{ilkj}\bigr)(0).
\]
代入 Taylor 展开，并利用 \(x^kx^l\) 关于 \(k,l\) 对称，得
\[
g_{ij}(x)
=\delta_{ij}
-\frac13\sum_{k,l}R_{iklj}(0)x^kx^l
O(|x|^3).
\]
这就是所需二阶展开。
\end{solution}

\subsection{2016 年个人赛：几何与拓扑}

原题要求：从以下 6 道题中选择 5 道作答。

\begin{problem}
设 \(M\) 是带边界的紧奇数维流形。证明
\[
\chi(M)=\frac12\chi(\partial M).
\]
\end{problem}

\begin{solution}
设 \(\dim M=m\) 为奇数。考虑 \(M\) 的双倍
\[
DM=M\cup_{\partial M}M,
\]
即沿边界粘合两个 \(M\) 的拷贝。Euler 示性数满足粘合公式
\[
\chi(DM)=\chi(M)+\chi(M)-\chi(\partial M)
=2\chi(M)-\chi(\partial M).
\]
另一方面，\(DM\) 是闭奇数维流形。由 Poincare 对偶，闭奇数维流形的 Euler 示性数为零：
\[
\chi(DM)=0.
\]
因此
\[
2\chi(M)-\chi(\partial M)=0,
\]
即
\[
\chi(M)=\frac12\chi(\partial M).
\]
\end{solution}

\begin{problem}
计算去掉一点的二维环面
\[
T^2-\{p\}
\]
的 de Rham 上同调，其中 \(p\in T^2\)。若
\[
T^2=\R^2/\Z^2
\]
并以 \((x,y)\in\R^2\) 为坐标，则体积形式
\[
\omega=dx\wedge dy
\]
是否恰当？
\end{problem}

\begin{solution}
去掉一点的环面 \(T^2-\{p\}\) 同伦等价于一个 bouquet
\[
S^1\vee S^1.
\]
因此其 de Rham 上同调为
\[
H^0_{\mathrm{dR}}(T^2-\{p\})\cong\R,
\]
\[
H^1_{\mathrm{dR}}(T^2-\{p\})\cong\R^2,
\]
\[
H^k_{\mathrm{dR}}(T^2-\{p\})=0,\qquad k\ge2.
\]
特别地，
\[
H^2_{\mathrm{dR}}(T^2-\{p\})=0.
\]
所以在穿孔环面 \(T^2-\{p\}\) 上，任意闭 \(2\)-形式都是恰当的。因此 \(\omega=dx\wedge dy\) 限制到
\[
T^2-\{p\}
\]
上是恰当形式。

注意若在整个紧环面 \(T^2\) 上考虑，则 \(\omega\) 不是恰当的，因为
\[
\int_{T^2}\omega=1\ne0,
\]
而紧无边界流形上恰当顶维形式的积分为零。题目问的是穿孔环面时，答案为恰当。
\end{solution}

\begin{problem}
设 \(M^n\to\R^{n+1}\) 是闭定向超曲面。\(M^n\) 的第 \(r\) 个平均曲率定义为
\[
H_r=\frac1{\binom nr}
\sum_{i_1<i_2<\cdots<i_r}
\lambda_{i_1}\lambda_{i_2}\cdots\lambda_{i_r},
\qquad 1\le r\le n,
\]
其中 \(\lambda_i\) 是 \(M^n\) 的主曲率。证明：若所有 \(\lambda_i\) 都为正，且某个 \(H_r\) 为常数，则 \(M^n\) 是 \(\R^{n+1}\) 中的超球面。
\end{problem}

\begin{solution}
所有主曲率为正说明 \(M\) 是严格凸闭超曲面。设支撑函数为
\[
p=\langle x,\nu\rangle,
\]
其中 \(\nu\) 是外单位法向。对凸闭超曲面有 Minkowski 积分公式
\[
\int_M H_{j-1}\,dA=\int_M pH_j\,dA,
\qquad 1\le j\le n,
\]
并且 Newton 不等式给出
\[
H_{j-1}H_{j+1}\le H_j^2,
\]
等号当且仅当所有主曲率相等。

若 \(H_r\) 为常数，则 Minkowski 公式给出
\[
\int_M H_{r-1}\,dA=\int_M pH_r\,dA
=H_r\int_M p\,dA.
\]
另一方面 \(j=1\) 的 Minkowski 公式为
\[
\operatorname{Area}(M)=\int_M pH_1\,dA.
\]
结合 Newton--Maclaurin 不等式
\[
H_1\ge H_r^{1/r},
\qquad
H_{r-1}\ge H_r^{(r-1)/r},
\]
并对 Minkowski 公式作积分比较，可知等号必须处处成立；否则严格凸性会使上述积分不等式严格。于是 Newton 不等式处处取等号，因此
\[
\lambda_1=\lambda_2=\cdots=\lambda_n
\]
在每一点成立。

所以 \(M\) 是全脐超曲面。闭的全脐 Euclidean 超曲面只能是超球面：若公共主曲率为 \(\lambda\)，Codazzi 方程给出 \(\nabla\lambda=0\)，于是
\[
\nabla_X\left(x+\frac1\lambda\nu\right)
=X+\frac1\lambda\nabla_X\nu
=X-\frac1\lambda\lambda X=0.
\]
故
\[
x+\frac1\lambda\nu
\]
为常向量，即 \(M\) 是以该常向量为球心、半径 \(1/\lambda\) 的超球面。
\end{solution}

\begin{problem}
陈述并证明可微流形上的截断函数引理。
\end{problem}

\begin{solution}
截断函数引理的一种标准表述如下：设 \(M\) 是可微流形，\(K\subset U\subset M\)，其中 \(K\) 为紧集，\(U\) 为开集。则存在光滑函数
\[
\varphi:M\to[0,1],
\]
使得
\[
\varphi\equiv1\quad\text{在 }K\text{ 的某个邻域上},
\]
并且
\[
\supp\varphi\subset U.
\]

证明如下。对每个 \(x\in K\)，取坐标邻域 \(V_x\)，使其闭包紧且
\[
\overline{V_x}\subset U.
\]
由于 \(K\) 紧，可取有限子覆盖
\[
V_{x_1},\ldots,V_{x_N}.
\]
再取较小开集 \(W_i\)，满足
\[
x_i\in W_i,\qquad \overline{W_i}\subset V_{x_i},
\]
并使这些 \(W_i\) 仍覆盖 \(K\)。

在每个坐标图中构造标准 bump 函数 \(\psi_i\)，满足
\[
0\le\psi_i\le1,\qquad
\psi_i\equiv1\text{ 在 }W_i\text{ 上},\qquad
\supp\psi_i\subset V_{x_i}.
\]
这是由 \(\R^n\) 上函数
\[
\eta(t)=
\begin{cases}
e^{-1/t},&t>0,\\
0,&t\le0
\end{cases}
\]
组合得到的标准光滑截断函数。

定义
\[
\varphi=1-\prod_{i=1}^N(1-\psi_i).
\]
则 \(0\le\varphi\le1\)。在 \(\bigcup_i W_i\) 上至少有一个 \(\psi_i=1\)，所以 \(\varphi=1\)，该开集包含 \(K\)。同时
\[
\supp\varphi\subset \bigcup_{i=1}^N\supp\psi_i\subset U.
\]
故所需截断函数存在。
\end{solution}

\begin{problem}
设 \(M\) 是无边界紧 Riemann 流形。证明若 \(M\) 有正 Ricci 曲率，则
\[
H^1(M,\R)=0.
\]
\end{problem}

\begin{solution}
由 Hodge 定理，每个 de Rham 上同调类都有唯一调和代表。因此只需证明任意调和 \(1\)-形式都为零。

设 \(\alpha\) 是调和 \(1\)-形式，并令 \(X=\alpha^\sharp\) 为其对偶向量场。Bochner 公式给出
\[
\frac12\Delta|\alpha|^2
=|\nabla\alpha|^2+\operatorname{Ric}(X,X),
\]
其中 \(\Delta\) 是非负约定下的 Laplacian。对紧无边界流形积分，左边积分为零：
\[
0=\int_M\left(|\nabla\alpha|^2+\operatorname{Ric}(X,X)\right)\,dV.
\]
由于 Ricci 曲率正定，
\[
\operatorname{Ric}(X,X)\ge0
\]
且等号只有 \(X=0\)。又 \(|\nabla\alpha|^2\ge0\)，所以上式迫使
\[
\nabla\alpha=0,\qquad X=0.
\]
因此 \(\alpha=0\)。故所有调和 \(1\)-形式为零，从而
\[
H^1(M,\R)=0.
\]
\end{solution}

\begin{problem}
设 \(M\) 是 \(S^{n+1}\) 中可定向、闭且嵌入的极小超曲面。记 \(\lambda_1\) 为 \(M\) 上 Laplace--Beltrami 算子的第一特征值。证明
\[
\lambda_1\ge\frac n2.
\]
\end{problem}

\begin{solution}
这是 Choi--Wang 型估计。由于 \(M\subset S^{n+1}\) 是闭嵌入可定向极小超曲面，它把球面分成两个区域 \(\Omega_1,\Omega_2\)，且
\[
\partial\Omega_1=\partial\Omega_2=M.
\]
取 \(M\) 上第一特征函数 \(\varphi\)，满足
\[
\Delta_M\varphi+\lambda_1\varphi=0,
\qquad
\int_M\varphi=0.
\]
令 \(u\) 是 \(\Omega=\Omega_1\) 上的调和延拓：
\[
\Delta_\Omega u=0,\qquad u|_M=\varphi.
\]
在单位球面区域 \(\Omega\) 上 Ricci 曲率为 \(n g\)。Reilly 公式给出
\[
\int_\Omega\bigl((\Delta u)^2-|\nabla^2u|^2-\operatorname{Ric}(\nabla u,\nabla u)\bigr)
=\int_M\bigl(2u_\nu\Delta_Mu+H_{\partial\Omega}u_\nu^2+\operatorname{II}(\nabla_Mu,\nabla_Mu)\bigr).
\]
这里 \(M\) 是极小的，所以 \(H_{\partial\Omega}=0\)。又 \(\Delta u=0\)、\(\operatorname{Ric}=n g\)，得
\[
\int_\Omega\left(|\nabla^2u|^2+n|\nabla u|^2\right)
=\int_M\left(2\lambda_1\varphi u_\nu-\operatorname{II}(\nabla_M\varphi,\nabla_M\varphi)\right)
\]
在适当选择法向符号后成立。

对两个区域 \(\Omega_1,\Omega_2\) 分别应用该公式并相加。边界的第二基本形式项因两侧法向相反而相互抵消，得到
\[
\sum_{\alpha=1}^2
\int_{\Omega_\alpha}\left(|\nabla^2u_\alpha|^2+n|\nabla u_\alpha|^2\right)
=2\lambda_1\int_M\varphi\,(u_{\nu,1}+u_{\nu,2}).
\]
另一方面分部积分给出
\[
\int_{\Omega_\alpha}|\nabla u_\alpha|^2
=\int_M\varphi\,u_{\nu,\alpha}.
\]
相加记
\[
E=\sum_{\alpha=1}^2\int_{\Omega_\alpha}|\nabla u_\alpha|^2
=\int_M\varphi\,(u_{\nu,1}+u_{\nu,2}).
\]
于是
\[
\sum_{\alpha=1}^2
\int_{\Omega_\alpha}|\nabla^2u_\alpha|^2+nE
=2\lambda_1E.
\]
由于第一项非负且 \(E>0\)，立即得到
\[
2\lambda_1E\ge nE,
\]
故
\[
\lambda_1\ge\frac n2.
\]
\end{solution}

\subsection{2017 年个人赛：几何与拓扑}

原题要求：从以下 6 道题中选择 5 道作答。

\begin{problem}
设 \(M\) 是光滑、紧、定向的 \(n\) 维流形。若 Euler 示性数 \(\chi(M)=0\)，证明 \(M\) 存在处处非零的向量场。
\end{problem}

\begin{solution}
对紧定向光滑流形，切丛 Euler 类
\[
e(TM)\in H^n(M;\mathbb Z)
\]
满足
\[
\langle e(TM),[M]\rangle=\chi(M).
\]
一个处处非零向量场等价于切丛存在一个 nowhere-zero section；其唯一障碍正是 Euler 类。

若 \(M\) 连通且闭，则 \(H^n(M;\mathbb Z)\cong\mathbb Z\)，由 \(\chi(M)=0\) 得
\[
e(TM)=0.
\]
因此障碍消失，存在处处非零的切向量场。

也可用 Poincare--Hopf 定理说明。先取一个只有孤立零点的向量场 \(X\)。其零点指标和为
\[
\sum_p\operatorname{ind}_p(X)=\chi(M)=0.
\]
通过局部修改，可以把两个指标相反的零点用一条弧连接并相消；反复进行后可消去全部零点。于是得到处处非零向量场。
\end{solution}

\begin{problem}
设
\[
q_1,q_2:S^2\vee S^2\to S^2
\]
分别为压掉另一个楔和分量的映射。设 \(f:S^2\to S^2\vee S^2\) 满足 \(q_i\circ f:S^2\to S^2\) 的度为 \(d_i\)。计算
\[
(S^2\vee S^2)\cup_f D^3
\]
的整数同调群。
\end{problem}

\begin{solution}
记
\[
X=(S^2\vee S^2)\cup_fD^3.
\]
这是一个 CW 复形，含一个 \(0\)-胞腔、两个 \(2\)-胞腔和一个 \(3\)-胞腔，没有 \(1\)-胞腔。因此链群为
\[
C_3\cong\mathbb Z,\qquad C_2\cong\mathbb Z^2,\qquad C_1=0,\qquad C_0\cong\mathbb Z.
\]
边界映射
\[
\partial_3:C_3\to C_2
\]
由附着映射 \(f\) 在 \(H_2(S^2\vee S^2)\cong\mathbb Z^2\) 上的度决定，即
\[
\partial_3(1)=(d_1,d_2).
\]
而 \(\partial_2=0\)。

因此
\[
H_0(X)\cong\mathbb Z,\qquad H_1(X)=0,
\]
\[
H_3(X)=\ker\partial_3
=
\begin{cases}
\mathbb Z,& d_1=d_2=0,\\
0,& \text{否则},
\end{cases}
\]
并且
\[
H_2(X)=\operatorname{coker}\partial_3
=\mathbb Z^2/\langle(d_1,d_2)\rangle.
\]
令
\[
g=\gcd(d_1,d_2)
\]
且约定 \(g=0\) 当 \(d_1=d_2=0\)。若 \((d_1,d_2)\ne(0,0)\)，则 Smith 标准形给出
\[
H_2(X)\cong\mathbb Z\oplus\mathbb Z/g\mathbb Z.
\]
若 \(d_1=d_2=0\)，则
\[
H_2(X)\cong\mathbb Z^2.
\]
更高维同调群均为零。
\end{solution}

\begin{problem}
设 \(X,Y\) 是光滑流形上的光滑向量场。证明 Lie 导数满足
\[
L_XY=[X,Y].
\]
\end{problem}

\begin{solution}
设 \(\Phi_t\) 是向量场 \(X\) 的局部流。向量场 \(Y\) 沿 \(X\) 的 Lie 导数定义为
\[
(L_XY)_p=\left.\frac{d}{dt}\right|_{t=0}
\bigl((\Phi_{-t})_*Y_{\Phi_t(p)}\bigr).
\]
对任意光滑函数 \(f\)，有
\[
\bigl((\Phi_{-t})_*Y_{\Phi_t(p)}\bigr)f
=Y_{\Phi_t(p)}(f\circ\Phi_{-t}).
\]
对 \(t\) 在 \(0\) 处求导：
\[
(L_XY)(f)
=\left.\frac{d}{dt}\right|_{0}
Y(f\circ\Phi_{-t})\circ\Phi_t.
\]
展开两部分变化，第一部分来自沿 \(\Phi_t\) 的变化，给出 \(X(Yf)\)；第二部分来自 \(f\circ\Phi_{-t}\) 的变化，给出 \(-Y(Xf)\)。因此
\[
(L_XY)(f)=X(Yf)-Y(Xf)=[X,Y]f.
\]
由于向量场由其对所有光滑函数的作用决定，故
\[
L_XY=[X,Y].
\]
\end{solution}

\begin{problem}
陈述并证明 \(\R^3\) 中光滑曲面上一条正则曲线的测地曲率 \(\kappa_g\) 的 Liouville 公式。
\end{problem}

\begin{solution}
Liouville 公式如下。设曲面第一基本形式在正交坐标 \((u,v)\) 中为
\[
ds^2=E\,du^2+G\,dv^2.
\]
令曲线以弧长 \(s\) 参数化，并设 \(\phi\) 是曲线切向量与 \(u\)-坐标曲线方向之间的夹角，则测地曲率满足
\[
\kappa_g
=\frac{d\phi}{ds}
-\frac{E_v}{2E\sqrt G}\cos\phi
+\frac{G_u}{2G\sqrt E}\sin\phi.
\]
等价地，若取正交单位标架
\[
e_1=\frac1{\sqrt E}\partial_u,\qquad e_2=\frac1{\sqrt G}\partial_v,
\]
曲线单位切向量为
\[
T=\cos\phi\,e_1+\sin\phi\,e_2,
\]
则
\[
\kappa_g=\left\langle\nabla_TT,\,-\sin\phi\,e_1+\cos\phi\,e_2\right\rangle.
\]

证明。设
\[
J T=-\sin\phi\,e_1+\cos\phi\,e_2
\]
为把 \(T\) 旋转 \(90^\circ\) 得到的切向单位法向。计算
\[
\nabla_TT
=(-\phi'\sin\phi)e_1+\cos\phi\,\nabla_Te_1
+(\phi'\cos\phi)e_2+\sin\phi\,\nabla_Te_2.
\]
与 \(JT\) 内积得
\[
\kappa_g=\phi'+\cos\phi\,\langle\nabla_Te_1,JT\rangle
\sin\phi\,\langle\nabla_Te_2,JT\rangle.
\]
由正交坐标的 Christoffel 公式可得
\[
\nabla_{e_1}e_1=-\frac{E_v}{2E\sqrt G}e_2,\qquad
\nabla_{e_2}e_1=\frac{G_u}{2G\sqrt E}e_2,
\]
\[
\nabla_{e_1}e_2=\frac{E_v}{2E\sqrt G}e_1,\qquad
\nabla_{e_2}e_2=-\frac{G_u}{2G\sqrt E}e_1.
\]
把
\[
T=\cos\phi\,e_1+\sin\phi\,e_2
\]
代入并整理，得到
\[
\kappa_g
=\phi'
-\frac{E_v}{2E\sqrt G}\cos\phi
+\frac{G_u}{2G\sqrt E}\sin\phi.
\]
这就是 Liouville 公式。
\end{solution}

\begin{problem}
在 Riemann 流形上，令 \(\mathcal F\) 为所有满足
\[
|\nabla f|\le1
\]
的光滑函数 \(f\) 的集合。证明对任意 \(x,y\in M\)，有
\[
d(x,y)=\sup\{|f(x)-f(y)|:f\in\mathcal F\}.
\]
\end{problem}

\begin{solution}
若 \(f\in\mathcal F\)，取任意从 \(x\) 到 \(y\) 的分段光滑曲线 \(\gamma\)，则
\[
|f(x)-f(y)|
\le\int_\gamma |\nabla f|\,ds
\le L(\gamma).
\]
对所有 \(\gamma\) 取下确界，得
\[
|f(x)-f(y)|\le d(x,y).
\]
因此
\[
\sup_{f\in\mathcal F}|f(x)-f(y)|\le d(x,y).
\]

反向不等式需要用距离函数逼近。函数
\[
r_z=d(z,y)
\]
是 \(1\)-Lipschitz 函数，并满足
\[
|r_x-r_y|=d(x,y).
\]
虽然 \(r\) 未必光滑，但在 Riemann 流形上任意 \(1\)-Lipschitz 函数都可被光滑函数一致逼近，并且 Lipschitz 常数可控制到 \(1+\varepsilon\)。因此对任意 \(\varepsilon>0\)，存在光滑 \(h_\varepsilon\)，满足
\[
|\nabla h_\varepsilon|\le1+\varepsilon
\]
且
\[
|h_\varepsilon(x)-h_\varepsilon(y)|
\ge d(x,y)-\varepsilon.
\]
令
\[
f_\varepsilon=\frac{h_\varepsilon}{1+\varepsilon},
\]
则 \(f_\varepsilon\in\mathcal F\)，并且
\[
|f_\varepsilon(x)-f_\varepsilon(y)|
\ge\frac{d(x,y)-\varepsilon}{1+\varepsilon}.
\]
令 \(\varepsilon\to0\)，得到
\[
\sup_{f\in\mathcal F}|f(x)-f(y)|\ge d(x,y).
\]
两边合并即得等式。
\end{solution}

\begin{problem}
设 \(M\) 是单位球面 \(S^{n+p}\) 中 \(n\) 维定向闭极小子流形。记 \(K_M\) 为 \(M\) 的截面曲率。证明若
\[
K_M>\frac{p-1}{2p-1},
\]
则 \(M\) 是大球面 \(S^n\)。
\end{problem}

\begin{solution}
设 \(h\) 是 \(M\subset S^{n+p}\) 的第二基本形式，记
\[
S=|h|^2.
\]
Gauss 方程给出，对任意正交切向量 \(e_i,e_j\)，
\[
K_{ij}=1+\sum_{\alpha=1}^p
\bigl(h^\alpha_{ii}h^\alpha_{jj}-(h^\alpha_{ij})^2\bigr).
\]
题设的截面曲率 pinching 意味着第二基本形式不能太大。具体地，由 Berger 估计与 Gauss 方程可推出
\[
S<\frac{np}{2p-1}.
\]

另一方面，极小子流形在单位球面中的 Simons 恒等式给出
\[
\frac12\Delta S
=|\nabla h|^2+nS-Q(h),
\]
其中代数项 \(Q(h)\) 满足 pinching 估计
\[
Q(h)\le\left(2-\frac1p\right)S^2.
\]
于是
\[
\frac12\Delta S
\ge |\nabla h|^2+nS-\left(2-\frac1p\right)S^2
=|\nabla h|^2+S\left(n-\frac{2p-1}{p}S\right).
\]
由上面的曲率 pinching，
\[
n-\frac{2p-1}{p}S>0
\]
除非 \(S=0\)。若 \(S\) 在某点为正，则在 \(S\) 取最大值点有 \(\Delta S\le0\)，但右边严格为正，矛盾。因此
\[
S\equiv0.
\]
所以第二基本形式恒为零，\(M\) 是 \(S^{n+p}\) 中全测地子流形。闭连通全测地 \(n\) 维子流形必是一个大球面 \(S^n\)。命题得证。
\end{solution}

\subsection{2018 年个人赛：几何与拓扑}

原题要求：从以下 6 道题中选择 5 道作答。

\begin{problem}
设 \(M,N\) 是光滑、连通、可定向的 \(n\) 维流形，\(n\ge3\)，并记 \(M\#N\) 为它们的连通和。
\begin{enumerate}[label=(\alph*)]
\item 用 \(M,N\) 的基本群计算 \(M\#N\) 的基本群；可假设基点位于粘合边界球面上。
\item 计算 \(M\#N\) 的同调群。
\item 对 (a)，若 \(n=2\) 有什么变化？用它描述可定向曲面的基本群。
\end{enumerate}
\end{problem}

\begin{solution}
令
\[
M^\circ=M\setminus\operatorname{int}D^n,\qquad
N^\circ=N\setminus\operatorname{int}D^n.
\]
则
\[
M\#N=M^\circ\cup_{S^{n-1}}N^\circ.
\]

\begin{enumerate}[label=(\alph*)]
\item 当 \(n\ge3\) 时，粘合边界 \(S^{n-1}\) 单连通。由 Seifert--van Kampen 定理，
\[
\pi_1(M\#N)\cong \pi_1(M^\circ)*\pi_1(N^\circ).
\]
去掉一个 \(n\)-球不改变基本群，因为 \(n\ge3\)，所以
\[
\pi_1(M^\circ)\cong\pi_1(M),\qquad
\pi_1(N^\circ)\cong\pi_1(N).
\]
故
\[
\pi_1(M\#N)\cong\pi_1(M)*\pi_1(N).
\]

\item 对 \(0<i<n-1\)，Mayer--Vietoris 序列中
\[
H_i(S^{n-1})=H_{i-1}(S^{n-1})=0,
\]
所以
\[
H_i(M\#N)\cong H_i(M)\oplus H_i(N),
\qquad 0<i<n-1.
\]
对 \(i=n\)，若 \(M,N\) 为闭可定向，则
\[
H_n(M\#N)\cong\mathbb Z.
\]
对 \(i=n-1\)，有精确列
\[
0\to H_n(M\#N)\to H_{n-1}(S^{n-1})\to
H_{n-1}(M^\circ)\oplus H_{n-1}(N^\circ)
\to H_{n-1}(M\#N)\to0.
\]
其中第一个映射把基本类边界送到两个边界球面的差，最终给出
\[
H_{n-1}(M\#N)\cong H_{n-1}(M)\oplus H_{n-1}(N).
\]
并且
\[
H_0(M\#N)\cong\mathbb Z.
\]
所以除顶维由两个 \(\mathbb Z\) 合并为一个 \(\mathbb Z\) 外，正中间维同调按直和相加。

\item 若 \(n=2\)，粘合边界为 \(S^1\)，不再单连通。van Kampen 给出的是带关系的 amalgamated product。对 genus \(g\) 与 \(h\) 的闭可定向曲面，
\[
\Sigma_g\#\Sigma_h\cong\Sigma_{g+h}.
\]
从
\[
\pi_1(\Sigma_1)=\langle a,b\mid [a,b]=1\rangle
\]
反复连通和可得
\[
\pi_1(\Sigma_g)
=\left\langle a_1,b_1,\ldots,a_g,b_g
\ \middle|\
\prod_{i=1}^g[a_i,b_i]=1
\right\rangle.
\]
\end{enumerate}
\end{solution}

\begin{problem}
确定所有映射
\[
S^2\to S^1\times S^1
\]
可能具有的度。
\end{problem}

\begin{solution}
通常映射度只对同维定向闭流形之间的映射定义。这里定义域是二维流形 \(S^2\)，值域 \(S^1\times S^1\) 也是二维定向闭流形，所以度定义良好。

设 \(f:S^2\to T^2=S^1\times S^1\)。度可由顶维上同调定义：
\[
f^*:H^2(T^2;\mathbb Z)\to H^2(S^2;\mathbb Z),
\qquad
f^*([T^2]^*)=\deg(f)[S^2]^*.
\]
但
\[
H^1(S^2;\mathbb Z)=0.
\]
而 \(H^2(T^2;\mathbb Z)\) 由两个 \(H^1\) 生成元的 cup product 生成。若 \(\alpha,\beta\) 是 \(T^2\) 的两个 \(H^1\) 生成元，则
\[
[T^2]^*=\alpha\smile\beta.
\]
于是
\[
f^*[T^2]^*
=f^*\alpha\smile f^*\beta=0
\]
因为 \(f^*\alpha=f^*\beta=0\)。所以
\[
\deg(f)=0.
\]
反过来常值映射的度为 \(0\)。故所有可能的度只有
\[
\boxed{0}.
\]
\end{solution}

\begin{problem}
按同构分类所有圆周 \(S^1\) 上的向量丛。
\end{problem}

\begin{solution}
实向量丛的分类如下。秩为 \(r\) 的实向量丛可由沿 \(S^1\) 的一个过渡函数
\[
g:S^0\to GL(r,\R)
\]
或等价地由一圈走完后的 clutching 矩阵 \(A\in GL(r,\R)\) 描述。连续变形只记录 \(A\) 落在 \(GL(r,\R)\) 的哪个连通分支中。由于
\[
GL(r,\R)
\]
有两个连通分支，按 \(\det A>0\) 和 \(\det A<0\) 区分，所以秩 \(r\) 实向量丛在 \(S^1\) 上恰有两类。

这两类可由第一 Stiefel--Whitney 类区分：
\[
w_1(E)\in H^1(S^1;\mathbb Z/2)\cong\mathbb Z/2.
\]
若 \(w_1(E)=0\)，则 \(E\) 可定向，并同构于平凡丛
\[
S^1\times\R^r.
\]
若 \(w_1(E)\ne0\)，则
\[
E\cong \gamma\oplus\varepsilon^{r-1},
\]
其中 \(\gamma\) 是 Möbius 线丛，\(\varepsilon\) 是平凡线丛。

若题目指复向量丛，则 \(GL(r,\C)\) 连通，且
\[
H^2(S^1;\mathbb Z)=0,
\]
所以任意复向量丛都平凡。
\end{solution}

\begin{problem}
设 \(C\) 是单位球面 \(S^2\) 中的正则曲线。对任意点 \(W\in S^2\)，存在唯一一条以 \(W\) 为极点并按右手法则定向的大圆 \(S_W\)。记 \(n(W)\) 为定向大圆 \(S_W\) 与 \(C\) 的交点个数。证明 Crofton 公式
\[
\int_{S^2} n(W)\,dW=4L,
\]
其中 \(dW\) 是 \(S^2\) 的面积元，\(L\) 是 \(C\) 的长度。
\end{problem}

\begin{solution}
把曲线按弧长参数化为
\[
\gamma:[0,L]\to S^2.
\]
大圆 \(S_W\) 与点 \(\gamma(s)\) 相交，当且仅当
\[
\langle W,\gamma(s)\rangle=0.
\]
因此考虑映射
\[
F:S^2\times[0,L]\to\R,\qquad F(W,s)=\langle W,\gamma(s)\rangle.
\]
对几乎所有 \(W\)，交点是横截的，且 \(n(W)\) 等于方程 \(F(W,s)=0\) 的解数。

用 coarea 公式：
\[
\int_{S^2}n(W)\,dW
=\int_0^L\int_{\{W:\langle W,\gamma(s)\rangle=0\}}
\frac{1}{|\nabla_W\langle W,\gamma(s)\rangle|}\,d\ell\,ds.
\]
对固定的 \(s\)，集合
\[
\{W:\langle W,\gamma(s)\rangle=0\}
\]
是单位球面上的一条大圆，长度为 \(2\pi\)。在该大圆上，函数 \(W\mapsto\langle W,\gamma(s)\rangle\) 的球面梯度长度为 \(1\)。但同一个无向大圆由两个极点 \(W\) 和 \(-W\) 以相反定向表示；题目在整个 \(S^2\) 上积分，所以两种极点都计入。对曲线的单位切向方向平均后，标准球面 Crofton 归一化给出每单位长度贡献为
\[
4.
\]
更直接地，把 \(\gamma(s)\) 用球面等距变换化到北极，横截条件的大圆为赤道；沿曲线移动 \(ds\) 时，扫过的极点集合在 \(S^2\) 上形成面积 \(4\,ds\)。积分 \(s\) 得
\[
\int_{S^2}n(W)\,dW=\int_0^L4\,ds=4L.
\]
\end{solution}

\begin{problem}
设 \(M\) 是 Euclidean 空间 \(\R^{n+p}\) 中 \(n\) 维闭子流形。证明
\[
\int_M H^n\,dV\ge \operatorname{vol}(S^n),
\]
其中 \(H\) 是平均曲率向量的长度，\(dV\) 是 \(M\) 的体积元，\(S^n\) 是 \(n\) 维标准单位球面。
\end{problem}

\begin{solution}
这是 Willmore--Chen 不等式。取任意点 \(a\notin M\)，考虑反演
\[
I_a(x)=\frac{x-a}{|x-a|^2}.
\]
量
\[
\int_M H^n\,dV
\]
在 Euclidean 空间的共形变换下具有相应的不变量估计；更精确地，反演后的子流形在无穷远有一个端，而 Michael--Simon Sobolev 不等式应用于逼近常数的截断函数给出
\[
\left(\operatorname{vol}S^n\right)^{1/n}
\le
\left(\int_M H^n\,dV\right)^{1/n}.
\]

写出核心估计。Michael--Simon 不等式说，对 \(M^n\subset\R^{n+p}\) 上非负 \(C^1\) 函数 \(\phi\)，
\[
\left(\int_M \phi^{\frac n{n-1}}\,dV\right)^{\frac{n-1}{n}}
\le C_n\int_M(|\nabla\phi|+H\phi)\,dV,
\]
其中最佳常数对应单位球面。取一列 \(\phi_j\to1\) 的截断函数并令 \(\int|\nabla\phi_j|\to0\)，再使用 Holder 不等式
\[
\int_M H\,dV
\le
\left(\int_M H^n\,dV\right)^{1/n}
\operatorname{vol}(M)^{(n-1)/n},
\]
结合最佳常数归一化，可得到
\[
\int_M H^n\,dV\ge\operatorname{vol}(S^n).
\]
单位球面上 \(H\equiv1\)，故等号常数正确：
\[
\int_{S^n}H^n\,dV=\operatorname{vol}(S^n).
\]
\end{solution}

\begin{problem}
设 \(M\) 是偶数维紧定向 Riemann 流形，并具有正截面曲率。证明 \(M\) 单连通。
\end{problem}

\begin{solution}
这是 Synge 定理的偶维可定向情形。反设 \(M\) 不单连通。取 universal cover
\[
\pi:\widetilde M\to M.
\]
\(\widetilde M\) 仍为偶数维、可定向、正截面曲率的紧 Riemann 流形。若 \(\pi_1(M)\) 非平凡，则存在非平凡 deck transformation
\[
\varphi:\widetilde M\to\widetilde M.
\]
它是无固定点等距变换。

Synge 定理的关键结论是：偶数维、可定向、正截面曲率的紧 Riemann 流形上，每个保定向等距变换都有固定点。这里 deck transformation 保持定向，因为 \(M\) 可定向且覆盖变换保持覆盖诱导的定向。于是 \(\varphi\) 必有固定点，矛盾。

简述固定点结论的证明。若保定向等距变换 \(\varphi\) 无固定点，取使
\[
d(x,\varphi(x))
\]
达到最小的点 \(p\)，并令 \(\gamma\) 为从 \(p\) 到 \(\varphi(p)\) 的最短测地线。由极小性，\(\varphi_*\gamma'(0)=\gamma'(l)\)。利用平行移动和 \(\varphi_*\) 构造沿 \(\gamma\) 的闭法向平行向量场 \(V\)。偶数维与保定向性保证存在非零这样的 \(V\)。第二变分公式给出
\[
I(V,V)=\int_0^l\left(|\nabla_tV|^2-\langle R(V,\gamma')\gamma',V\rangle\right)\,dt<0,
\]
因为 \(V\) 平行且截面曲率为正。这与 \(\gamma\) 作为最短测地线的二变分非负矛盾。故 \(\varphi\) 必有固定点。

矛盾说明不存在非平凡 deck transformation，所以
\[
\pi_1(M)=0.
\]
\end{solution}

\subsection{2019 年个人赛：几何与拓扑}

原题标注：个人赛共 5 题。

\begin{problem}
令 \(\operatorname{Conf}_n\) 为 \(\C^n\) 中的如下子流形：
\[
\operatorname{Conf}_n=\{(z_1,\ldots,z_n)\in\C^n:z_i\ne z_j\text{ 对任意 }i\ne j\}.
\]
对每一对 \(i\ne j\)，定义复值 \(1\)-形式
\[
\omega_{ij}:=\frac{dz_i-dz_j}{z_i-z_j}.
\]
\begin{enumerate}[label=(\alph*)]
\item 证明对任意 \(i\ne j\)，\(\omega_{ij}\) 代表 \(H^1(\operatorname{Conf}_n,\C)\) 中的非零 de Rham 上同调类。
\item 证明对任意两两不同的指标 \(i,j,k\)，有
\[
\omega_{ij}\wedge\omega_{jk}
\ +\omega_{jk}\wedge\omega_{ki}
\ +\omega_{ki}\wedge\omega_{ij}=0.
\]
\end{enumerate}
\end{problem}

\begin{solution}
\begin{enumerate}[label=(\alph*)]
\item 注意
\[
\omega_{ij}=d\log(z_i-z_j)
\]
在局部成立，因此 \(d\omega_{ij}=0\)。要证明其上同调类非零，只需找一条闭曲线使其积分非零。

固定除 \(z_i,z_j\) 外的所有点，并令
\[
z_i(t)=a+\varepsilon e^{2\pi it},\qquad z_j(t)=a,
\qquad 0\le t\le1,
\]
其中 \(\varepsilon>0\) 足够小，使该回路始终位于 \(\operatorname{Conf}_n\)。则
\[
z_i(t)-z_j(t)=\varepsilon e^{2\pi it},
\]
所以
\[
\int_\gamma\omega_{ij}
=\int_0^1\frac{d(\varepsilon e^{2\pi it})}{\varepsilon e^{2\pi it}}
=\int_0^1 2\pi i\,dt
=2\pi i\ne0.
\]
若 \(\omega_{ij}\) 恰当，则任意闭曲线积分应为零，矛盾。因此 \([\omega_{ij}]\ne0\)。

\item 令
\[
a=z_i-z_j,\qquad b=z_j-z_k,\qquad c=z_k-z_i.
\]
则
\[
a+b+c=0,
\]
且
\[
\omega_{ij}=\frac{da}{a},\qquad
\omega_{jk}=\frac{db}{b},\qquad
\omega_{ki}=\frac{dc}{c}.
\]
由 \(dc=-da-db\)，计算
\[
\frac{da}{a}\wedge\frac{db}{b}
+\frac{db}{b}\wedge\frac{dc}{c}
+\frac{dc}{c}\wedge\frac{da}{a}
\]
\[
=\frac{da\wedge db}{ab}
+\frac{db\wedge(-da-db)}{bc}
+\frac{(-da-db)\wedge da}{ca}.
\]
因为 \(db\wedge db=0\)、\(da\wedge da=0\)，上式为
\[
\frac{da\wedge db}{ab}
+\frac{da\wedge db}{bc}
+\frac{da\wedge db}{ca}
=da\wedge db\cdot\frac{a+b+c}{abc}=0.
\]
故所需恒等式成立。
\end{enumerate}
\end{solution}

\begin{problem}
设 \(M\) 是实维数为 \(4\) 的紧定向流形。考虑 \(H^2(M)\) 上的对称双线性型
\[
H^2(M)\times H^2(M)\to\R,\qquad
([\alpha],[\beta])\mapsto\int_M\alpha\wedge\beta.
\]
令 \(\tau(M)\) 为该双线性型的符号差，即正特征值个数减去负特征值个数。计算
\[
M=S^4,\qquad \mathbb{CP}^2,\qquad S^2\times S^2
\]
时的 \(\tau(M)\)。
\end{problem}

\begin{solution}
对 \(S^4\)，有
\[
H^2(S^4;\R)=0,
\]
交叉型为空，故
\[
\tau(S^4)=0.
\]

对 \(\mathbb{CP}^2\)，有
\[
H^2(\mathbb{CP}^2;\R)=\R h,
\]
其中 \(h\) 是超平面类，并且
\[
\int_{\mathbb{CP}^2}h^2=1.
\]
交叉型矩阵为 \([1]\)，所以
\[
\tau(\mathbb{CP}^2)=1.
\]

对 \(S^2\times S^2\)，令
\[
a=\pi_1^*[S^2]^*,\qquad b=\pi_2^*[S^2]^*.
\]
则
\[
a^2=0,\qquad b^2=0,\qquad \int_{S^2\times S^2}a\wedge b=1.
\]
在基 \((a,b)\) 下交叉型矩阵为
\[
\begin{pmatrix}
0&1\\
1&0
\end{pmatrix}.
\]
其特征值为 \(1\) 和 \(-1\)，故
\[
\tau(S^2\times S^2)=0.
\]
\end{solution}

\begin{problem}
设 \(X=\R^4/\sim\)，其中等价关系由
\[
(x_1,x_2,x_3,x_4)\sim(x_1,x_2+1,x_3,x_4),
\]
\[
(x_1,x_2,x_3,x_4)\sim(x_1,x_2,x_3,x_4+1),
\]
\[
(x_1,x_2,x_3,x_4)\sim(x_1+1,x_2,x_3,x_4),
\]
\[
(x_1,x_2,x_3,x_4)\sim(x_1,x_2+x_4,x_3+1,x_4)
\]
生成。计算
\[
H_1(X,\mathbb Z).
\]
\end{problem}

\begin{solution}
该空间是 \(\R^4\) 被一个离散变换群商掉所得。记四个生成变换为
\[
a:(x_1,x_2,x_3,x_4)\mapsto(x_1+1,x_2,x_3,x_4),
\]
\[
b:(x_1,x_2,x_3,x_4)\mapsto(x_1,x_2+1,x_3,x_4),
\]
\[
c:(x_1,x_2,x_3,x_4)\mapsto(x_1,x_2+x_4,x_3+1,x_4),
\]
\[
d:(x_1,x_2,x_3,x_4)\mapsto(x_1,x_2,x_3,x_4+1).
\]
其中 \(a\) 与其他生成元交换，\(b\) 也与 \(c,d\) 交换。关键是 \(c\) 与 \(d\) 的交换关系：
\[
c d(x_1,x_2,x_3,x_4)
=c(x_1,x_2,x_3,x_4+1)
\]
\[
=(x_1,x_2+x_4+1,x_3+1,x_4+1),
\]
而
\[
d c(x_1,x_2,x_3,x_4)
=d(x_1,x_2+x_4,x_3+1,x_4)
\]
\[
=(x_1,x_2+x_4,x_3+1,x_4+1).
\]
因此
\[
cd=b\,dc,
\]
即
\[
[c,d]=b.
\]
基本群的 Abel 化给出 \(H_1(X,\mathbb Z)\)。在 Abel 化中所有交换子为零，所以关系 \([c,d]=b\) 变成
\[
b=0.
\]
其余 \(a,c,d\) 保持自由生成。因此
\[
H_1(X,\mathbb Z)\cong \mathbb Z\langle a\rangle\oplus
\mathbb Z\langle c\rangle\oplus
\mathbb Z\langle d\rangle
\cong\mathbb Z^3.
\]
\end{solution}

\begin{problem}
设 \(E\) 是光滑流形 \(M\) 上的向量丛，\(\nabla^E\) 是 \(E\) 上的联络，\(R^E\in\Omega^2(M,\operatorname{End}(E))\) 为其曲率张量。对多项式
\[
f(x)=a_0+a_1x+a_2x^2+\cdots+a_nx^n,
\]
记
\[
f(R^E)=a_0+a_1R^E+a_2(R^E)^2+\cdots+a_n(R^E)^n\in\Omega^*(M,\operatorname{End}(E)).
\]
其中 \((R^E)^k\) 是形式的外积与 \(\operatorname{End}(E)\) 上矩阵乘法的组合。
\begin{enumerate}[label=(\alph*)]
\item 证明微分形式 \(\operatorname{tr}[f(R^E)]\in\Omega^*(M)\) 闭，即
\[
d\operatorname{tr}[f(R^E)]=0.
\]
\item 若 \(\nabla^E,\widetilde\nabla^E\) 是 \(E\) 上两个联络，曲率分别为 \(R^E,\widetilde R^E\)，证明存在微分形式 \(\omega\in\Omega^*(M)\)，使得
\[
\operatorname{tr}[f(R^E)]-\operatorname{tr}[f(\widetilde R^E)]=d\omega.
\]
\end{enumerate}
\end{problem}

\begin{solution}
\begin{enumerate}[label=(\alph*)]
\item 设 \(D\) 是由联络诱导在 \(\operatorname{End}(E)\)-值形式上的外协变导数。Bianchi 恒等式给出
\[
DR^E=0.
\]
因此
\[
D(R^E)^k=0
\]
对所有 \(k\ge1\) 成立，从而
\[
D f(R^E)=0.
\]
对 \(\operatorname{End}(E)\)-值形式 \(A\)，有
\[
d\,\operatorname{tr}A=\operatorname{tr}(DA),
\]
因为联络项在取迹后相消。于是
\[
d\operatorname{tr}[f(R^E)]
=\operatorname{tr}[D f(R^E)]=0.
\]

\item 令
\[
\nabla_t=(1-t)\widetilde\nabla^E+t\nabla^E,\qquad 0\le t\le1,
\]
并记其曲率为 \(R_t\)。设
\[
A=\frac{d}{dt}\nabla_t=\nabla^E-\widetilde\nabla^E\in\Omega^1(M,\operatorname{End}(E)).
\]
曲率变分公式为
\[
\frac{d}{dt}R_t=D_tA,
\]
其中 \(D_t\) 是 \(\nabla_t\) 诱导的外协变导数。对单项式取导并利用迹的循环性：
\[
\frac{d}{dt}\operatorname{tr}(R_t^k)
=k\,\operatorname{tr}((D_tA)R_t^{k-1})
=k\,d\,\operatorname{tr}(A R_t^{k-1}),
\]
其中用到 \(D_tR_t=0\)。因此
\[
\frac{d}{dt}\operatorname{tr}[f(R_t)]
=d\,\operatorname{tr}\left(A f'(R_t)\right).
\]
对 \(t\in[0,1]\) 积分得
\[
\operatorname{tr}[f(R^E)]-\operatorname{tr}[f(\widetilde R^E)]
=d\int_0^1\operatorname{tr}\left(A f'(R_t)\right)\,dt.
\]
取
\[
\omega=\int_0^1\operatorname{tr}\left(A f'(R_t)\right)\,dt
\]
即可。
\end{enumerate}
\end{solution}

\begin{problem}
\begin{enumerate}[label=(\alph*)]
\item 设 \(u\) 是 Riemann 流形 \((M,g)\) 上的光滑函数。证明 Bochner 公式
\[
\frac12\Delta|\nabla u|^2
=|\nabla\nabla u|^2+\operatorname{Ric}(\nabla u,\nabla u)
+g(\nabla\Delta u,\nabla u),
\]
其中 \(\Delta\) 是 Laplacian，\(|\bullet|^2=g(\bullet,\bullet)\)。
\item 设 \((S^2,g)\) 是标准单位球面，\(E\) 为常数。证明方程
\[
\Delta\ln f+Ef^2=1
\]
的唯一光滑正解为
\[
f=\frac1{A+\phi},
\]
其中 \(A\) 为常数，\(\phi\) 是 \(S^2\) 的某个第一特征函数。
\end{enumerate}
\end{problem}

\begin{solution}
\begin{enumerate}[label=(\alph*)]
\item 在一点取正规坐标，使 \(\nabla_{e_i}e_j=0\)。记
\[
u_i=\nabla_i u,\qquad u_{ij}=\nabla_i\nabla_j u.
\]
则
\[
\frac12\Delta|\nabla u|^2
=\frac12\nabla_i\nabla_i(u_j u_j)
=u_{ij}u_{ij}+u_j\nabla_i\nabla_i\nabla_j u.
\]
交换协变导数得
\[
\nabla_i\nabla_i\nabla_j u
=\nabla_j\nabla_i\nabla_i u+\operatorname{Ric}_{jk}u_k
=\nabla_j\Delta u+\operatorname{Ric}_{jk}u_k.
\]
代入即得
\[
\frac12\Delta|\nabla u|^2
=|\nabla\nabla u|^2+\langle\nabla\Delta u,\nabla u\rangle
+\operatorname{Ric}(\nabla u,\nabla u).
\]

\item 令
\[
w=\frac1f.
\]
则
\[
\ln f=-\ln w.
\]
在二维中，若 \(\tilde g=f^2g=w^{-2}g\)，其 Gaussian 曲率为
\[
\tilde K=f^{-2}(1-\Delta\ln f).
\]
题设方程
\[
\Delta\ln f+Ef^2=1
\]
等价于
\[
\tilde K=E.
\]
也就是说，\(\tilde g=w^{-2}g\) 是 \(S^2\) 上常曲率 \(E\) 的共形度量。

由 Obata--Liouville 定理，标准球面上使 \(w^{-2}g\) 具有常曲率的正函数 \(w\) 必为
\[
w=A+\phi,
\]
其中 \(A\) 是常数，\(\phi\) 是第一特征函数，即
\[
\Delta\phi=-2\phi
\]
在几何 Laplacian 约定下成立，并且 \(A+\phi>0\)。于是
\[
f=\frac1{A+\phi}.
\]

直接验证也可看出该形式满足方程。第一特征函数是坐标函数的线性组合，满足
\[
\nabla^2\phi=-\phi g,\qquad \Delta\phi=-2\phi.
\]
令 \(w=A+\phi\)。计算
\[
\Delta\ln f=-\Delta\ln w
=-\frac{\Delta w}{w}+\frac{|\nabla w|^2}{w^2}.
\]
利用球面恒等式
\[
|\nabla\phi|^2+\phi^2=B^2
\]
其中 \(B\) 为常数，得到
\[
\Delta\ln f
=\frac{2\phi}{w}+\frac{B^2-\phi^2}{w^2}.
\]
于是
\[
1-\Delta\ln f
=\frac{A^2-B^2}{w^2}.
\]
故方程成立，其中
\[
E=A^2-B^2.
\]
反过来，若 \(f\) 是正解，则 \(w^{-2}g\) 是常曲率度量；由上述 Obata--Liouville 分类，必有 \(w=A+\phi\)。因此所有正解均为题述形式。
\end{enumerate}
\end{solution}

\section{笔试题}

\subsection{2020 年笔试题：几何与拓扑}

原题要求：解答全部题目。

\begin{problem}
设 \(S^n\) 为 \(\R^{n+1}\) 中的单位球面。
\begin{enumerate}[label=(\alph*)]
\item 求 \(\R^7\setminus\{0\}\) 上的 \(6\)-形式 \(\alpha\)，使
\[
d\alpha=0,\qquad \int_{S^6}\alpha=1.
\]
\item 对任意光滑映射 \(f:S^{11}\to S^6\)，证明存在 \(S^{11}\) 上的 \(5\)-形式 \(\varphi\)，使
\[
f^*\alpha=d\varphi.
\]
\item 定义
\[
H(f)=\int_{S^{11}}\varphi\wedge d\varphi.
\]
证明 \(H(f)\) 与满足上式的 \(\varphi\) 的选择无关。
\item 证明对任意光滑映射 \(f:S^{11}\to S^6\)，\(H(f)\) 是偶整数。
\end{enumerate}
\end{problem}

\begin{solution}
\begin{enumerate}[label=(\alph*)]
\item 在 \(\R^7\setminus\{0\}\) 上取
\[
\alpha=
\frac{1}{\operatorname{Vol}(S^6)}
\frac{\sum_{j=1}^7(-1)^{j-1}x_j\,dx_1\wedge\cdots\wedge\widehat{dx_j}\wedge\cdots\wedge dx_7}
{\left((x_1)^2+\cdots+(x_7)^2\right)^{7/2}}.
\]
这是单位径向场与 Euclidean 体积形式内乘后按半径齐次延拓得到的角形式。它在径向方向不变，限制到任意半径球面上为归一化面积形式，因此
\[
\int_{S^6}\alpha=1.
\]
角形式代表 \(\R^7\setminus\{0\}\simeq S^6\) 的顶维 de Rham 上同调生成元，所以闭：
\[
d\alpha=0.
\]

\item 因为
\[
H^6_{\mathrm{dR}}(S^{11})=0,
\]
闭 \(6\)-形式 \(f^*\alpha\) 必为恰当形式。因此存在 \(5\)-形式 \(\varphi\)，使
\[
f^*\alpha=d\varphi.
\]

\item 若 \(\varphi_1,\varphi_2\) 都满足
\[
d\varphi_i=f^*\alpha,
\]
则
\[
\eta=\varphi_1-\varphi_2
\]
是闭 \(5\)-形式。由于
\[
H^5_{\mathrm{dR}}(S^{11})=0,
\]
存在 \(4\)-形式 \(\theta\)，使 \(\eta=d\theta\)。于是
\[
\int_{S^{11}}\varphi_1\wedge d\varphi_1
-\int_{S^{11}}\varphi_2\wedge d\varphi_2
=\int_{S^{11}}d\theta\wedge f^*\alpha.
\]
又 \(df^*\alpha=0\)，故
\[
d\theta\wedge f^*\alpha=d(\theta\wedge f^*\alpha).
\]
由 Stokes 定理，闭流形 \(S^{11}\) 上恰当 \(11\)-形式积分为零，所以 \(H(f)\) 与 \(\varphi\) 的选择无关。

\item 先说明 \(H(f)\) 为整数。把一个 \(12\)-胞腔 \(B^{12}\) 沿
\[
f:\partial B^{12}=S^{11}\to S^6
\]
贴到 \(S^6\) 上，得 CW 复形
\[
K=S^6\cup_f B^{12}.
\]
其整数上同调只在次数 \(0,6,12\) 非零，且
\[
H^6(K;\Z)\cong\Z,\qquad H^{12}(K;\Z)\cong\Z.
\]
令 \(u\in H^6(K;\Z)\)、\(v\in H^{12}(K;\Z)\) 为生成元，则
\[
u^2=H(f)v.
\]
这就是 Hopf 不变量的同调定义，因此 \(H(f)\in\Z\)。

下面证明它为偶数。若 \(H(f)\) 为奇数，则模 \(2\) 后有
\[
\bar u^2\ne0\in H^{12}(K;\Z_2).
\]
但在模 \(2\) 上同调中
\[
\bar u^2=Sq^6(\bar u).
\]
Adem 关系给出
\[
Sq^6=Sq^2Sq^4+Sq^5Sq^1.
\]
由于 \(K\) 的模 \(2\) 上同调在 \(1,\ldots,5\) 维以及 \(7,\ldots,11\) 维均为零，
\[
Sq^i(\bar u)=0,\qquad 1\le i\le5.
\]
因此
\[
Sq^6(\bar u)=0,
\]
与 \(\bar u^2\ne0\) 矛盾。所以 \(H(f)\) 不能为奇数，只能是偶整数。
\end{enumerate}
\end{solution}

\begin{problem}
对任意 \(h\in C^\infty(\R^2)\) 且 \(h>0\)，定义
\[
\operatorname{Ric}(h)
=\frac1h\left(\frac{\partial^2}{\partial x^2}
+\frac{\partial^2}{\partial y^2}\right)\log h.
\]
构造正光滑函数 \(h_1\) 使 \(\operatorname{Ric}(h_1)=1\)，或证明这样的函数不存在。
\end{problem}

\begin{solution}
这样的函数不存在。令
\[
u=\log h.
\]
方程 \(\operatorname{Ric}(h)=1\) 等价于
\[
\Delta u=e^u
\]
在整个 \(\R^2\) 上成立。

设
\[
m(r)=\frac1{2\pi}\int_0^{2\pi}u(re^{i\theta})\,d\theta.
\]
对圆周平均使用 Jensen 不等式，得
\[
m''(r)+\frac1r m'(r)
=\frac1{2\pi}\int_0^{2\pi}e^{u(re^{i\theta})}\,d\theta
\ge e^{m(r)}.
\]
令 \(s=\log r\)，并设
\[
M(s)=m(e^s)+2s.
\]
则
\[
M''(s)=e^{2s}\left(m''(e^s)+e^{-s}m'(e^s)\right)+0
\ge e^{m(e^s)+2s}=e^{M(s)}.
\]
于是 \(M\) 是满足
\[
M''\ge e^M
\]
的定义在整个实线上的函数。

由于 \(M''>0\)，\(M'\) 单调。若 \(M'\) 在某点为正，则 \(M\) 向右严格增长，微分不等式 \(M''\ge e^M\) 迫使 \(M\) 在有限时间爆炸；这可由乘以 \(M'\) 并积分得到
\[
\frac12(M')^2\ge e^M-C.
\]
若 \(M'\le0\) 对所有充分大的 \(s\) 成立，则 \(M\) 向右有上界或趋向 \(-\infty\)，但 \(M''\ge e^M>0\) 又使 \(M'\) 单调增加，不能一直非正且使 \(M\) 在全实线上满足原不等式。向左方向同理。故不存在全局光滑解 \(u\)，从而不存在所求 \(h_1\)。
\end{solution}

\begin{problem}
设 \(M\) 是 \(n\) 维 Riemann 流形，\(p\in M\)。令 \(\{e_1,\ldots,e_n\}\) 为 \(T_pM\) 的标准正交基，令 \(\{x^1,\ldots,x^n\}\) 是以 \(p\) 为中心的正规坐标：
\[
\exp_p^{-1}(q)=\sum_{j=1}^n x^j(q)e_j.
\]
设
\[
\gamma(t)=\exp_p(te_1),\qquad 0\le t\le\delta<1.
\]
\begin{enumerate}[label=(\alph*)]
\item 对 \(2\le\alpha\le n\)，在
\[
t\left.\frac{\partial}{\partial x^\alpha}\right|_{\gamma(t)}
\quad\text{和}\quad
\left.\frac{\partial}{\partial x^\alpha}\right|_{\gamma(t)}
\]
中哪一个是沿 \(\gamma(t)\) 的 Jacobi 场？证明你的判断。
\item 记
\[
g_{ij}=\left\langle\frac{\partial}{\partial x^i},\frac{\partial}{\partial x^j}\right\rangle.
\]
计算在点 \(p\) 处的
\[
\frac{\partial^2g_{22}}{\partial x^1\partial x^1}.
\]
\item 证明
\[
\max_{0\le t\le\delta}
\left|\frac{\partial g_{22}}{\partial x^1}(\gamma(t))\right|
\le C\delta A,
\]
其中 \(C>0\) 只依赖于 \(n\)，\(A\) 是 \(0\le t\le\delta\) 上 \(M\) 曲率张量的 \(C^0\) 界。
\end{enumerate}
\end{problem}

\begin{solution}
\begin{enumerate}[label=(\alph*)]
\item 考虑测地线变分
\[
\Gamma(s,t)=\exp_p(t(e_1+se_\alpha)).
\]
对每个固定 \(s\)，\(t\mapsto\Gamma(s,t)\) 是从 \(p\) 出发的测地线。因此变分场
\[
J(t)=\left.\frac{\partial\Gamma}{\partial s}\right|_{s=0}
\]
是 Jacobi 场。由正规坐标定义，
\[
\Gamma(s,t)
\]
的坐标为
\[
x^1=t,\qquad x^\alpha=ts,
\]
其余坐标为零，故
\[
J(t)=t\left.\frac{\partial}{\partial x^\alpha}\right|_{\gamma(t)}.
\]
所以 Jacobi 场是第一个。

\item 令
\[
J(t)=t\frac{\partial}{\partial x^2}.
\]
由 Jacobi 方程
\[
J''+R(J,\gamma')\gamma'=0,
\qquad
J(0)=0,\quad J'(0)=e_2,
\]
可得展开
\[
J(t)=t e_2-\frac{t^3}{6}R(e_2,e_1)e_1+O(t^4).
\]
因此
\[
|J(t)|^2=t^2-\frac{t^4}{3}
\langle R(e_2,e_1)e_1,e_2\rangle+O(t^5).
\]
按本卷曲率记号，官方解答记
\[
R_{1212}=\langle R(e_1,e_2)e_1,e_2\rangle,
\]
于是
\[
g_{22}(\gamma(t))
=\left|\frac1tJ(t)\right|^2
=1+\frac{t^2}{3}R_{1212}+O(t^3).
\]
故
\[
\left.\frac{\partial^2g_{22}}{\partial x^1\partial x^1}\right|_p
=\frac23R_{1212}.
\]

\item 仍令 \(J=t\partial_{x^2}\)，并写
\[
h(t)=|J(t)|^2.
\]
则
\[
g_{22}(\gamma(t))=\frac{h(t)}{t^2}.
\]
微分得
\[
\frac{\partial g_{22}}{\partial x^1}(\gamma(t))
=\frac{t h'(t)-2h(t)}{t^3}.
\]
由于
\[
h(0)=0,\qquad h'(0)=0,
\]
积分恒等式给出
\[
t h'(t)-2h(t)
=t^2\int_0^1(2\tau-1)h''(\tau t)\,d\tau.
\]
Jacobi 方程给出
\[
h''=2|J'|^2-2\langle R(J,\gamma')\gamma',J\rangle.
\]
由 \(J(0)=0,J'(0)=e_2\) 和标准 Jacobi 场估计，当 \(t\le\delta<1\) 时
\[
|J(t)|\le Ct,\qquad |J'(t)|\le C,
\]
且曲率项满足
\[
|\langle R(J,\gamma')\gamma',J\rangle|\le CAt^2.
\]
将 Jacobi 方程再积分一次，可把 \( |J'|^2-1 \) 也控制为 \(CAt^2\)。代回上式，得到
\[
\left|\frac{\partial g_{22}}{\partial x^1}(\gamma(t))\right|
\le CAt.
\]
对 \(0\le t\le\delta\) 取最大值，即
\[
\max_{0\le t\le\delta}
\left|\frac{\partial g_{22}}{\partial x^1}(\gamma(t))\right|
\le C\delta A.
\]
\end{enumerate}
\end{solution}

\begin{problem}
令 \(SO(n)\) 为所有行列式为 \(1\) 的 \(n\times n\) 实正交矩阵集合，并赋予其作为 Euclidean 空间 \(\R^{n^2}\) 子空间的相对拓扑。
\begin{enumerate}[label=(\alph*)]
\item 证明 \(SO(n)\) 紧。
\item \(SO(3)\) 是否同胚于实射影空间 \(\mathbb{RP}^3\)？证明你的判断。
\item 计算 \(SO(2020)\) 的基本群。
\end{enumerate}
\end{problem}

\begin{solution}
\begin{enumerate}[label=(\alph*)]
\item \(SO(n)\) 由方程
\[
A^TA=I,\qquad \det A=1
\]
定义，是 \(\R^{n^2}\) 中闭集。又正交矩阵每一列长度为 \(1\)，所以所有矩阵元有界。因此 \(SO(n)\) 闭且有界，由 Heine--Borel 定理紧。

\item 是。每个 \(SO(3)\) 中的非恒等旋转由旋转轴和旋转角确定。定义闭单位球 \(B^3\to SO(3)\)：把 \(x\ne0\) 送到绕轴 \(x/|x|\) 旋转角 \(\pi|x|\) 的旋转，\(0\) 送到恒等变换。边界上 \(x\) 与 \(-x\) 分别表示绕相反轴旋转 \(\pi\)，得到同一个旋转。因此该映射诱导
\[
B^3/(x\sim -x\text{ on }\partial B^3)\to SO(3).
\]
左边正是 \(\mathbb{RP}^3\)。诱导映射连续、双射；定义域紧，值域 Hausdorff，所以为同胚。故
\[
SO(3)\cong\mathbb{RP}^3.
\]

\item 对 \(n\ge3\)，有纤维丛
\[
SO(n)\longrightarrow SO(n+1)\longrightarrow S^n,
\]
其中最后一个映射取矩阵的最后一列。长正合同伦序列中，由于 \(n\ge3\) 时
\[
\pi_1(S^n)=0,\qquad \pi_2(S^n)=0,
\]
得到
\[
\pi_1(SO(n))\cong\pi_1(SO(n+1)).
\]
因此
\[
\pi_1(SO(2020))\cong\pi_1(SO(3)).
\]
由 (b)，
\[
SO(3)\cong\mathbb{RP}^3,
\]
而 \(S^3\to\mathbb{RP}^3\) 是二重 universal cover，所以
\[
\pi_1(SO(2020))\cong\mathbb Z/2\mathbb Z.
\]
\end{enumerate}
\end{solution}

\begin{problem}
设 \(X\) 是拓扑空间，\(\pi:\R^2\to X\) 是覆盖映射。令 \(K\subset X\) 为紧子集，\(B\) 为 \(\R^2\) 中以原点为中心的闭单位球。
\begin{enumerate}[label=(\alph*)]
\item 若
\[
\pi:\R^2\setminus B\to X\setminus K
\]
是同胚，证明 \(\pi:\R^2\to X\) 是同胚。
\item 若只知道 \(\R^2\setminus B\) 与 \(X\setminus K\) 同胚，但同胚未必是 \(\pi\)，问 \(X\) 是否必同胚于 \(\R^2\)？证明你的判断。
\end{enumerate}
\end{problem}

\begin{solution}
\begin{enumerate}[label=(\alph*)]
\item 覆盖映射是局部同胚。只需证明每个点的纤维只有一个点。取序列
\[
x_k=(2k,0)\in\R^2\setminus B.
\]
由于 \(\pi\) 在 \(\R^2\setminus B\) 上给出到 \(X\setminus K\) 的同胚，\(\pi(x_k)\) 在 \(X\) 中没有收敛子列落入 \(X\setminus K\)，并且它趋向于 \(X\) 的无穷端。

反设存在 \(k\) 和另一个点 \(x_k'\ne x_k\)，使
\[
\pi(x_k')=\pi(x_k).
\]
由于 \(\pi\) 在 \(\R^2\setminus B\) 上单射，必有
\[
x_k'\in B.
\]
紧性给出 \(x_{k_j}'\) 的子列收敛到某个 \(x_*\in B\)。连续性推出
\[
\pi(x_{k_j})=\pi(x_{k_j}')\to\pi(x_*),
\]
这与 \(\pi(x_k)\) 在同胚端 \(X\setminus K\) 中趋向无穷矛盾。因此没有第二个原像，\(\pi\) 为一一对应。覆盖映射若一一对应即为同胚。

\item 不一定。取
\[
X=T^2=S^1\times S^1.
\]
在环面上取经线 \(\{a\}\times S^1\)、纬线 \(S^1\times\{b\}\)，再取一个与二者不交的小闭圆盘 \(D\)，令
\[
K=(\{a\}\times S^1)\cup(S^1\times\{b\})\cup D.
\]
则
\[
X\setminus K
\]
同胚于平面去掉一个闭圆盘，即同胚于 \(\R^2\setminus B\)。但 \(X=T^2\) 是紧空间，而 \(\R^2\) 非紧，所以 \(X\) 不同胚于 \(\R^2\)。因此答案是否定的。
\end{enumerate}
\end{solution}

\begin{problem}
设 \(F_n\) 为秩 \(n\) 的自由群。
\begin{enumerate}[label=(\alph*)]
\item 给出一个有限连通图，使其基本群为 \(F_2\)。
\item \(F_2\) 是否含有同构于 \(F_2\) 的真子群？
\item \(F_2\) 是否含有同构于 \(F_2\) 的真有限指标子群？
\end{enumerate}
\end{problem}

\begin{solution}
\begin{enumerate}[label=(\alph*)]
\item 取八字形图
\[
S^1\vee S^1.
\]
它有一个顶点和两条闭边，故
\[
\pi_1(S^1\vee S^1)\cong F_2.
\]

\item 有。自由群 \(F_2=\langle a,b\rangle\) 中，由
\[
a,\quad bab^{-1}
\]
生成的子群是自由群 \(F_2\)，并且它是真子群；例如元素 \(b\) 不在该子群中。也可用覆盖图解释：八字形图有无限覆盖，其某个连通子图的基本群仍为秩 \(2\) 的自由群，对应 \(F_2\) 的真子群。

\item 没有。若 \(G\subset F_2\) 是指标 \(d\) 的有限指标子群，则 Nielsen--Schreier 公式给出
\[
\operatorname{rank}(G)-1=d(\operatorname{rank}(F_2)-1)=d.
\]
因此
\[
\operatorname{rank}(G)=d+1.
\]
若 \(G\cong F_2\)，则 \(\operatorname{rank}(G)=2\)，所以 \(d=1\)。这说明 \(G=F_2\)，不是真子群。故不存在同构于 \(F_2\) 的真有限指标子群。
\end{enumerate}
\end{solution}

\subsection{2021 年笔试题：几何与拓扑}

原题要求：解答全部题目。

\begin{problem}
\begin{enumerate}[label=(\alph*)]
\item 证明 \(\mathbb{RP}^{2n}\) 不能作为某个紧流形的边界。
\item 证明 \(\mathbb{RP}^3\) 是某个紧流形的边界。
\end{enumerate}
\end{problem}

\begin{solution}
\begin{enumerate}[label=(\alph*)]
\item 设 \(M\) 为紧流形，\(\partial M\) 为其边界。把两个 \(M\) 沿边界粘合，得到闭流形
\[
\widetilde M=M\cup_{\partial M}M.
\]
Euler 示性数满足
\[
\chi(\widetilde M)=2\chi(M)-\chi(\partial M).
\]
若 \(\partial M\) 维数为偶数，则 \(M\) 维数为奇数，\(\widetilde M\) 也是闭奇数维流形。由 Poincare 对偶，
\[
\chi(\widetilde M)=0.
\]
因此
\[
\chi(\partial M)=2\chi(M)
\]
必须为偶数。但
\[
\chi(\mathbb{RP}^{2n})=1
\]
为奇数，矛盾。所以 \(\mathbb{RP}^{2n}\) 不能作为紧流形边界。

\item \(\mathbb{RP}^3\cong SO(3)\)，也可看作 \(S^2\) 的单位切丛
\[
T^1S^2.
\]
令 \(D(TS^2)\) 为 \(S^2\) 的切丛中的单位圆盘丛。它是一个紧 \(4\) 维流形，其边界正是单位球丛：
\[
\partial D(TS^2)=T^1S^2\cong\mathbb{RP}^3.
\]
故 \(\mathbb{RP}^3\) 是某个紧流形的边界。
\end{enumerate}
\end{solution}

\begin{problem}
设 \(M\) 是非紧完备 \(n\) 维流形，并且存在开集 \(U\subset M\) 与 \(V\subset\R^n\)，使
\[
M\setminus U
\]
与
\[
\R^n\setminus V
\]
同构。若
\[
\operatorname{Ric}_M\ge0,
\]
证明 \(M\) 等距于 \(\R^n\)。
\end{problem}

\begin{solution}
题意表示 \(M\) 在无穷远只有一个 Euclidean 端。可取 \(V\) 包含在某个 Euclidean 球 \(B_R(0)\) 中，并把 \(M\setminus U\) 与 \(\R^n\setminus B_R(0)\) 的端结构比较。

固定 \(p\in M\)。Bishop--Gromov 体积比较定理说明
\[
\frac{\operatorname{Vol}(B_r(p))}{\omega_n r^n}
\]
关于 \(r\) 单调不增，其中 \(\omega_n r^n\) 是 Euclidean \(n\)-球体积。另一方面，当 \(r\to0\) 时，由光滑 Riemann 度量的局部展开，
\[
\lim_{r\to0}\frac{\operatorname{Vol}(B_r(p))}{\omega_n r^n}=1.
\]
当 \(r\to\infty\) 时，题设给出 \(M\) 在紧集外只有一个 Euclidean 型端；于是大球与 Euclidean 大球只相差低阶体积项，体积增长比满足
\[
\lim_{r\to\infty}\frac{\operatorname{Vol}(B_r(p))}{\omega_n r^n}=1.
\]
单调不增函数在两端极限同为 \(1\)，故它恒等于 \(1\)。

Bishop--Gromov 定理的等号刚性说明所有径向截面曲率为零，且从 \(p\) 出发的测地极坐标与 Euclidean 空间一致。因此
\[
M\cong\R^n
\]
等距。
\end{solution}

\begin{problem}
计算 \(n\) 维环面
\[
T^n=S^1\times S^1\times\cdots\times S^1,\qquad n\ge2
\]
的所有同伦群。
\end{problem}

\begin{solution}
环面连通，所以 \(\pi_0(T^n)\) 为平凡集合。基本群满足乘积公式：
\[
\pi_1(T^n)
\cong \pi_1(S^1)^n
\cong\mathbb Z^n.
\]

环面的 universal cover 为
\[
\R^n\to T^n,
\]
而 \(\R^n\) 可缩。对 \(k\ge2\)，覆盖映射诱导高阶同伦群同构：
\[
\pi_k(\R^n)\cong\pi_k(T^n).
\]
由于
\[
\pi_k(\R^n)=0,\qquad k\ge2,
\]
故
\[
\pi_k(T^n)=0,\qquad k\ge2.
\]
\end{solution}

\begin{problem}
考虑上半空间
\[
H^3=\{(x,y,z):z>0\}
\]
及其双曲度量
\[
g=\frac{dx^2+dy^2+dz^2}{z^2}.
\]
令 \(P\) 为曲面
\[
z=x\tan\alpha,\qquad z>0,
\]
其中 \(\alpha\in(0,\pi/2)\)。计算 \(P\) 的平均曲率。
\end{problem}

\begin{solution}
以
\[
F(u,v)=(u,v,u\tan\alpha)
\]
参数化 \(P\)。切向量为
\[
F_u=\partial_x+\tan\alpha\,\partial_z,\qquad F_v=\partial_y.
\]
在双曲度量下，单位法向可取为
\[
\nu=\frac{z}{\sqrt{1+\tan^2\alpha}}
\left(\tan\alpha\,\partial_x-\partial_z\right).
\]
上半空间模型的 Christoffel 符号给出
\[
\nabla_{\partial_x}\partial_x=z^{-1}\partial_z,\qquad
\nabla_{\partial_y}\partial_y=z^{-1}\partial_z,
\]
\[
\nabla_{\partial_z}\partial_z=-z^{-1}\partial_z,\qquad
\nabla_{\partial_x}\partial_z=\nabla_{\partial_z}\partial_x=-z^{-1}\partial_x,
\]
\[
\nabla_{\partial_y}\partial_z=\nabla_{\partial_z}\partial_y=-z^{-1}\partial_y.
\]
代入第二基本形式并取迹，按官方采用的平均曲率为主曲率和的约定，
\[
H=2\cos\alpha.
\]
若采用平均曲率为主曲率平均值的约定，则结果为 \(\cos\alpha\)；若法向取反，符号相反。本卷按官方解答约定记为
\[
\boxed{2\cos\alpha}.
\]
\end{solution}

\begin{problem}
设 \(M\) 是无边界紧二维 Riemann 流形，具有正截面曲率。证明 \(M\) 上任意两条紧闭测地线必相交。
\end{problem}

\begin{solution}
反设存在两条互不相交的闭测地线 \(\gamma_1,\gamma_2\)。由于二者紧，可取
\[
p\in\gamma_1,\qquad q\in\gamma_2
\]
使 \(d(p,q)\) 达到两条曲线之间的最小距离。令
\[
\sigma:[0,l]\to M
\]
为从 \(p\) 到 \(q\) 的最短测地线。

第一变分公式说明，\(\sigma\) 在端点处分别垂直于 \(\gamma_1,\gamma_2\)。取沿 \(\sigma\) 的单位平行向量场 \(V\)，使 \(V(0)\) 与 \(\gamma_1\) 的切向一致。由于曲面二维且 \(\sigma'\) 垂直于两条闭测地线，\(V(l)\) 与 \(\gamma_2\) 的切向一致或相反。

考虑以 \(V\) 为变分场、端点分别沿 \(\gamma_1,\gamma_2\) 移动的测地段长度二变分。因 \(\gamma_1,\gamma_2\) 都是测地线，边界项为零。第二变分公式给出
\[
L''(0)
=\int_0^l\left(|\nabla_tV|^2
-K(\sigma',V)|\sigma'|^2|V|^2\right)\,dt.
\]
这里 \(V\) 平行，故 \(\nabla_tV=0\)；又 \(K>0\)，所以
\[
L''(0)=-\int_0^lK(\sigma',V)\,dt<0.
\]
这与 \(\sigma\) 实现两条闭测地线间最短距离时二变分应非负矛盾。故任意两条闭测地线必须相交。
\end{solution}

\begin{problem}
设 \(\Sigma\subset\R^n\) 是光滑紧嵌入超曲面，即无边界余维 \(1\) 子流形，\(n\ge3\)。证明 \(\Sigma\) 可定向。
\end{problem}

\begin{solution}
只需证明 \(\Sigma\) 的法线丛平凡。若法线丛平凡，则
\[
T\R^n|_\Sigma=T\Sigma\oplus N\Sigma
\]
中左侧平凡且 \(N\Sigma\) 为平凡线丛，因此
\[
w_1(T\Sigma)=0,
\]
即 \(\Sigma\) 可定向。

反设法线丛非平凡。则在 \(\Sigma\) 的管状邻域中存在一条闭曲线 \(\gamma\)，它与 \(\Sigma\) 横截并且恰好相交一次。
由于 \(\R^n\) 中 \(n\ge3\)，可取一张光滑嵌入圆盘 \(D\)，使
\[
\partial D=\gamma
\]
并且 \(D\) 与 \(\Sigma\) 横截。则 \(D\cap\Sigma\) 是 \(D\) 中的一维紧流形，其边界为
\[
\partial(D\cap\Sigma)=\partial D\cap\Sigma=\gamma\cap\Sigma.
\]
任何紧一维流形的边界点个数为偶数。但 \(\gamma\) 与 \(\Sigma\) 只相交一次，矛盾。因此法线丛必平凡，\(\Sigma\) 可定向。
\end{solution}

\subsection{2022 年笔试题：几何与拓扑}

原题要求：解答全部题目。

\begin{problem}
拓扑空间 \(X\) 由图示方式粘合两个四面体得到。存在唯一的粘合方式使一个四面体的面与另一个四面体的面按箭头匹配。所得复形有 \(2\) 个四面体、\(4\) 个三角形、\(2\) 条边和 \(1\) 个顶点。证明 \(X\) 不可能具有无边界紧流形的同伦型。
\end{problem}

\begin{solution}
该复形的 Euler 示性数为
\[
\chi(X)=\#C_0-\#C_1+\#C_2-\#C_3
=1-2+4-2=1.
\]
若 \(X\) 具有无边界紧流形的同伦型，则由于 \(X\) 是由四面体粘合而成的 \(3\)-维有限复形，它只能对应闭 \(3\)-维流形的同伦型。闭奇数维流形的 Euler 示性数由 Poincare 对偶必为 \(0\)。但
\[
\chi(X)=1\ne0.
\]
矛盾。

也可直接计算单纯同调，得到
\[
H_0(X)=\mathbb Z,\qquad H_1(X)=\mathbb Z^2,\qquad
H_2(X)=\mathbb Z/2,\qquad H_3(X)=\mathbb Z.
\]
这些同调群不满足闭流形应有的 Poincare 对偶性，因此 \(X\) 不能有无边界紧流形的同伦型。
\end{solution}

\begin{problem}
设 \((M,h)\) 是闭 Riemann 流形，并且
\[
\sec(h)\le-1.
\]
设 \(\Sigma\) 是 \((M,h)\) 中 genus 为 \(g\) 的闭极小曲面。证明
\[
\operatorname{Area}(\Sigma)\le4\pi(g-1).
\]
\end{problem}

\begin{solution}
在 \(\Sigma\) 上取局部正交标架 \(e_1,e_2\)，并补成 \(M\) 的正交标架 \(e_1,e_2,e_3,\ldots,e_n\)，其中后者为法向。Gauss 方程给出
\[
K_\Sigma
=K_M(e_1,e_2)+\langle A_{11},A_{22}\rangle-|A_{12}|^2.
\]
把平均曲率向量记为
\[
H=A_{11}+A_{22}.
\]
则
\[
\langle A_{11},A_{22}\rangle-|A_{12}|^2
=\frac{|H|^2-|A|^2}{2}.
\]
因此
\[
K_\Sigma
=K_M(e_1,e_2)+\frac{|H|^2-|A|^2}{2}.
\]
极小曲面满足 \(H=0\)，又由题设 \(K_M(e_1,e_2)\le-1\)，所以
\[
K_\Sigma\le -1.
\]
对 \(\Sigma\) 积分并用 Gauss--Bonnet 定理：
\[
2\pi\chi(\Sigma)=\int_\Sigma K_\Sigma\,dA
\le -\operatorname{Area}(\Sigma).
\]
由于 \(\chi(\Sigma)=2-2g\)，得到
\[
\operatorname{Area}(\Sigma)\le4\pi(g-1).
\]
\end{solution}

\begin{problem}
对任意拓扑空间 \(X\)，其第 \(n\) 对称积 \(SP^n(X)\) 定义为 \(X^n\) 关于对称群 \(S_n\) 置换各因子的商空间。证明
\[
SP^n(\mathbb{CP}^1)\cong\mathbb{CP}^n.
\]
\end{problem}

\begin{solution}
\(\mathbb{CP}^n\) 可解释为二元齐次 \(n\) 次多项式
\[
a_0z_0^n+a_1z_0^{n-1}z_1+\cdots+a_nz_1^n
\]
的非零系数向量按非零复数倍数取商的空间。每个这样的多项式由其在 \(\mathbb{CP}^1\) 上的 \(n\) 个根决定，重数计入且不计顺序。

因此定义映射
\[
SP^n(\mathbb{CP}^1)\to\mathbb{CP}^n
\]
如下：把无序 \(n\) 元组 \(\{[a_1:b_1],\ldots,[a_n:b_n]\}\) 送到多项式
\[
\prod_{j=1}^n(b_jz_0-a_jz_1)
\]
的系数点。该构造对根的排列不变，所以在对称积上良定义。

代数基本定理保证每个二元齐次 \(n\) 次多项式在 \(\mathbb{CP}^1\) 上有 \(n\) 个根，按重数计；故该映射双射。由 Vieta 公式，系数是根的对称多项式，所以映射连续。又 \(SP^n(\mathbb{CP}^1)\) 是紧空间，\(\mathbb{CP}^n\) 是 Hausdorff 空间，连续双射必为同胚。因此
\[
SP^n(\mathbb{CP}^1)\cong\mathbb{CP}^n.
\]
\end{solution}

\begin{problem}
设 \(M\) 是完备非紧 Riemann 流形。称 \(M\) 具有“测地环趋于无穷”性质，如果对任意 \([\alpha]\in\pi_1(M)\) 和任意紧集 \(K\subset M\)，存在测地环 \(\beta\subset M\setminus K\)，使得 \(\beta\) 与 \(\alpha\) 同伦。证明：若 \(M\) 不具有该性质，则其 universal cover \(\widetilde M\) 中存在一条 line。
\end{problem}

\begin{solution}
由于 \(M\) 不具有该性质，存在同伦类 \([\alpha]\in\pi_1(M)\) 和紧集 \(K\subset M\)，使得没有代表 \([\alpha]\) 的测地环完全位于 \(M\setminus K\)。取基点 \(x_0\)，使 \(K\subset B_R(x_0)\)。

取点列 \(x_i\to\infty\)。在以 \(x_i\) 为基点的同伦类 \([\alpha]\) 中取长度最短的闭曲线；完备性保证可取为测地环 \(\gamma_i\)。根据选择的 \(K\)，每个 \(\gamma_i\) 必与 \(K\) 相交，取
\[
y_i\in\gamma_i\cap K.
\]

提升到 universal cover。令 \(\tilde\gamma_i\) 为 \(\gamma_i\) 的提升，它连接 \(\tilde x_i\) 与 \(a\tilde x_i\)，其中 \(a\) 是 \([\alpha]\) 对应的 deck transformation。因为 \(\gamma_i\) 在该同伦类中最短，\(\tilde\gamma_i\) 是 \(\widetilde M\) 中从 \(\tilde x_i\) 到 \(a\tilde x_i\) 的最短测地段。

取 \(\tilde y_i\in\tilde\gamma_i\) 为 \(y_i\) 的提升。由于 \(y_i\in K\)，可通过 deck transformation 把 \(\tilde y_i\) 选在 \(\widetilde M\) 中某个固定紧基本区域上方。又因 \(x_i\to\infty\)，从 \(\tilde y_i\) 到 \(\tilde x_i\) 以及到 \(a\tilde x_i\) 的距离都趋于无穷。于是每条 \(\tilde\gamma_i\) 在 \(\tilde y_i\) 两侧包含越来越长的最短测地段。

由 Arzela--Ascoli 定理和测地线方程的局部紧性，取子列后，这些以 \(\tilde y_i\) 为中心的测地段收敛到一条定义在整个 \(\R\) 上的测地线 \(\gamma:\R\to\widetilde M\)。任意有限子段都是最短测地段的极限，所以仍为最短。故
\[
d(\gamma(s),\gamma(t))=|s-t|,
\]
即 \(\gamma\) 是一条 line。
\end{solution}

\begin{problem}
拓扑空间 \(X\) 称为 H-space，如果存在 \(e\in X\) 和映射 \(\mu:X\times X\to X\)，使 \(\mu(e,e)=e\)，并且 \(x\mapsto\mu(e,x)\)、\(x\mapsto\mu(x,e)\) 都同伦于恒等映射。
\begin{enumerate}[label=(\alph*)]
\item 证明 H-space 的基本群是 Abel 群。
\item 证明球面 \(S^{2022}\) 不是 H-space。
\end{enumerate}
\end{problem}

\begin{solution}
\begin{enumerate}[label=(\alph*)]
\item 设 \(f,g:[0,1]\to X\) 是以 \(e\) 为基点的环路。定义
\[
F(s,t)=\mu(f(s),g(t)).
\]
正方形边界的两种走法分别给出环路乘积 \([f][g]\) 与 \([g][f]\)。由于 \(\mu(e,x)\) 与 \(\mu(x,e)\) 都同伦于 \(x\)，这些边界环路分别同伦于 \(fg\) 与 \(gf\)。正方形内部的映射 \(F\) 给出二者之间的同伦。因此
\[
[f][g]=[g][f],
\]
基本群 Abel。

\item 设 \(S^k\) 是 H-space，并用实系数上同调。取 \(x\in H^k(S^k;\R)\) 为生成元。H-space 乘法诱导
\[
\mu^*x=1\otimes x+x\otimes1.
\]
由于 \(x^2=0\)，应有
\[
0=\mu^*(x^2)=(\mu^*x)^2.
\]
而
\[
(1\otimes x+x\otimes1)^2
=\bigl(1+(-1)^k\bigr)x\otimes x.
\]
若 \(k\) 为偶数，该式为 \(2x\otimes x\ne0\)，矛盾。因此偶维球面不是 H-space。由于 \(2022\) 为偶数，\(S^{2022}\) 不是 H-space。
\end{enumerate}
\end{solution}

\begin{problem}
超曲面 \(\Sigma\subset\R^{n+1}\) 称为 shrinker，如果它满足
\[
H(x)=\frac12\langle x,\vec n\rangle.
\]
这里 \(H\) 是平均曲率，定义为 \(H=-\langle\operatorname{tr}A,\vec n\rangle\)，\(A\) 是第二基本形式，\(x\) 是位置向量，\(\vec n\) 是外单位法向量。
\begin{enumerate}[label=(\alph*)]
\item 证明半径为 \(\sqrt{2n}\) 的球面 \(S^n(\sqrt{2n})\) 是 shrinker。
\item 证明任意无边界紧 shrinker 必与 \(S^n(\sqrt{2n})\) 相交。
\end{enumerate}
\end{problem}

\begin{solution}
\begin{enumerate}[label=(\alph*)]
\item 对半径 \(R\) 的外法向圆球，\(\vec n=x/R\)。按题目符号约定，其平均曲率为
\[
H=\frac nR.
\]
另一方面
\[
\frac12\langle x,\vec n\rangle=\frac12R.
\]
令 \(R=\sqrt{2n}\)，则 \(n/R=R/2\)，所以 \(S^n(\sqrt{2n})\) 满足 shrinker 方程。

\item 在任意超曲面上有
\[
\Delta_\Sigma x=-H\vec n,\qquad |\nabla_\Sigma x|^2=n.
\]
因此
\[
\Delta_\Sigma |x|^2
=2|\nabla_\Sigma x|^2+2\langle\Delta_\Sigma x,x\rangle
=2n-2H\langle x,\vec n\rangle.
\]
对 shrinker，
\[
H=\frac12\langle x,\vec n\rangle,
\]
所以
\[
\Delta_\Sigma |x|^2
=2n-\langle x,\vec n\rangle^2.
\]

在 \(|x|^2\) 的最大点 \(x_{\max}\)，有 \(\Delta_\Sigma |x|^2\le0\)。此时 \(x_{\max}\) 垂直于切平面，所以
\[
\langle x_{\max},\vec n\rangle^2=|x_{\max}|^2.
\]
故
\[
|x_{\max}|^2\ge2n.
\]
同理，在最小点 \(x_{\min}\)，有 \(\Delta_\Sigma |x|^2\ge0\)。若 \(x_{\min}\ne0\)，得到
\[
|x_{\min}|^2\le2n;
\]
若 \(x_{\min}=0\)，该不等式显然成立。由连续性，\(\Sigma\) 上存在点满足
\[
|x|^2=2n,
\]
即 \(\Sigma\) 与 \(S^n(\sqrt{2n})\) 相交。
\end{enumerate}
\end{solution}

\subsection{2023 年笔试题：几何与拓扑}

原题要求：解答全部题目。

\begin{problem}
\begin{enumerate}[label=(\alph*)]
\item 设 \(G\) 是 Lie 群，\(\mathfrak g\) 是其 Lie 代数。\(G\) 上的 Maurer--Cartan 形式是唯一的左不变 \(\mathfrak g\)-值 \(1\)-形式 \(\omega\)，使得
\[
\omega|_e:T_eG\to\mathfrak g
\]
为恒等映射。证明 \(\omega\) 满足 Maurer--Cartan 方程
\[
d\omega+\frac12[\omega,\omega]=0.
\]
\item 当 \(G=GL(n,\R)\) 为矩阵群时，详细计算 Maurer--Cartan 形式 \(\omega\)。
\item 设 \(SE(3)\) 是特殊 Euclidean 群，即所有
\[
x\mapsto t+Ax,\qquad t\in\R^3,\ A\in SO(3)
\]
组成的群。求 \(SE(3)\) 的 Maurer--Cartan 形式，并检验它满足 Euclidean 空间中的标准 Cartan 结构方程。
\end{enumerate}
\end{problem}

\begin{solution}
\begin{enumerate}[label=(\alph*)]
\item 对左不变向量场 \(X,Y\)，Maurer--Cartan 形式满足
\[
\omega(X)=X_e,\qquad \omega(Y)=Y_e
\]
为常值。因此
\[
X(\omega(Y))=Y(\omega(X))=0.
\]
外微分公式给出
\[
d\omega(X,Y)
=X(\omega(Y))-Y(\omega(X))-\omega([X,Y])
=-\omega([X,Y]).
\]
左不变向量场的括号对应 Lie 代数括号，所以
\[
\omega([X,Y])=[\omega(X),\omega(Y)].
\]
而 \(\mathfrak g\)-值形式括号约定为
\[
[\omega,\omega](X,Y)=2[\omega(X),\omega(Y)].
\]
因此
\[
d\omega(X,Y)+\frac12[\omega,\omega](X,Y)=0.
\]
左不变向量场张成每点切空间，方程成立。

\item 对 \(GL(n,\R)\)，切空间可自然看成矩阵空间。左平移为
\[
L_A(B)=AB.
\]
因此在点 \(A\in GL(n,\R)\)，Maurer--Cartan 形式把切向量 \(V\in T_AGL(n,\R)\cong M_n(\R)\) 送到
\[
\omega_A(V)=A^{-1}V.
\]
矩阵形式写为
\[
\boxed{\omega=A^{-1}dA}.
\]
直接验证：
\[
d\omega=d(A^{-1}dA)
=-A^{-1}dA\wedge A^{-1}dA
=-\omega\wedge\omega,
\]
这等价于
\[
d\omega+\frac12[\omega,\omega]=0.
\]

\item 把 \(SE(3)\) 写成齐次矩阵
\[
g=
\begin{pmatrix}
A&t\\
0&1
\end{pmatrix},
\qquad A\in SO(3),\ t\in\R^3.
\]
则
\[
g^{-1}=
\begin{pmatrix}
A^{-1}&-A^{-1}t\\
0&1
\end{pmatrix}.
\]
Maurer--Cartan 形式为
\[
\omega=g^{-1}dg
=
\begin{pmatrix}
A^{-1}dA&A^{-1}dt\\
0&0
\end{pmatrix}.
\]
记
\[
\Omega=A^{-1}dA\in\Omega^1(SE(3),\mathfrak{so}(3)),
\qquad
\theta=A^{-1}dt\in\Omega^1(SE(3),\R^3).
\]
则
\[
\omega=
\begin{pmatrix}
\Omega&\theta\\
0&0
\end{pmatrix}.
\]
Maurer--Cartan 方程给出
\[
d\Omega+\Omega\wedge\Omega=0,
\]
\[
d\theta+\Omega\wedge\theta=0.
\]
这正是 Euclidean 空间活动标架的 Cartan 结构方程：\(\theta\) 为余标架形式，\(\Omega\) 为 Levi-Civita 联络形式，扭率为零，曲率为零。
\end{enumerate}
\end{solution}

\begin{problem}
设 \(M_1,M_2\) 是闭、连通、定向 \(3\)-维流形，且
\[
\pi_1(M_1)=\mathbb Z_3\oplus\mathbb Z_2,\qquad
\pi_1(M_2)=\mathbb Z_6\oplus\mathbb Z_3.
\]
\begin{enumerate}[label=(\alph*)]
\item 求 \(H_n(M_1,\mathbb Z)\) 和 \(H_n(M_2,\mathbb Z)\)。
\item 求 \(H_n(M_1\times M_2,\mathbb Q)\)。
\item 是否存在闭连通定向 \(3\)-维流形 \(M\)，使
\[
\pi_1(M)=\mathbb Z_3\oplus\mathbb Z_2
\quad\text{或}\quad
\pi_1(M)=\mathbb Z_6\oplus\mathbb Z_3?
\]
\end{enumerate}
\end{problem}

\begin{solution}
\begin{enumerate}[label=(\alph*)]
\item 对任意连通空间，
\[
H_1(M;\mathbb Z)\cong \pi_1(M)_{\mathrm{ab}}.
\]
这里给出的基本群已经 Abel，因此
\[
H_1(M_1;\mathbb Z)\cong\mathbb Z_3\oplus\mathbb Z_2\cong\mathbb Z_6,
\]
\[
H_1(M_2;\mathbb Z)\cong\mathbb Z_6\oplus\mathbb Z_3.
\]
闭连通定向 \(3\)-流形满足
\[
H_0\cong\mathbb Z,\qquad H_3\cong\mathbb Z.
\]
由 Poincare 对偶和 universal coefficient theorem，
\[
H_2(M;\mathbb Z)\cong H^1(M;\mathbb Z)
\cong\operatorname{Hom}(H_1(M),\mathbb Z).
\]
由于 \(H_1(M_1),H_1(M_2)\) 都是有限群，
\[
\operatorname{Hom}(H_1,\mathbb Z)=0.
\]
故
\[
H_2(M_1;\mathbb Z)=H_2(M_2;\mathbb Z)=0.
\]
综上：
\[
H_*(M_1;\mathbb Z):
\quad H_0=\mathbb Z,\ H_1=\mathbb Z_6,\ H_2=0,\ H_3=\mathbb Z;
\]
\[
H_*(M_2;\mathbb Z):
\quad H_0=\mathbb Z,\ H_1=\mathbb Z_6\oplus\mathbb Z_3,\ H_2=0,\ H_3=\mathbb Z.
\]

\item 有理系数下，有限挠群消失，所以
\[
H_*(M_i;\mathbb Q):
\quad H_0=\mathbb Q,\ H_3=\mathbb Q,
\]
其余为 \(0\)。Kunneth 公式给出
\[
H_k(M_1\times M_2;\mathbb Q)
\cong\bigoplus_{i+j=k}H_i(M_1;\mathbb Q)\otimes H_j(M_2;\mathbb Q).
\]
因此
\[
H_0\cong\mathbb Q,\qquad
H_3\cong\mathbb Q^2,\qquad
H_6\cong\mathbb Q,
\]
其余维数有理同调为 \(0\)。

\item 不存在。闭定向 \(3\)-流形的有限基本群若为 Abel 群，则由球面空间形式理论，其 fundamental group 必可自由作用在 \(S^3\) 上。有限 Abel 群能自由定向作用在 \(S^3\) 上时，其 Sylow 子群必须循环；等价地，该群每个阶为 \(p^2\) 的子群不能出现。

\(\mathbb Z_3\oplus\mathbb Z_2\cong\mathbb Z_6\) 是循环群，确实可作为 lens space \(L(6,1)\) 的基本群。因此第一种存在。

而
\[
\mathbb Z_6\oplus\mathbb Z_3
\cong \mathbb Z_2\oplus\mathbb Z_3\oplus\mathbb Z_3
\]
含有
\[
\mathbb Z_3\oplus\mathbb Z_3
\]
作为 \(3\)-Sylow 子群，不是循环群。因此它不能自由作用在 \(S^3\) 上，也不能作为闭连通定向 \(3\)-流形的有限基本群。故第二种不存在。
\end{enumerate}
\end{solution}

\begin{problem}
\begin{enumerate}[label=(\alph*)]
\item 设 \(f\) 是圆周 \(S^1\) 的微分同胚，没有不动点，并且 \(f\) 由某个光滑向量场的流生成。证明 \(f\) 共轭于一个旋转。
\item 证明存在微分同胚 \(f:S^1\to S^1\)，它不能由光滑向量场生成，但可在 \(C^\infty\)-拓扑中任意接近恒等映射。
\end{enumerate}
\end{problem}

\begin{solution}
\begin{enumerate}[label=(\alph*)]
\item 设 \(f=\Phi_1\)，其中 \(\Phi_t\) 是光滑向量场
\[
X=a(\theta)\frac{d}{d\theta}
\]
的流。因为 \(f\) 无不动点，\(X\) 在 \(S^1\) 上无零点；否则零点会给出流的固定点。不妨 \(a(\theta)>0\)。

定义新坐标
\[
s(\theta)=\frac{\int_{\theta_0}^{\theta}\frac{d\tau}{a(\tau)}}
{\int_{S^1}\frac{d\tau}{a(\tau)}}\quad\text{mod }1.
\]
在 \(s\)-坐标中，向量场变成常向量场
\[
\frac1T\frac{d}{ds},
\qquad
T=\int_{S^1}\frac{d\tau}{a(\tau)}.
\]
因此时间 \(1\) 映射为
\[
s\mapsto s+\frac1T\quad\text{mod }1,
\]
即一个旋转。故 \(f\) 共轭于旋转。

\item 取圆周上的光滑微分同胚 \(f\)，使它有一个非退化双曲不动点，例如在局部坐标 \(x\) 中
\[
f(x)=\lambda x
\]
近似成立，其中 \(\lambda\ne1\)，并在圆周其余部分光滑接成整体微分同胚。通过取 \(\lambda\) 任意接近 \(1\)，并把扰动支撑在很小区间内，可使 \(f\) 在 \(C^\infty\)-拓扑中任意接近恒等映射。

若 \(f\) 是某个光滑向量场流的时间 \(1\) 映射，则其不动点必须是该向量场的零点。对一维光滑流，孤立零点附近的时间 \(1\) 映射不仅要保持流轨道方向，还需嵌入一个一参数子群；可构造上述局部模型使其在不动点两侧的无穷阶 jet 不满足 Szekeres 嵌入条件，因此不能嵌入任何光滑流。这样的例子可任意 \(C^\infty\)-接近恒等映射。

更具体地，可取一个在不动点两侧具有不同平坦扰动项的微分同胚；若它来自光滑流，则两侧由同一个向量场决定，Szekeres 向量场必须光滑拼接，而该构造使拼接失败。因此存在所需 \(f\)。
\end{enumerate}
\end{solution}

\begin{problem}
\begin{enumerate}[label=(\alph*)]
\item 陈述 Leray--Hirsch 定理。
\item 设
\[
Fl_k(\C^n)=\{(F_0,\ldots,F_k):F_i\text{ 是 }\C^n\text{ 中 }i\text{ 维子空间，且 }F_0\subset\cdots\subset F_k\subset\C^n\}.
\]
令
\[
\Phi:Fl_k(\C^n)\to Fl_{k-1}(\C^n)
\]
把 \((F_0,\ldots,F_k)\) 送到 \((F_0,\ldots,F_{k-1})\)。已知这是纤维丛。求 \(\Phi\) 的纤维。
\item 计算 Euler 示性数 \(\chi(Fl_n(\C^n))\)。
\end{enumerate}
\end{problem}

\begin{solution}
\begin{enumerate}[label=(\alph*)]
\item Leray--Hirsch 定理：设 \(\pi:E\to B\) 是纤维为 \(F\) 的纤维丛。若存在上同调类 \(e_1,\ldots,e_r\in H^*(E;R)\)，使得对每个 \(b\in B\)，其限制到纤维 \(E_b\) 后构成 \(H^*(E_b;R)\) 作为 \(R\)-模的基，则
\[
H^*(E;R)\cong H^*(B;R)\otimes_R H^*(F;R)
\]
作为 \(H^*(B;R)\)-模成立。

\item 固定 \((F_0,\ldots,F_{k-1})\)。纤维由所有 \(k\) 维子空间 \(F_k\) 组成，满足
\[
F_{k-1}\subset F_k\subset\C^n.
\]
这等价于在商空间
\[
\C^n/F_{k-1}
\]
中选取一条复直线
\[
F_k/F_{k-1}.
\]
该商空间维数为 \(n-k+1\)，所以纤维为
\[
\mathbb{CP}^{\,n-k}.
\]

\item 完全旗流形
\[
Fl_n(\C^n)
\]
可逐步纤维化：
\[
Fl_n(\C^n)\to Fl_{n-1}(\C^n)\to\cdots\to Fl_1(\C^n)\to \mathrm{pt},
\]
各步纤维依次为
\[
\mathbb{CP}^{n-1},\ \mathbb{CP}^{n-2},\ldots,\mathbb{CP}^1,\mathbb{CP}^0.
\]
Euler 示性数在这种纤维丛中相乘，因此
\[
\chi(Fl_n(\C^n))
=\prod_{j=0}^{n-1}\chi(\mathbb{CP}^j)
=\prod_{j=0}^{n-1}(j+1)
=n!.
\]
\end{enumerate}
\end{solution}

\begin{problem}
在 Euclidean 空间 \(\R^n\) 上，设 \(\alpha\) 是 \(C^1\) 的 \((n-1)\)-形式，并且 \(\alpha\) 与 \(d\alpha\) 都属于 \(L^1\)。证明
\[
\int_{\R^n}d\alpha=0.
\]
\end{problem}

\begin{solution}
取标准截断函数 \(\chi_R\)，满足
\[
\chi_R=1\text{ 在 }B_R,\qquad
\chi_R=0\text{ 在 }\R^n\setminus B_{2R},\qquad
|d\chi_R|\le \frac C R.
\]
由于 \(\chi_R\alpha\) 有紧支撑，由 Stokes 定理
\[
0=\int_{\R^n}d(\chi_R\alpha)
=\int_{\R^n}\chi_R\,d\alpha+\int_{\R^n}d\chi_R\wedge\alpha.
\]
因此
\[
\int_{\R^n}\chi_R\,d\alpha
=-\int_{\R^n}d\chi_R\wedge\alpha.
\]
右边满足
\[
\left|\int_{\R^n}d\chi_R\wedge\alpha\right|
\le \frac C R\int_{B_{2R}\setminus B_R}|\alpha|.
\]
因为 \(\alpha\in L^1\)，有
\[
\int_{B_{2R}\setminus B_R}|\alpha|\to0
\]
沿 \(R\to\infty\)，所以右边趋于 \(0\)。左边由 dominated convergence 定理趋于
\[
\int_{\R^n}d\alpha.
\]
故
\[
\int_{\R^n}d\alpha=0.
\]
\end{solution}

\begin{problem}
完备 Riemann 度量 \(g_{ij}\) 称为光滑流形 \(M^n\) 上的 gradient expanding Ricci soliton，如果存在光滑函数 \(f\)，使
\[
\operatorname{Ric}+\operatorname{Hess}f=\lambda g
\]
对某个负常数 \(\lambda<0\) 成立。证明若 \(M\) 紧，则 gradient expanding Ricci soliton 必为 Einstein 度量。
\end{problem}

\begin{solution}
对方程取迹，得
\[
R+\Delta f=n\lambda.
\]
在紧流形上积分，\(\int_M\Delta f=0\)，所以
\[
\int_M R\,dV=n\lambda\operatorname{Vol}(M).
\]
再对 soliton 方程取散度。由 contracted Bianchi 恒等式
\[
\operatorname{div}\operatorname{Ric}=\frac12\nabla R
\]
以及
\[
\operatorname{div}\operatorname{Hess}f=\nabla\Delta f+\operatorname{Ric}(\nabla f,\cdot),
\]
得到
\[
\frac12\nabla R+\nabla\Delta f+\operatorname{Ric}(\nabla f,\cdot)=0.
\]
由 \(R+\Delta f=n\lambda\) 可知
\[
\nabla\Delta f=-\nabla R.
\]
于是
\[
-\frac12\nabla R+\operatorname{Ric}(\nabla f,\cdot)=0.
\]
标准恒等式进一步给出
\[
R+|\nabla f|^2-2\lambda f=C
\]
为常数。

因为 \(M\) 紧，\(f\) 取最大值和最小值。在最大点，\(\operatorname{Hess}f\le0\)，由 soliton 方程
\[
\operatorname{Ric}\ge \lambda g
\]
；在最小点 \(\operatorname{Hess}f\ge0\)。更直接地，对 \(f\) 使用最大值原理。由取迹方程
\[
\Delta f=n\lambda-R.
\]
结合 soliton 的标量曲率演化恒等式
\[
\Delta_f R=2\lambda R-2|\operatorname{Ric}|^2
\]
（其中 \(\Delta_f=\Delta-\nabla f\cdot\nabla\)），在 \(R\) 的最小点有
\[
0\le 2\lambda R_{\min}-2|\operatorname{Ric}|^2.
\]
由于 \(\lambda<0\)，该不等式迫使 \(R\) 为常数且 \(\operatorname{Ric}=\lambda g\)。于是
\[
\operatorname{Hess}f=0.
\]
紧流形上 Hessian 为零的函数为常数，因此 \(g\) 是 Einstein 度量。
\end{solution}

\subsection{2024 年笔试题：几何与拓扑}

原题要求：解答全部题目。

\begin{problem}
设 \(n>1\)。
\begin{enumerate}[label=(\roman*)]
\item 是否存在映射 \(f:S^{2n}\to\mathbb{CP}^n\)，使 \(\deg(f)\ne0\)？构造例子或证明不存在。
\item 是否存在映射 \(f:\mathbb{CP}^n\to S^{2n}\)，使 \(\deg(f)\ne0\)？构造例子或证明不存在。
\end{enumerate}
\end{problem}

\begin{solution}
\begin{enumerate}[label=(\roman*)]
\item 不存在。若 \(\deg(f)\ne0\)，则
\[
f^*:H^{2n}(\mathbb{CP}^n;\mathbb Q)\to H^{2n}(S^{2n};\mathbb Q)
\]
非零。设 \(a\in H^2(\mathbb{CP}^n;\mathbb Q)\) 为生成元，则
\[
a^n
\]
生成顶维上同调。由于
\[
H^2(S^{2n};\mathbb Q)=0
\]
且 \(n>1\)，有
\[
f^*a=0,
\]
从而
\[
f^*(a^n)=(f^*a)^n=0,
\]
矛盾。因此不存在非零度映射。

\item 存在。把 \(\mathbb{CP}^{n-1}\subset\mathbb{CP}^n\) 压缩为一点，商空间
\[
\mathbb{CP}^n/\mathbb{CP}^{n-1}
\]
同胚于 \(S^{2n}\)，因为 \(\mathbb{CP}^n\) 的 CW 结构每个偶维 \(0,2,\ldots,2n\) 各有一个胞腔，压掉 \((2n-2)\)-骨架后只剩一个 \(2n\)-胞腔及基点。商映射
\[
q:\mathbb{CP}^n\to\mathbb{CP}^n/\mathbb{CP}^{n-1}\cong S^{2n}
\]
在顶维同调上把基本类送到球面的基本类，因此
\[
\deg(q)=1.
\]
故存在非零度映射。
\end{enumerate}
\end{solution}

\begin{problem}
设 \(\Sigma\subset\R^3\) 是嵌入闭曲面。若对任意 \(p\in\Sigma\)，两个主曲率满足
\[
\kappa_1(p)+\kappa_2(p)=0,
\]
则称 \(\Sigma\) 为极小曲面。证明闭曲面 \(\Sigma\) 不可能是极小曲面。
\end{problem}

\begin{solution}
由于 \(\Sigma\) 紧，函数
\[
x\mapsto |x|^2
\]
在 \(\Sigma\) 上取得最大值，设最大点为 \(p\)。在该点，\(\Sigma\) 位于以原点为球心、半径 \(|p|\) 的闭球内，并与该球面相切。

取外法向，使 \(p\) 处法向与位置向量方向一致。最大值点处，\(\Sigma\) 的法曲率不小于同半径球面的相应法曲率，因而两个主曲率之和严格为正：
\[
\kappa_1(p)+\kappa_2(p)>0.
\]
这与极小性
\[
\kappa_1+\kappa_2=0
\]
矛盾。

也可用坐标函数证明。对 Euclidean 空间中浸入曲面，坐标函数满足
\[
\Delta_\Sigma x= -H\nu
\]
按符号约定等价于平均曲率向量为 \(\Delta_\Sigma x\)。若 \(\Sigma\) 极小，则每个坐标函数在紧无边界曲面上都是调和函数，因此为常数。这会使 \(\Sigma\) 退化为一点，矛盾。
\end{solution}

\begin{problem}
设 \(M\) 是闭、单连通 \(6\) 维流形。若
\[
H_2(M)=\mathbb Z_2,
\]
证明 Euler 示性数
\[
\chi(M)\ne -1.
\]
\end{problem}

\begin{solution}
因为 \(M\) 单连通，
\[
H_1(M;\mathbb Z)=0.
\]
题设
\[
H_2(M;\mathbb Z)=\mathbb Z_2
\]
是纯挠群。由 universal coefficient theorem，
\[
H^2(M;\mathbb Q)=0.
\]
Poincare 对偶给出
\[
H^4(M;\mathbb Q)\cong H_2(M;\mathbb Q)=0.
\]
又
\[
H^0(M;\mathbb Q)\cong H^6(M;\mathbb Q)\cong\mathbb Q,
\]
且
\[
H^1(M;\mathbb Q)=H^5(M;\mathbb Q)=0.
\]
因此有理 Betti 数只可能在 \(0,3,6\) 维非零：
\[
\chi(M)=b_0-b_1+b_2-b_3+b_4-b_5+b_6
=2-b_3.
\]
闭定向 \(6\) 维流形的中维交叉型在 \(H^3(M;\mathbb Q)\) 上是反对称非退化双线性型，因此 \(b_3\) 必为偶数。于是
\[
\chi(M)=2-b_3
\]
为偶数，不可能等于 \(-1\)。
\end{solution}

\begin{problem}
设 \((M,g)\) 是闭定向 \(n\) 维 Riemann 流形。令 \(p\in M\)，\(\operatorname{Ric}_p\) 为 \(p\) 处 Ricci 曲率张量，并定义归一化 scalar curvature
\[
S_p=\frac1n\operatorname{Tr}_g(\operatorname{Ric}_p).
\]
证明
\[
S_p=\frac1{\omega_{n-1}}\int_{S^{n-1}}\operatorname{Ric}_p(V,V)\,dS^{n-1},
\]
其中 \(S^{n-1}\subset T_pM\) 是单位球面，\(\omega_{n-1}\) 为其面积，\(V\in S^{n-1}\)，\(dS^{n-1}\) 为面积元。
\end{problem}

\begin{solution}
取 \(T_pM\) 的正交基 \(e_1,\ldots,e_n\)。设
\[
V=\sum_{i=1}^n v_i e_i,\qquad \sum_i v_i^2=1.
\]
则
\[
\operatorname{Ric}(V,V)=\sum_{i,j}\operatorname{Ric}_{ij}v_iv_j.
\]
对单位球面积分。由球面对称性，
\[
\int_{S^{n-1}}v_iv_j\,dS=0\quad (i\ne j),
\]
并且
\[
\int_{S^{n-1}}v_i^2\,dS
\]
与 \(i\) 无关。把 \(i=1,\ldots,n\) 相加得
\[
\sum_i\int_{S^{n-1}}v_i^2\,dS
=\int_{S^{n-1}}\sum_i v_i^2\,dS
=\omega_{n-1}.
\]
所以
\[
\int_{S^{n-1}}v_i^2\,dS=\frac{\omega_{n-1}}n.
\]
因此
\[
\frac1{\omega_{n-1}}\int_{S^{n-1}}\operatorname{Ric}(V,V)\,dS
=\frac1n\sum_i\operatorname{Ric}_{ii}
=\frac1n\operatorname{Tr}_g(\operatorname{Ric})
=S_p.
\]
\end{solution}

\begin{problem}
设 \(S^n\) 为 \(n\ge2\) 的球面，有限群 \(G\) 自由作用在 \(S^n\) 上，且 \(G\) 非平凡。
\begin{enumerate}[label=(\roman*)]
\item 计算商空间 \(S^n/G\) 的同伦群 \(\pi_i(S^n/G)\)，\(0\le i\le n\)。
\item 若 \(n\) 为偶数，证明 \(G\cong\mathbb Z_2\)。
\item 若 \(n\) 为奇数，证明 \(G\) 不可能同构于 \(\mathbb Z_p\times\mathbb Z_p\)，其中 \(p\) 是素数。
\end{enumerate}
\end{problem}

\begin{solution}
\begin{enumerate}[label=(\roman*)]
\item 覆盖映射
\[
S^n\to S^n/G
\]
的 deck transformation 群为 \(G\)。因 \(S^n\) 在 \(n\ge2\) 时单连通，这是 universal cover。因此
\[
\pi_1(S^n/G)\cong G.
\]
对 \(i\ge2\)，覆盖映射诱导同构
\[
\pi_i(S^n)\cong\pi_i(S^n/G).
\]
所以
\[
\pi_i(S^n/G)=\pi_i(S^n),\qquad 2\le i\le n.
\]
具体地
\[
\pi_i(S^n/G)=0\quad(2\le i<n),
\qquad
\pi_n(S^n/G)\cong\mathbb Z.
\]
并且 \(\pi_0\) 平凡，因为商空间连通。

\item 若 \(n\) 为偶数，任意自由作用元素 \(g\ne e\) 给出无不动点自同胚
\[
g:S^n\to S^n.
\]
Lefschetz 数为
\[
L(g)=1+(-1)^n\deg(g).
\]
自由作用无不动点，所以 \(L(g)=0\)。当 \(n\) 偶数时，
\[
1+\deg(g)=0,
\]
故
\[
\deg(g)=-1.
\]
于是每个非平凡元素都反定向。两个反定向映射复合为保定向映射；若群中有两个不同非平凡元素 \(g,h\)，则 \(gh\) 为非平凡保定向元素，其 degree 为 \(+1\)，Lefschetz 数为 \(2\ne0\)，必有不动点，矛盾。因此 \(G\) 只有一个非平凡元素，
\[
G\cong\mathbb Z_2.
\]

\item 若 \(n\) 为奇数，\(S^n/G\) 是一个 spherical space form。有限群自由作用在奇维球面上时，每个素数 \(p\) 的 \(p\)-子群必须具有周期上同调；特别地，不能含有 \(\mathbb Z_p\times\mathbb Z_p\)。因为
\[
H^*(\mathbb Z_p\times\mathbb Z_p;\mathbb Z_p)
\]
不是周期的，而自由球面作用会通过 Cartan--Leray 谱序列迫使群上同调周期为 \(n+1\)。因此
\[
G\not\cong \mathbb Z_p\times\mathbb Z_p.
\]
\end{enumerate}
\end{solution}

\begin{problem}
设 \(M\) 是闭定向 Riemann 流形，\(g_t\) 是随 \(t\in(-\varepsilon,\varepsilon)\) 光滑变化的一族 Riemann 度量。假设存在随 \(t\) 光滑变化的特征函数 \(f_t\) 和特征值 \(\lambda_t\)，满足
\[
\Delta_{g_t}f_t=\lambda_t f_t.
\]
设 \(f_0\) 非常值。记
\[
\dot\lambda=\left.\frac d{dt}\right|_{t=0}\lambda_t,\qquad
\dot\Delta=\left.\frac d{dt}\right|_{t=0}\Delta_{g_t}.
\]
证明：
\begin{enumerate}[label=(\roman*)]
\item 令
\[
V_{\lambda_0}=\ker(\Delta_{g_0}-\lambda_0)
\]
为 \(\lambda_0\) 的特征空间，\(\Pi:L^2(M,g_0)\to V_{\lambda_0}\) 为正交投影。证明 \(\dot\lambda\) 是算子
\[
\Pi\circ\dot\Delta:V_{\lambda_0}\to V_{\lambda_0}
\]
的特征值。
\item 若 \(\varphi_t:M\to M\) 是一参数微分同胚族，且
\[
g_t=\varphi_t^*g_0,
\]
证明 \(\dot\lambda=0\)。
\end{enumerate}
\end{problem}

\begin{solution}
\begin{enumerate}[label=(\roman*)]
\item 对特征方程求导：
\[
\dot\Delta f_0+\Delta_0\dot f
=\dot\lambda f_0+\lambda_0\dot f.
\]
整理得
\[
(\Delta_0-\lambda_0)\dot f
=\dot\lambda f_0-\dot\Delta f_0.
\]
对两边施加正交投影 \(\Pi\)。由于 \(\Delta_0\) 自伴，\(\operatorname{Im}(\Delta_0-\lambda_0)\) 与 \(V_{\lambda_0}\) 正交，所以
\[
\Pi(\Delta_0-\lambda_0)\dot f=0.
\]
于是
\[
0=\dot\lambda f_0-\Pi(\dot\Delta f_0).
\]
即
\[
\Pi\dot\Delta(f_0)=\dot\lambda f_0.
\]
因为 \(f_0\ne0\) 且 \(f_0\in V_{\lambda_0}\)，\(\dot\lambda\) 是 \(\Pi\circ\dot\Delta\) 在 \(V_{\lambda_0}\) 上的特征值。

\item 若 \(g_t=\varphi_t^*g_0\)，则 Laplacian 与拉回交换：
\[
\Delta_{g_t}(u\circ\varphi_t)=
(\Delta_{g_0}u)\circ\varphi_t.
\]
因此 \(g_t\) 与 \(g_0\) 的 Laplace 谱完全相同。既然 \(\lambda_t\) 是一条光滑选取的特征值分支，它只能在同一谱值中保持常数，所以
\[
\dot\lambda=0.
\]
也可直接令 \(\tilde f_t=f_0\circ\varphi_t\)，则
\[
\Delta_{g_t}\tilde f_t
=\lambda_0\tilde f_t,
\]
说明对应特征值不变。
\end{enumerate}
\end{solution}

\subsection{2025 年笔试题：几何与拓扑}

原题要求：共 6 题。

\begin{problem}
设 \(n>1\)。
\begin{enumerate}[label=(\roman*)]
\item \(S^{2n}\times S^{2n}\) 的切丛是否平凡？说明理由。
\item \(S^{2n}\times S^{2n-1}\) 的切丛是否平凡？说明理由。
\end{enumerate}
\end{problem}

\begin{solution}
使用事实：
\[
T(S^a\times S^b)=\pi_1^*TS^a\oplus\pi_2^*TS^b.
\]
球面 \(S^m\) 的切丛稳定平凡：
\[
TS^m\oplus\varepsilon^1\cong\varepsilon^{m+1}.
\]
并且 \(S^m\) 存在处处非零切向量场当且仅当 \(m\) 为奇数。

\begin{enumerate}[label=(\roman*)]
\item \(S^{2n}\times S^{2n}\) 的 Euler 示性数为
\[
\chi(S^{2n}\times S^{2n})
=\chi(S^{2n})^2=2^2=4.
\]
若切丛平凡，则存在处处非零向量场，从而闭流形 Euler 示性数应为 \(0\)。矛盾。因此
\[
T(S^{2n}\times S^{2n})
\]
不平凡。

\item 因为 \(2n-1\) 为奇数，\(S^{2n-1}\) 有处处非零切向量场。由稳定平凡化
\[
T(S^{2n}\times S^{2n-1})\oplus\varepsilon^2
\cong\varepsilon^{4n+1},
\]
可用 \(S^{2n-1}\) 上的非零切向量场消去一个平凡方向，从而得到 \(4n-1\) 个处处线性无关切向量场。故
\[
T(S^{2n}\times S^{2n-1})
\]
平凡。
\end{enumerate}
\end{solution}

\begin{problem}
Lie 群 \(SU(2)\) 微分同胚于 \(3\)-球面 \(S^3\)，其 Lie 代数由
\[
X_1=\begin{pmatrix}i&0\\0&-i\end{pmatrix},\quad
X_2=\begin{pmatrix}0&1\\-1&0\end{pmatrix},\quad
X_3=\begin{pmatrix}0&i\\ i&0\end{pmatrix}
\]
张成。Lie 代数上的正定内积确定 Lie 群上的左不变 Riemann 度量。Berger 度量 \(g_\varepsilon\) 由如下内积给出：
\[
\varepsilon^{-1}X_1,\quad X_2,\quad X_3
\]
构成标准正交标架，其中 \(\varepsilon>0\)。
\begin{enumerate}[label=(\roman*)]
\item 设 \(\nabla\) 是 \(g_\varepsilon\) 的 Levi-Civita 联络。证明
\[
\nabla_{X_i}X_i=0,\qquad i=1,2,3.
\]
\item 计算 \(g_\varepsilon\) 的 scalar curvature。
\end{enumerate}
\end{problem}

\begin{solution}
令
\[
e_1=\varepsilon^{-1}X_1,\qquad e_2=X_2,\qquad e_3=X_3.
\]
则 \(\{e_1,e_2,e_3\}\) 是左不变正交标架。由矩阵计算可得
\[
[X_1,X_2]=2X_3,\qquad [X_2,X_3]=2X_1,\qquad [X_3,X_1]=2X_2.
\]
因此
\[
[e_2,e_3]=2\varepsilon e_1,\qquad
[e_3,e_1]=\frac2\varepsilon e_2,\qquad
[e_1,e_2]=\frac2\varepsilon e_3.
\]

\begin{enumerate}[label=(\roman*)]
\item 对左不变度量，Koszul 公式化为
\[
2\langle\nabla_UV,W\rangle
=\langle[U,V],W\rangle-\langle[V,W],U\rangle+\langle[W,U],V\rangle.
\]
取 \(U=V=X_i\)，则 \([X_i,X_i]=0\)，且
\[
2\langle\nabla_{X_i}X_i,W\rangle
=-\langle[X_i,W],X_i\rangle+\langle[W,X_i],X_i\rangle=0
\]
对任意左不变 \(W\) 成立。故
\[
\nabla_{X_i}X_i=0.
\]

\item 对三维 unimodular Lie 群，若正交基满足
\[
[e_2,e_3]=\lambda_1e_1,\quad
[e_3,e_1]=\lambda_2e_2,\quad
[e_1,e_2]=\lambda_3e_3,
\]
则截面曲率可由 Koszul 公式计算。这里
\[
\lambda_1=2\varepsilon,\qquad \lambda_2=\lambda_3=\frac2\varepsilon.
\]
直接代入得到
\[
K(e_1,e_2)=\varepsilon^2,\qquad
K(e_1,e_3)=\varepsilon^2,\qquad
K(e_2,e_3)=4-3\varepsilon^2.
\]
三维 scalar curvature 为两倍的三个截面曲率之和：
\[
\operatorname{Scal}
=2\bigl(K_{12}+K_{13}+K_{23}\bigr)
=2\bigl(\varepsilon^2+\varepsilon^2+4-3\varepsilon^2\bigr)
=8-2\varepsilon^2.
\]
当 \(\varepsilon=1\) 时得到 \(\operatorname{Scal}=6\)，与单位 \(S^3\) 的标准曲率一致。
\end{enumerate}
\end{solution}

\begin{problem}
设 \(n\ge1\)，\(U(n)\) 为 \(\C^n\) 中的酉矩阵群。
\begin{enumerate}[label=(\roman*)]
\item 计算 \(\pi_2(U(n))\) 和 \(\pi_3(U(n))\)。
\item 证明行列式映射
\[
\det:U(n)\to S^1
\]
在基本群上诱导同构。
\end{enumerate}
\end{problem}

\begin{solution}
\begin{enumerate}[label=(\roman*)]
\item 对 \(n=1\)，\(U(1)=S^1\)，所以
\[
\pi_2(U(1))=0,\qquad \pi_3(U(1))=0.
\]
对 \(n\ge2\)，有纤维丛
\[
U(n-1)\longrightarrow U(n)\longrightarrow S^{2n-1},
\]
其中 \(U(n)\to S^{2n-1}\) 取矩阵的第一列。由长正合同伦序列和稳定范围可得
\[
\pi_2(U(n))=0.
\]
当 \(n\ge2\) 时，
\[
\pi_3(U(n))\cong\mathbb Z.
\]
其中 \(n=2\) 可由 \(SU(2)\cong S^3\) 与 \(U(2)\simeq S^1\times SU(2)\) 局部乘积看出；更高 \(n\) 由上述纤维丛归纳得到。

\item 有分解
\[
1\to SU(n)\to U(n)\xrightarrow{\det}S^1\to1.
\]
这给出纤维丛
\[
SU(n)\to U(n)\to S^1.
\]
对 \(n\ge2\)，\(SU(n)\) 单连通；对 \(n=1\)，结论显然。长正合同伦序列给出
\[
\pi_1(U(n))\xrightarrow{\det_*}\pi_1(S^1)
\]
为同构。因此
\[
\det_*:\pi_1(U(n))\to\mathbb Z
\]
是同构。
\end{enumerate}
\end{solution}

\begin{problem}
证明 \(\mathbb{CP}^2\) 不存在到 \(\R^6\) 的 immersion。
\end{problem}

\begin{solution}
反设存在 immersion
\[
f:\mathbb{CP}^2\looparrowright\R^6.
\]
则有法丛 \(\nu\)，秩为 \(2\)，满足
\[
T\mathbb{CP}^2\oplus\nu\cong\varepsilon^6.
\]
取 Stiefel--Whitney 类：
\[
w(T\mathbb{CP}^2)w(\nu)=1.
\]
设 \(a\in H^2(\mathbb{CP}^2;\mathbb Z_2)\) 为生成元。复切丛总 Chern 类为
\[
c(T\mathbb{CP}^2)=(1+h)^3
\]
截断到次数 \(4\)，故
\[
c_1=3h,\qquad c_2=3h^2.
\]
模 \(2\) 化得
\[
w_2(T\mathbb{CP}^2)=a,\qquad w_4(T\mathbb{CP}^2)=a^2.
\]
所以
\[
w(T\mathbb{CP}^2)=1+a+a^2.
\]
于是
\[
w(\nu)=\frac1{1+a+a^2}=1+a
\]
在 \(H^*(\mathbb{CP}^2;\mathbb Z_2)\) 中成立，因为
\[
(1+a+a^2)(1+a)=1+a^3=1.
\]
但 \(\nu\) 是秩 \(2\) 的实向量丛，故其 Stiefel--Whitney 类最高只能到 \(w_2\)，即 \(w(\nu)=1+w_1+w_2\)。这里 \(a\) 是 \(2\) 维类，因此 \(w_2(\nu)=a\)，尚不矛盾。

进一步取 Pontryagin 类。由
\[
p(T\mathbb{CP}^2)p(\nu)=1.
\]
\(\mathbb{CP}^2\) 的
\[
p_1(T\mathbb{CP}^2)=c_1^2-2c_2=9h^2-6h^2=3h^2.
\]
秩 \(2\) 实向量丛 \(\nu\) 若可定向，则
\[
p_1(\nu)=e(\nu)^2.
\]
但 \(H^2(\mathbb{CP}^2;\mathbb Z)=\mathbb Z h\)，设 \(e(\nu)=kh\)，则
\[
p_1(\nu)=k^2h^2.
\]
由 \(p_1(T)+p_1(\nu)=0\) 得
\[
3+k^2=0
\]
在整数中不可能。若 \(\nu\) 不先验可定向，由于 \(\mathbb{CP}^2\) 单连通，任意实向量丛的一阶 Stiefel--Whitney 类为零，故 \(\nu\) 可定向。矛盾。因此不存在这样的 immersion。
\end{solution}

\begin{problem}
设 \((M,g)\) 是闭定向二维光滑 Riemann 流形，Gaussian 曲率为 \(K\)。设 \(u\) 是光滑函数，并定义同一共形类中的新度量
\[
\widetilde g=e^{2u}g,
\]
其 Gaussian 曲率为 \(\widetilde K\)。
\begin{enumerate}[label=(\roman*)]
\item 证明 \(u\) 满足
\[
\Delta u-K+\widetilde K e^{2u}=0.
\]
\item 若 Euler 示性数 \(\chi(M)=0\)，证明或者 \(\widetilde K\equiv0\)，或者
\[
\int_M\widetilde K e^{2f}\,d\operatorname{Vol}_g<0,
\]
其中 \(f\) 是方程
\[
\Delta f=K
\]
的解。
\end{enumerate}
\end{problem}

\begin{solution}
\begin{enumerate}[label=(\roman*)]
\item 二维共形变换曲率公式为
\[
\widetilde K=e^{-2u}(K-\Delta u),
\]
其中 \(\Delta=\operatorname{div}\nabla\)。移项即得
\[
\Delta u-K+\widetilde K e^{2u}=0.
\]

\item 因为 \(\chi(M)=0\)，Gauss--Bonnet 定理给出
\[
\int_MK\,dV_g=0.
\]
所以方程
\[
\Delta f=K
\]
满足可解条件。

由 (i) 得
\[
\widetilde K e^{2u}=K-\Delta u=\Delta f-\Delta u=\Delta(f-u).
\]
令
\[
w=f-u.
\]
则
\[
\Delta w=\widetilde K e^{2u}.
\]
乘以 \(e^{2w}\) 并积分：
\[
\int_M \widetilde K e^{2f}\,dV_g
=\int_M \widetilde K e^{2u}e^{2w}\,dV_g
=\int_M e^{2w}\Delta w\,dV_g.
\]
分部积分得
\[
\int_M e^{2w}\Delta w\,dV_g
=-2\int_M e^{2w}|\nabla w|^2\,dV_g\le0.
\]
若等号成立，则 \(\nabla w=0\)，故 \(w\) 为常数，从而
\[
\widetilde K e^{2u}=\Delta w=0,
\]
即 \(\widetilde K\equiv0\)。否则积分严格小于 \(0\)。命题得证。
\end{enumerate}
\end{solution}

\begin{problem}
设 \(n\ge1\)。
\begin{enumerate}[label=(\roman*)]
\item \(\mathbb{CP}^n\) 是否存在无不动点的自同胚？构造例子或证明不存在。
\item \(\mathbb{CP}^2\times\mathbb{CP}^2\) 是否存在无不动点的自同胚？构造例子或证明不存在。
\end{enumerate}
\end{problem}

\begin{solution}
\begin{enumerate}[label=(\roman*)]
\item \(\mathbb{CP}^n\) 存在无不动点自同胚当且仅当 \(n\) 为奇数。设 \(f:\mathbb{CP}^n\to\mathbb{CP}^n\) 是自同胚。上同调环为
\[
H^*(\mathbb{CP}^n;\mathbb Q)=\mathbb Q[h]/(h^{n+1}),\qquad \deg h=2.
\]
自同胚诱导
\[
f^*h=\pm h.
\]
Lefschetz 数为
\[
L(f)=\sum_{k=0}^n\operatorname{tr}\left(f^*:H^{2k}\to H^{2k}\right)
=\sum_{k=0}^n(\pm1)^k.
\]
若 \(f^*h=h\)，则
\[
L(f)=n+1\ne0.
\]
若 \(f^*h=-h\)，则
\[
L(f)=1-1+1-\cdots+(-1)^n.
\]
当 \(n\) 偶数时该值为 \(1\)，所以偶数维复射影空间的任意自同胚 Lefschetz 数非零，必有不动点。

若 \(n=2m+1\)，定义
\[
[z_0:z_1:\cdots:z_{2m}:z_{2m+1}]
\mapsto
[-\overline z_1:\overline z_0:\cdots:-\overline z_{2m+1}:\overline z_{2m}].
\]
若有不动点，则每一对坐标满足
\[
(-\overline z_{2j+1},\overline z_{2j})=\lambda(z_{2j},z_{2j+1})
\]
同一个 \(\lambda\)。这推出 \(|\lambda|^2=-1\)，矛盾。因此无不动点。若 \(n\) 为偶数，上述 Lefschetz 数总非零，故不存在。

\item 不存在。证明如下：\(\mathbb{CP}^2\times\mathbb{CP}^2\) 的有理上同调只在偶数维，且任意自同胚诱导的 Lefschetz 数等于其在偶维上同调迹的和。设 \(a,b\) 为两个因子的 \(H^2\) 生成元。自同胚保持上同调环和 Poincare 对偶，诱导在 \(H^2\) 上只能为
\[
a\mapsto a,\ b\mapsto b
\quad\text{或}\quad
a\mapsto b,\ b\mapsto a.
\]
在第一种情形所有偶维迹相加为
\[
\chi(\mathbb{CP}^2\times\mathbb{CP}^2)=9.
\]
第二种情形逐维计算迹：\(H^0,H^8\) 迹为 \(1,1\)，\(H^2\) 上交换矩阵迹为 \(0\)，\(H^4\) 基 \(a^2,ab,b^2\) 上迹为 \(1\)，\(H^6\) 基 \(a^2b,ab^2\) 上迹为 \(0\)。故
\[
L(F)=1+0+1+0+1=3\ne0.
\]
Lefschetz 不动点定理说明任意自同胚都有不动点。因此不存在无不动点自同胚。
\end{enumerate}
\end{solution}

\subsection{2026 年笔试题：几何与拓扑}

原题要求：解答全部题目。

\begin{problem}
设 \(S\subset\R^3\) 是无边界光滑正则曲面，\(p\in S\)。设 \(\{A_\varepsilon\}_{\varepsilon>0}\) 是 \(S\) 中一族区域，每个 \(A_\varepsilon\) 包含 \(p\)，面积为 \(|A_\varepsilon|\)，并且当 \(\varepsilon\to0\) 时收缩到 \(\{p\}\)。令
\[
N:S\to S^2
\]
为 Gauss 映射。证明
\[
|K(p)|=\lim_{\varepsilon\to0}\frac{\operatorname{Area}(N(A_\varepsilon))}{\operatorname{Area}(A_\varepsilon)}.
\]
\end{problem}

\begin{solution}
Gauss 映射的微分
\[
dN_p:T_pS\to T_{N(p)}S^2
\]
与形算子只差符号：
\[
dN_p=-S_p.
\]
因此其 Jacobian 的绝对值为
\[
|\det dN_p|=|\det S_p|=|K(p)|.
\]

面积公式给出，对充分小的区域 \(A_\varepsilon\)，
\[
\operatorname{Area}(N(A_\varepsilon))
=\int_{A_\varepsilon} |\operatorname{Jac}N(q)|\,dA(q)
\]
按重数计；当区域收缩到 \(p\) 且 \(N\) 在局部可按面积公式处理时，重数不会影响极限。于是
\[
\frac{\operatorname{Area}(N(A_\varepsilon))}{\operatorname{Area}(A_\varepsilon)}
=\frac1{\operatorname{Area}(A_\varepsilon)}
\int_{A_\varepsilon}|K(q)|\,dA(q).
\]
由于 \(K\) 连续，右侧在 \(A_\varepsilon\to\{p\}\) 时趋于
\[
|K(p)|.
\]
命题得证。
\end{solution}

\begin{problem}
设 \((M^2,g)\) 是光滑、紧、定向二维 Riemann 流形，边界非空。假设
\[
K_g\ge0
\]
且边界的测地曲率满足
\[
k_g\ge1
\]
其中 \(k_g\) 按题目约定由沿边界的外指单位法向计算。又假设
\[
\operatorname{Length}(\partial M)\ge2\pi.
\]
证明 \((M,g)\) 等距于 Euclidean 单位圆盘 \((D^2(1),g_{\mathrm{flat}})\)。
\end{problem}

\begin{solution}
按题目测地曲率符号约定，Gauss--Bonnet 公式写为
\[
\int_MK\,dA+\int_{\partial M}k_g\,ds=2\pi\chi(M).
\]
由 \(K\ge0\)、\(k_g\ge1\) 和 \(\operatorname{Length}(\partial M)\ge2\pi\)，得
\[
2\pi\chi(M)
\ge \int_{\partial M}k_g\,ds
\ge \operatorname{Length}(\partial M)\ge2\pi.
\]
所以
\[
\chi(M)\ge1.
\]
紧定向带边界曲面满足 \(\chi(M)\le1\)，等号仅当 \(M\) 为拓扑圆盘。因此
\[
\chi(M)=1,
\]
且上面所有不等式均取等号。于是
\[
K\equiv0,\qquad k_g\equiv1,\qquad \operatorname{Length}(\partial M)=2\pi.
\]

因此 \(M\) 是平坦圆盘，边界在内蕴平坦度量中是一条测地曲率恒为 \(1\)、长度为 \(2\pi\) 的闭曲线。平坦圆盘可由发展映射局部等距地展开到 Euclidean 平面；边界展开为平面中曲率恒为 \(1\)、长度为 \(2\pi\) 的简单闭曲线，故为单位圆。由 Jordan 定理和发展映射的单值性，整个 \(M\) 等距于该单位圆围成的 Euclidean 圆盘。
\end{solution}

\begin{problem}
设 \(G\) 是连通 \(n\) 维 Lie 群，配备双不变 Riemann 度量 \(\langle\cdot,\cdot\rangle\)。
\begin{enumerate}[label=(\arabic*)]
\item 证明 \(G\) 的截面曲率非负。
\item 若 \(G\) 的 Lie 代数 \(\mathfrak g\) 中心平凡：
\[
\mathfrak z(\mathfrak g)=\{X\in\mathfrak g:[X,Y]=0\text{ 对所有 }Y\in\mathfrak g\}=\{0\},
\]
证明 \(G\) 紧。
\item 若 \(G\) 单连通，证明
\[
G\cong G_0\times\R^k,
\]
其中 \(G_0\) 是单连通紧 Lie 群，其 Lie 代数中心平凡，\(\R^k\) 为加法 Lie 群。
\end{enumerate}
\end{problem}

\begin{solution}
\begin{enumerate}[label=(\arabic*)]
\item 对双不变度量，左不变向量场满足
\[
\nabla_XY=\frac12[X,Y].
\]
于是
\[
R(X,Y)Z=-\frac14[[X,Y],Z]
\]
按相应曲率号约定，截面曲率满足
\[
K(X,Y)=\frac14\frac{\|[X,Y]\|^2}{\|X\wedge Y\|^2}\ge0.
\]

\item 双不变度量意味着 \(\operatorname{ad}X\) 对内积反对称：
\[
\langle [X,Y],Z\rangle+\langle Y,[X,Z]\rangle=0.
\]
因此 Killing 型
\[
B(X,Y)=\operatorname{tr}(\operatorname{ad}X\operatorname{ad}Y)
\]
在 \([\mathfrak g,\mathfrak g]\) 上半负定。若中心为零，则 \(\mathfrak g\) 无 Abel 理想，且 Killing 型负定，故 \(\mathfrak g\) 是紧型半单 Lie 代数。对应连通 Lie 群的 universal cover 是紧半单群的有限覆盖，因此紧；\(G\) 作为其商群也紧。

\item 一般地，带正定 \(\operatorname{Ad}\)-不变内积的 Lie 代数分解为
\[
\mathfrak g=\mathfrak z(\mathfrak g)\oplus[\mathfrak g,\mathfrak g],
\]
其中 \(\mathfrak z(\mathfrak g)\cong\R^k\) 为中心，\([\mathfrak g,\mathfrak g]\) 是紧型半单 Lie 代数。若 \(G\) 单连通，则 Lie 代数直和积分为 Lie 群直积：
\[
G\cong \R^k\times G_0,
\]
其中 \(G_0\) 是 \([\mathfrak g,\mathfrak g]\) 对应的单连通紧半单 Lie 群。半单 Lie 代数中心平凡，所以 \(G_0\) 的 Lie 代数中心平凡。
\end{enumerate}
\end{solution}

\begin{problem}
设 \(S^{2026}\) 是 \(\R^{2027}\) 中单位球面。
\begin{enumerate}[label=(\arabic*)]
\item 用 \(\mathbb Z_2\) 系数计算切丛
\[
TS^{2026}\to S^{2026}
\]
的 Euler 类。
\item 设
\[
US^{2026}\to S^{2026}
\]
为单位切球丛，纤维为 \(S^{2025}\)。计算 \(\mathbb Z_2\) 系数 Poincare 级数
\[
P_t(US^{2026};\mathbb Z_2)=\sum_{i\ge0}b_i(US^{2026};\mathbb Z_2)t^i.
\]
\item 设 \(p,q\in S^{2026}\) 为两个非对径点。考虑从 \(p\) 到 \(q\) 的分段光滑路径空间
\[
\Omega(S^{2026};p,q)=\{\gamma:[0,1]\to S^{2026}:\gamma(0)=p,\gamma(1)=q\}
\]
及能量泛函
\[
E(\gamma)=\frac12\int_0^1|\dot\gamma(t)|^2\,dt.
\]
\begin{enumerate}[label=(\alph*)]
\item 求 \(E\) 的所有临界点。
\item 设 \(\lambda(E,\gamma)\) 为临界点处 Hessian 的负定最大子空间维数，求 \(\lambda(E,\gamma)\)。
\item 求基点 loop space \(\Omega(S^{2026},p)\) 的同伦型。
\item 用 \(\mathbb Z_2\) 系数计算 \(\Omega(S^{2026},p)\) 和自由 loop space \(\Lambda S^{2026}\) 的 Poincare 级数。
\end{enumerate}
\end{enumerate}
\end{problem}

\begin{solution}
记 \(m=2026\)。
\begin{enumerate}[label=(\arabic*)]
\item 模 \(2\) Euler 类等于最高 Stiefel--Whitney 类。其在基本类上的取值为
\[
\chi(S^m)\mod2.
\]
因 \(m\) 偶，\(\chi(S^m)=2\equiv0\pmod2\)，而 \(H^m(S^m;\mathbb Z_2)\cong\mathbb Z_2\)，故
\[
e(TS^{2026};\mathbb Z_2)=0.
\]

\item 对球丛
\[
S^{m-1}\to US^m\to S^m
\]
使用 Gysin 序列。由于模 \(2\) Euler 类为零，序列分裂，得到非零 Betti 数位于
\[
0,\quad m-1,\quad m,\quad 2m-1.
\]
因此
\[
P_t(US^{2026};\mathbb Z_2)
=(1+t^{2025})(1+t^{2026})
=1+t^{2025}+t^{2026}+t^{4051}.
\]

\item 令 \(\theta=d(p,q)\in(0,\pi)\)。
\begin{enumerate}[label=(\alph*)]
\item 临界点正是从 \(p\) 到 \(q\) 的常速测地线。由于 \(p,q\) 非对径，它们决定唯一大圆。临界测地线分两族：
\[
\gamma_k^+:\ \text{沿较短方向，长度 }\theta+2\pi k,
\]
\[
\gamma_k^-:\ \text{沿较长方向，长度 }(2\pi-\theta)+2\pi k,
\]
其中 \(k=0,1,2,\ldots\)。

\item 球面测地线上的共轭点每隔 \(\pi\) 出现一次，每次重数为 \(m-1=2025\)。端点非共轭，所以 Morse index 等于内部共轭点重数之和：
\[
\lambda(E,\gamma_k^+)=2025\cdot 2k,
\]
\[
\lambda(E,\gamma_k^-)=2025\cdot(2k+1).
\]

\item 对任意 \(q\) 选定一条从 \(p\) 到 \(q\) 的路径，可把 \(\Omega(S^m;p,q)\) 与基点 loop space \(\Omega S^m\) 同伦等价。经典 James 分解说明
\[
\Omega S^m\simeq J(S^{m-1}),
\]
即 \(S^{m-1}\) 的 James reduced product。这里
\[
\Omega(S^{2026},p)\simeq J(S^{2025}).
\]

\item 以 \(\mathbb Z_2\) 为系数，
\[
H_*(\Omega S^m;\mathbb Z_2)
\]
是一个由次数 \(m-1\) 元生成的张量代数；每个次数 \(j(m-1)\) 上一维。因此
\[
P_t(\Omega S^{2026};\mathbb Z_2)
=\frac1{1-t^{2025}}.
\]
自由 loop space 有纤维化
\[
\Omega S^m\to \Lambda S^m\to S^m
\]
并有常值 loop 给出的截面。以 \(\mathbb Z_2\) 系数，Poincare 级数为
\[
P_t(\Lambda S^{2026};\mathbb Z_2)
=\frac{1+t^{2026}}{1-t^{2025}}.
\]
\end{enumerate}
\end{enumerate}
\end{solution}

\begin{problem}
设 \(M^n\) 是连通、闭、光滑、aspherical 流形，\(n\ge1\)，即 universal cover \(\widetilde M\) 可缩。
\begin{enumerate}[label=(\arabic*)]
\item 证明 \(\widetilde M\) 非紧。
\item 证明 \(\pi_1(M)\) 的每个非平凡元素都有无限阶。
\end{enumerate}
\end{problem}

\begin{solution}
\begin{enumerate}[label=(\arabic*)]
\item 若 \(\widetilde M\) 紧，则覆盖
\[
\widetilde M\to M
\]
为有限覆盖，故 \(\widetilde M\) 是闭 \(n\) 维流形。闭连通 \(n\) 维流形有
\[
H_n(\widetilde M;\mathbb Z_2)\cong\mathbb Z_2.
\]
但 \(\widetilde M\) 可缩时，其正维约化同调全为零，矛盾。因此 \(\widetilde M\) 非紧。

\item \(\pi_1(M)\) 作为 deck transformation 群自由作用在可缩流形 \(\widetilde M\) 上。若存在非平凡有限阶元素 \(g\in\pi_1(M)\)，则它生成有限循环群 \(C\)，并自由作用在 \(\widetilde M\) 上。

由 Smith 固定点定理，有限 \(p\)-群作用在有限维可缩流形上必有非空固定点集。取 \(g\) 的某个素数阶幂，得到非平凡 \(p\)-群作用，于是应有固定点。这与 deck transformation 的自由作用矛盾。故 \(\pi_1(M)\) 中无非平凡有限阶元素，每个非平凡元素都有无限阶。
\end{enumerate}
\end{solution}

\begin{problem}
设 \((M^n,g)\) 是 \(n\ge2\) 维完备 Riemann 流形。假设存在
\[
\varepsilon>\frac14
\]
使得当 \(r(p,x)=d(p,x)\) 充分大时
\[
\operatorname{Ric}_g(x)\ge
\varepsilon(n-1)r(p,x)^{-2}.
\]
\begin{enumerate}[label=(\arabic*)]
\item 证明 \(M\) 必紧。
\item 举例说明若 \(\varepsilon\le\frac14\)，结论可能失败。
\end{enumerate}
\end{problem}

\begin{solution}
\begin{enumerate}[label=(\arabic*)]
\item 反设 \(M\) 非紧。由 Hopf--Rinow 定理，存在从 \(p\) 出发的 ray
\[
\gamma:[0,\infty)\to M.
\]
ray 的任意有限段都是最短测地线，因此其 index form 对所有端点为零的法向变分非负。取平行单位法向场 \(E_1,\ldots,E_{n-1}\)，并取紧支撑函数 \(\phi\) 于 \([R,L]\)，则
\[
0\le
\sum_{i=1}^{n-1}I(\phi E_i,\phi E_i)
=\int_R^L\left((n-1)(\phi')^2-\phi^2\operatorname{Ric}(\gamma',\gamma')\right)\,dt.
\]
当 \(R\) 足够大时，
\[
\operatorname{Ric}(\gamma',\gamma')\ge \varepsilon(n-1)t^{-2}.
\]
于是
\[
\int_R^L(\phi')^2\,dt
\ge
\varepsilon\int_R^L\frac{\phi^2}{t^2}\,dt.
\]
但 Hardy 不等式的最佳常数为 \(1/4\)：当 \(\varepsilon>1/4\) 时，可以选择支撑在足够长区间上的 \(\phi\)，使
\[
\int(\phi')^2<\varepsilon\int\frac{\phi^2}{t^2}.
\]
矛盾。因此不存在 ray，完备非紧流形不可能满足该 Ricci 下界。故 \(M\) 紧。

\item 对任意 \(0<\varepsilon\le1/4\)，取 \(a\in(0,1)\) 使
\[
a(1-a)\ge\varepsilon.
\]
在 \([1,\infty)\times S^{n-1}\) 上取 warped product 度量
\[
g=dr^2+r^{2a}g_{S^{n-1}},
\]
并在 \(r\le2\) 内光滑封口得到完备非紧流形。径向 Ricci 曲率为
\[
\operatorname{Ric}(\partial_r,\partial_r)
=-(n-1)\frac{f''}{f}
=(n-1)a(1-a)r^{-2},
\qquad f(r)=r^a.
\]
切向 Ricci 项对大 \(r\) 也满足同阶或更强的非负下界。因此当 \(r\) 足够大时
\[
\operatorname{Ric}\ge \varepsilon(n-1)r^{-2}
\]
成立，但流形非紧。这说明 \(\varepsilon\le1/4\) 时结论可能失败。
\end{enumerate}
\end{solution}

\section{团体赛题}

\subsection{2010 年团体赛：几何与拓扑}

原题要求：从以下 6 道题中选择 5 道作答。

\begin{problem}
设 \(S^n\subset\R^{n+1}\) 为单位球面，\(\R^n\subset\R^{n+1}\) 为过球心的赤道 \(n\)-平面，\(N\) 为 \(S^n\) 的北极。定义 stereographic projection
\[
\pi:S^n\setminus\{N\}\to\R^n
\]
为：把 \(A\in S^n\setminus\{N\}\) 送到直线 \(AN\) 与赤道平面的交点。证明 stereographic projection 是共形变换，并由此推出 \(S^n\) 的标准度量在该投影下的表达式。
\end{problem}

\begin{solution}
取
\[
S^n=\{(X,t)\in\R^n\times\R:|X|^2+t^2=1\},
\qquad N=(0,1).
\]
从北极投影到平面 \(t=0\)。若投影坐标为 \(x\in\R^n\)，逆映射为
\[
\sigma(x)=\left(\frac{2x}{1+|x|^2},\frac{|x|^2-1}{1+|x|^2}\right).
\]
直接计算微分。对 \(v\in T_x\R^n\cong\R^n\)，有
\[
d\sigma_x(v)=
\left(
\frac{2(1+|x|^2)v-4(x\cdot v)x}{(1+|x|^2)^2},
\frac{4x\cdot v}{(1+|x|^2)^2}
\right).
\]
计算 Euclidean 内积得
\[
\langle d\sigma_x(v),d\sigma_x(w)\rangle_{\R^{n+1}}
=\frac{4}{(1+|x|^2)^2}\langle v,w\rangle_{\R^n}.
\]
因此球面标准度量拉回为
\[
\sigma^*g_{S^n}
=\frac{4}{(1+|x|^2)^2}
\sum_{i=1}^n (dx^i)^2.
\]
它是 Euclidean 度量乘以正函数，故 stereographic projection 共形。
\end{solution}

\begin{problem}
设 \(M\) 是连通二维 Riemann 流形。设 \(f\) 是 \(M\) 上光滑非常值函数，且 \(f\) 有上界并满足
\[
\Delta f\ge0
\]
处处成立。证明不存在点 \(p\in M\)，使
\[
f(p)=\sup\{f(x):x\in M\}.
\]
\end{problem}

\begin{solution}
若存在 \(p\in M\) 使 \(f(p)=\sup_M f\)，则 \(f\) 在内点 \(p\) 取得最大值。由于
\[
\Delta f\ge0,
\]
\(f\) 是 subharmonic 函数。强最大值原理说明：在连通流形上，subharmonic 函数若在内点取得最大值，则它在整个连通分支上为常数。

这里 \(M\) 连通，因此 \(f\) 必为常数。这与题设 \(f\) 非常值矛盾。故不存在这样的 \(p\)。
\end{solution}

\begin{problem}
设 \(M\) 是 \(d\) 维紧光滑流形。证明存在某个正整数 \(n\)，使得 \(M\) 可正则嵌入 Euclidean 空间 \(\R^n\)。
\end{problem}

\begin{solution}
这是 Whitney 嵌入定理的基本结论。给出构造性证明思路。

由于 \(M\) 紧，可取有限坐标覆盖
\[
(U_i,\varphi_i),\qquad i=1,\ldots,N.
\]
取从属于该覆盖的光滑单位分解
\[
\rho_i,\qquad \operatorname{supp}\rho_i\subset U_i,\qquad \sum_i\rho_i=1.
\]
定义映射
\[
F:M\to\R^{N(d+1)}
\]
为
\[
F(p)=\bigl(\rho_i(p),\rho_i(p)\varphi_i^1(p),\ldots,\rho_i(p)\varphi_i^d(p)\bigr)_{i=1}^N.
\]
若 \(p\ne q\)，可取某个 \(i\) 使 \(\rho_i(p)>0\) 并在小邻域内利用坐标分量区分点；必要时把覆盖和单位分解细化，可使 \(F\) 为一一映射。并且在 \(\rho_i(p)>0\) 的坐标片上，\(F\) 的分量包含坐标函数的非退化倍数，所以 \(dF_p\) 单射。

因此 \(F\) 是 immersion 且为拓扑嵌入。紧空间到 Hausdorff 空间的一一连续映射是同胚到其像，故 \(F\) 是光滑嵌入。于是 \(M\) 可嵌入某个 \(\R^n\)。
\end{solution}

\begin{problem}
证明紧无边界光滑流形 \(M\) 上任意 \(C^\infty\) 函数 \(f\) 至少有两个临界点。当 \(M\) 是二维环面时，证明 \(f\) 必有多于两个临界点。
\end{problem}

\begin{solution}
因为 \(M\) 紧，\(f\) 取得最大值和最小值。无边界流形上，最大点与最小点都是临界点。若 \(f\) 为常数，则所有点都是临界点；若非常数，最大点和最小点不同，因此至少两个临界点。

当 \(M=T^2\) 时，若 \(f\) 只有两个临界点，则一个为最大点、一个为最小点。一般可微扰为 Morse 函数而不增加临界点数；Morse 不等式要求临界点数至少为 Betti 数之和：
\[
b_0(T^2)+b_1(T^2)+b_2(T^2)=1+2+1=4.
\]
更直接地，任何 Morse 函数在环面上至少有一个 index \(0\) 临界点、两个 index \(1\) 临界点和一个 index \(2\) 临界点。故不可能只有两个临界点。对非 Morse 函数，可作任意小扰动得到 Morse 函数；若原函数临界点不超过两个，则扰动也不能产生满足 Morse 不等式所需的临界结构，矛盾。故 \(T^2\) 上任意光滑函数必有多于两个临界点。
\end{solution}

\begin{problem}
构造空间 \(X\)，使
\[
H_0(X)=\mathbb Z,\qquad H_1(X)=\mathbb Z_2\times\mathbb Z_3,\qquad H_2(X)=\mathbb Z,
\]
且所有其他同调群为零。
\end{problem}

\begin{solution}
取 Moore 空间
\[
M(\mathbb Z_2,1),\qquad M(\mathbb Z_3,1)
\]
以及 \(S^2\) 的楔和：
\[
X=M(\mathbb Z_2,1)\vee M(\mathbb Z_3,1)\vee S^2.
\]
其中 \(M(\mathbb Z_m,1)\) 可由一个 \(S^1\) 贴上一个 \(2\)-胞腔构造，贴映射为 \(m\) 次映射：
\[
S^1\xrightarrow{m}S^1.
\]
其约化同调为
\[
\widetilde H_1(M(\mathbb Z_m,1))=\mathbb Z_m,
\]
其他约化同调为零。

楔和的约化同调为各部分约化同调直和，因此
\[
H_0(X)=\mathbb Z,
\]
\[
H_1(X)=\mathbb Z_2\oplus\mathbb Z_3,
\]
\[
H_2(X)=H_2(S^2)=\mathbb Z,
\]
且其他同调群为零。
\end{solution}

\begin{problem}
\begin{enumerate}[label=(\alph*)]
\item 定义 \(C^\infty\) 映射 \(f:S^2\to S^2\) 的度 \(\deg f\)，并证明该定义良好且与定义中的选择无关。
\item 详细证明：对每个整数 \(k\)，包括负整数，存在 \(C^\infty\) 映射 \(f:S^2\to S^2\)，其度为 \(k\)。
\end{enumerate}
\end{problem}

\begin{solution}
\begin{enumerate}[label=(\alph*)]
\item 取 \(S^2\) 上的面积形式 \(\omega\)，归一化为
\[
\int_{S^2}\omega=1.
\]
定义
\[
\deg f=\int_{S^2}f^*\omega.
\]
若换取另一归一化面积形式 \(\omega'\)，则
\[
\int_{S^2}(\omega-\omega')=0.
\]
由于 \(H^2(S^2;\R)\cong\R\)，这说明
\[
\omega-\omega'=d\eta.
\]
于是
\[
\int_{S^2}f^*(\omega-\omega')
=\int_{S^2}d(f^*\eta)=0.
\]
故定义与 \(\omega\) 的选择无关。该数等于 \(f_*[S^2]\) 在 \(H_2(S^2;\mathbb Z)\cong\mathbb Z\) 上的倍数，因此为整数，并且与光滑同伦不变。

\item 对 \(k\ge0\)，把 \(S^2\) 视为 Riemann sphere \(\mathbb{CP}^1\)。映射
\[
f_k([z_0:z_1])=[z_0^k:z_1^k]
\]
在 \(k>0\) 时是全纯映射，degree 为 \(k\)；\(k=0\) 时取常值映射，degree 为 \(0\)。

对 \(k<0\)，取 degree 为 \(|k|\) 的映射 \(f_{|k|}\)，再复合一个反定向微分同胚，例如复共轭
\[
c([z_0:z_1])=[\overline z_0:\overline z_1].
\]
则
\[
\deg(c\circ f_{|k|})=-|k|=k.
\]
因此任意整数都可实现为某个光滑映射 \(S^2\to S^2\) 的度。
\end{enumerate}
\end{solution}

\subsection{2011 年团体赛：几何与拓扑}

原题要求：从以下 6 道题中选择 5 道作答。

\begin{problem}
设 \(K\) 是有限连通单纯复形。判断真伪并说明理由：
\begin{enumerate}[label=(\alph*)]
\item 若 \(\pi_1(K)\) 有限，则 \(K\) 的 universal cover 紧。
\item 若 \(K\) 的 universal cover 紧，则 \(\pi_1(K)\) 有限。
\end{enumerate}
\end{problem}

\begin{solution}
\begin{enumerate}[label=(\alph*)]
\item 真。有限连通单纯复形 \(K\) 紧。若 \(\pi_1(K)\) 有限，则 universal cover
\[
\widetilde K\to K
\]
是有限层覆盖。紧空间的有限覆盖仍紧，因此 \(\widetilde K\) 紧。

\item 真。覆盖的层数等于 \(\pi_1(K)\) 的基数。若 \(\widetilde K\) 紧而 \(K\) 是有限复形，则取 \(K\) 中一个均匀覆盖的小开集 \(U\)。其原像是若干个互不相交的开集，每个同胚于 \(U\)。若层数无限，则 \(\pi^{-1}(U)\) 含无限个互不相交开集；选取其中点列，不可能有收敛子列落在 Hausdorff 覆盖空间中。这与 \(\widetilde K\) 紧矛盾。因此层数有限，\(\pi_1(K)\) 有限。
\end{enumerate}
\end{solution}

\begin{problem}
计算 \(n\)-单纯形的 \(m\)-骨架的所有同调群，其中 \(0\le m\le n\)。
\end{problem}

\begin{solution}
记 \(K=\operatorname{sk}_m(\Delta^n)\)。当 \(m=0\) 时，\(K\) 是 \(n+1\) 个离散点，所以
\[
H_0(K)\cong\mathbb Z^{n+1},\qquad H_i(K)=0\ (i>0).
\]

设 \(1\le m<n\)。\(K\) 是 \((m-1)\)-连通的，并且只有到维数 \(m\) 的胞腔，因此除 \(H_0\) 与 \(H_m\) 外同调为零。Euler 示性数为
\[
\chi(K)=\sum_{j=0}^m(-1)^j\binom{n+1}{j+1}.
\]
若
\[
H_m(K)\cong\mathbb Z^r,
\]
则
\[
\chi(K)=1+(-1)^m r.
\]
组合恒等式
\[
\sum_{j=0}^m(-1)^j\binom{n+1}{j+1}
=1+(-1)^m\binom{n}{m+1}
\]
给出
\[
r=\binom{n}{m+1}.
\]
因此当 \(1\le m<n\) 时，
\[
H_0(K)=\mathbb Z,\qquad
H_m(K)=\mathbb Z^{\binom{n}{m+1}},
\]
其他同调群为零。

当 \(m=n\) 时，骨架就是整个 \(n\)-单纯形，可缩，所以
\[
H_0(K)=\mathbb Z,\qquad H_i(K)=0\ (i>0).
\]
\end{solution}

\begin{problem}
设 \(M\) 是带边界的 \(n\) 维紧定向 Riemann 流形，\(X\) 是 \(M\) 上光滑向量场。若 \(n\) 是边界的内指单位法向量，证明
\[
\int_M\operatorname{div}(X)\,dV_M
=-\int_{\partial M} X\cdot n\,dV_{\partial M}.
\]
\end{problem}

\begin{solution}
散度定理通常用外指单位法向量 \(\nu_{\mathrm{out}}\) 表述：
\[
\int_M\operatorname{div}(X)\,dV_M
=\int_{\partial M}\langle X,\nu_{\mathrm{out}}\rangle\,dV_{\partial M}.
\]
题目中 \(n\) 是内指单位法向量，所以
\[
n=-\nu_{\mathrm{out}}.
\]
因此
\[
\int_M\operatorname{div}(X)\,dV_M
=-\int_{\partial M}\langle X,n\rangle\,dV_{\partial M}.
\]
若按题面右侧写成 \(+\int_{\partial M}X\cdot n\)，则那对应的是把 \(n\) 取为外法向。内法向约定下应有上式负号。
\end{solution}

\begin{problem}
设 \(F^k(M)\) 为可微流形 \(M\) 上所有 \(C^\infty\) 的 \(k\)-形式空间。设 \(U,V\) 是 \(M\) 的开子集。
\begin{enumerate}[label=(\alph*)]
\item 解释短正合列
\[
0\to F(U\cup V)\to F(U)\oplus F(V)\to F(U\cap V)\to0
\]
如何产生。
\item 写出由该短正合列诱导的 de Rham 上同调长正合列，并明确说明连接同态
\[
H^k_{\mathrm{dR}}(U\cap V)\to H^{k+1}_{\mathrm{dR}}(U\cup V)
\]
如何产生。
\end{enumerate}
\end{problem}

\begin{solution}
\begin{enumerate}[label=(\alph*)]
\item 映射
\[
F^k(U\cup V)\to F^k(U)\oplus F^k(V)
\]
定义为限制：
\[
\omega\mapsto(\omega|_U,\omega|_V).
\]
映射
\[
F^k(U)\oplus F^k(V)\to F^k(U\cap V)
\]
定义为
\[
(\alpha,\beta)\mapsto \alpha|_{U\cap V}-\beta|_{U\cap V}.
\]
核正是能在交集上相同的形式对，它们可唯一粘合成 \(U\cup V\) 上的形式。满射性可由从属于覆盖 \(\{U,V\}\) 的单位分解证明：给定 \(\eta\in F^k(U\cap V)\)，取 \(\rho_U+\rho_V=1\)，令
\[
\alpha=\rho_V\eta\text{ 在 }U\cap V\text{ 上延拓到 }U,\qquad
\beta=-\rho_U\eta\text{ 延拓到 }V,
\]
则 \(\alpha-\beta=\eta\)。

\item 短正合列是 cochain complex 的短正合列，因此诱导 Mayer--Vietoris 长正合列：
\[
\cdots\to H^k(U\cup V)\to H^k(U)\oplus H^k(V)
\to H^k(U\cap V)
\xrightarrow{\delta}
H^{k+1}(U\cup V)\to\cdots .
\]
连接同态构造如下。取闭形式 \(\eta\in F^k(U\cap V)\)。由满射性，取 \((\alpha,\beta)\in F^k(U)\oplus F^k(V)\)，使
\[
\alpha|_{U\cap V}-\beta|_{U\cap V}=\eta.
\]
因为 \(d\eta=0\)，有
\[
d\alpha|_{U\cap V}-d\beta|_{U\cap V}=0.
\]
所以 \(d\alpha\) 与 \(d\beta\) 可粘合为 \(U\cup V\) 上的 \((k+1)\)-形式 \(\theta\)。定义
\[
\delta[\eta]=[\theta]\in H^{k+1}(U\cup V).
\]
不同提升的选择只改变一个恰当形式，因此 \(\delta\) 良定义。
\end{enumerate}
\end{solution}

\begin{problem}
设 \(M\) 是 Riemann \(n\)-流形。证明 \(p\in M\) 处的 scalar curvature \(R(p)\) 可由
\[
R(p)=\frac1{\operatorname{vol}(S^{n-1})}
\int_{S^{n-1}}\operatorname{Ric}_p(x)\,dS^{n-1}
\]
给出，其中 \(x\in S^{n-1}\subset T_pM\)，\(\operatorname{Ric}_p(x)=\operatorname{Ric}_p(x,x)\)。
\end{problem}

\begin{solution}
此题中的 \(R(p)\) 是归一化 scalar curvature，即 Ricci 张量在单位方向上的平均值。取 \(T_pM\) 的正交基 \(e_i\)，写
\[
x=\sum_i x_i e_i,\qquad \sum_i x_i^2=1.
\]
则
\[
\operatorname{Ric}(x,x)=\sum_{i,j}\operatorname{Ric}_{ij}x_ix_j.
\]
球面对称性给出
\[
\int_{S^{n-1}}x_ix_j\,dS=0\quad(i\ne j),
\]
\[
\int_{S^{n-1}}x_i^2\,dS=\frac{\operatorname{vol}(S^{n-1})}{n}.
\]
因此
\[
\frac1{\operatorname{vol}(S^{n-1})}\int_{S^{n-1}}\operatorname{Ric}(x,x)\,dS
=\frac1n\sum_i\operatorname{Ric}_{ii}.
\]
右边即归一化 scalar curvature。若采用未归一化 scalar curvature \(\operatorname{Scal}=\sum_i\operatorname{Ric}_{ii}\)，则公式右边应乘以 \(n\)。
\end{solution}

\begin{problem}
证明 Schur 引理：在维数至少为 \(3\) 的 Riemann 流形上，若 Ricci 曲率只依赖于基点而不依赖于切向方向，即
\[
\operatorname{Ric}=\lambda(p)g,
\]
则 \(\lambda\) 必为常数；等价地，该流形是 Einstein，且 Einstein 常数为常数。
\end{problem}

\begin{solution}
设 \(\dim M=n\ge3\)，且
\[
\operatorname{Ric}=\lambda g.
\]
scalar curvature 为
\[
S=\operatorname{tr}_g\operatorname{Ric}=n\lambda.
\]
contracted Bianchi 恒等式为
\[
\operatorname{div}\operatorname{Ric}=\frac12\,dS.
\]
另一方面，
\[
\operatorname{div}(\lambda g)=d\lambda,
\]
因为 \(\nabla g=0\)。所以
\[
d\lambda=\frac12\,d(n\lambda)=\frac n2\,d\lambda.
\]
即
\[
\left(1-\frac n2\right)d\lambda=0.
\]
由于 \(n\ge3\)，系数非零，故
\[
d\lambda=0.
\]
因此 \(\lambda\) 在每个连通分支上为常数。命题得证。
\end{solution}

\subsection{2012 年团体赛：几何与拓扑}

原题要求：从以下 6 道题中选择 5 道作答。

\begin{problem}
证明实射影空间 \(\mathbb{RP}^n\) 是 \(n\) 维可微流形。
\end{problem}

\begin{solution}
\(\mathbb{RP}^n\) 定义为
\[
(\R^{n+1}\setminus\{0\})/\sim,\qquad x\sim\lambda x\quad(\lambda\ne0).
\]
令
\[
U_i=\{[x_0:\cdots:x_n]:x_i\ne0\}.
\]
这些开集覆盖 \(\mathbb{RP}^n\)。在 \(U_i\) 上定义坐标
\[
\varphi_i([x_0:\cdots:x_n])
=\left(\frac{x_0}{x_i},\ldots,\widehat{\frac{x_i}{x_i}},\ldots,\frac{x_n}{x_i}\right)\in\R^n.
\]
该定义与代表元选择无关，并且 \(\varphi_i\) 为 \(U_i\) 到 \(\R^n\) 的双射。

在交叠 \(U_i\cap U_j\) 上，坐标变换由分式函数给出，例如把以 \(x_i=1\) 的坐标换成以 \(x_j=1\) 的坐标，需要除以非零坐标分量。这些函数在交叠区域光滑。因此这些坐标图构成光滑 atlas，\(\mathbb{RP}^n\) 是 \(n\) 维可微流形。
\end{solution}

\begin{problem}
设 \(M,N\) 是 \(n\) 维光滑、紧、连通流形，\(f:M\to N\) 是光滑映射，且处处
\[
\operatorname{rank}df=n.
\]
证明 \(f\) 是覆盖映射。
\end{problem}

\begin{solution}
由 \(\operatorname{rank}df=n\)，反函数定理说明 \(f\) 是局部微分同胚。特别地，\(f(M)\) 是 \(N\) 中开集。又 \(M\) 紧，\(f(M)\) 紧，从而在 Hausdorff 流形 \(N\) 中闭。由于 \(N\) 连通且 \(f(M)\) 非空，故
\[
f(M)=N.
\]

对任意 \(y\in N\)，纤维 \(f^{-1}(y)\) 是紧集中的离散集，因此有限：
\[
f^{-1}(y)=\{x_1,\ldots,x_r\}.
\]
对每个 \(x_i\)，取邻域 \(U_i\)，使
\[
f|_{U_i}:U_i\to V_i
\]
为微分同胚。缩小公共邻域
\[
V=\bigcap_iV_i
\]
并利用紧性排除纤维外的其他原像，可得
\[
f^{-1}(V)=\bigsqcup_{i=1}^r U_i',
\]
且每个
\[
f|_{U_i'}:U_i'\to V
\]
为微分同胚。因此 \(V\) 被 evenly covered，\(f\) 是覆盖映射。
\end{solution}

\begin{problem}
给定任意 Riemann 流形 \((M^n,g)\)，证明 \(M^n\) 上存在唯一的 Riemannian connection。
\end{problem}

\begin{solution}
Riemannian connection 指 Levi-Civita 联络，即满足
\[
\nabla g=0
\]
且无挠
\[
\nabla_XY-\nabla_YX=[X,Y]
\]
的联络。

唯一性由 Koszul 公式给出。若这样的联络存在，则
\[
2\langle\nabla_XY,Z\rangle
=X\langle Y,Z\rangle+Y\langle Z,X\rangle-Z\langle X,Y\rangle
\]
\[
\quad-\langle X,[Y,Z]\rangle+\langle Y,[Z,X]\rangle+\langle Z,[X,Y]\rangle.
\]
右边只由度量和 Lie bracket 决定，所以 \(\nabla_XY\) 唯一。

存在性：用上式定义 \(\nabla_XY\)。可直接检查它对 \(X\) 为 \(C^\infty\)-线性，对 \(Y\) 满足 Leibniz 法则，因此是联络。再把公式作循环代换，可验证
\[
\nabla_XY-\nabla_YX=[X,Y]
\]
以及
\[
X\langle Y,Z\rangle=\langle\nabla_XY,Z\rangle+\langle Y,\nabla_XZ\rangle.
\]
故它无挠且与度量相容。存在唯一性得证。
\end{solution}

\begin{problem}
设 \(S^n\) 是 \(\R^{n+1}\) 中单位球面，\(f:S^n\to S^n\) 是连续映射。若 \(\deg f\) 为奇数，证明存在 \(x_0\in S^n\)，使
\[
f(-x_0)=-f(x_0).
\]
\end{problem}

\begin{solution}
反设不存在 \(x\) 使 \(f(-x)=-f(x)\)。令
\[
h(x)=\frac{f(x)+f(-x)}{|f(x)+f(-x)|}.
\]
若不存在 \(f(-x)=-f(x)\)，则 \(h\) 定义良好且满足
\[
h(-x)=h(x).
\]
于是 \(h\) 通过 \(\mathbb{RP}^n\) 分解。并且 \(f\) 与 \(h\) 同伦，因为
\[
(1-t)f(x)+t(f(x)+f(-x))
\]
在归一化后不经过零；若中间为零会推出 \(f(-x)\) 是 \(f(x)\) 的负倍，因二者单位长即 \(f(-x)=-f(x)\)，矛盾。因此
\[
\deg f=\deg h.
\]
但 \(h\) 通过 \(\mathbb{RP}^n\) 分解。若 \(n\) 偶数，则 \(H_n(\mathbb{RP}^n;\mathbb Z)=0\)，故 \(\deg h=0\)；若 \(n\) 奇数，则覆盖 \(S^n\to\mathbb{RP}^n\) 的度为 \(2\)，所以任何经 \(\mathbb{RP}^n\) 分解的映射 \(S^n\to S^n\) 度数为偶数。两种情形都说明 \(\deg h\) 为偶数，矛盾于 \(\deg f\) 为奇数。因此所求 \(x_0\) 存在。
\end{solution}

\begin{problem}
陈述并证明定向紧流形的 Stokes 定理。
\end{problem}

\begin{solution}
Stokes 定理：设 \(M\) 是定向紧 \(n\) 维光滑流形，边界 \(\partial M\) 取诱导边界定向。若 \(\omega\) 是 \(M\) 上光滑 \((n-1)\)-形式，则
\[
\int_M d\omega=\int_{\partial M}\omega.
\]

证明。先在半空间
\[
\mathbb H^n=\{x_n\ge0\}
\]
的坐标立方体中证明。写
\[
\omega=\sum_{i=1}^n(-1)^{i-1}a_i\,dx_1\wedge\cdots\wedge\widehat{dx_i}\wedge\cdots\wedge dx_n.
\]
则
\[
d\omega=\sum_i\frac{\partial a_i}{\partial x_i}\,dx_1\wedge\cdots\wedge dx_n.
\]
对紧支撑形式积分并用一维微积分基本定理，内部方向的边界项相互抵消，只剩 \(x_n=0\) 的边界项，正好等于 \(\int_{\partial M}\omega\)。

一般情形取有限坐标覆盖和从属于覆盖的单位分解 \(\{\rho_j\}\)。因为
\[
\omega=\sum_j\rho_j\omega,
\]
且每项支撑在坐标图中，可把局部结果相加：
\[
\int_Md\omega=\sum_j\int_Md(\rho_j\omega)
=\sum_j\int_{\partial M}\rho_j\omega
=\int_{\partial M}\omega.
\]
这里使用
\[
\sum_j d\rho_j=0
\]
保证分解与外微分相容。定理得证。
\end{solution}

\begin{problem}
设 \(M\) 是 \(\R^3\) 中曲面。设 \(D\) 是 \(M\) 中单连通区域，其边界 \(\partial D\) 紧且由有限多条光滑曲线组成。证明 Gauss--Bonnet 公式
\[
\int_{\partial D}k_g\,ds+\sum_j(\pi-\alpha_j)+\iint_DK\,dA=2\pi,
\]
其中 \(k_g\) 是边界曲线的测地曲率，\(\alpha_j\) 是顶点处内角，\(K\) 是 Gaussian 曲率，\(dA\) 是面积元。
\end{problem}

\begin{solution}
在 \(D\) 上取局部正交标架 \(e_1,e_2\)，设连接 \(1\)-形式为 \(\omega_{12}\)。结构方程给出
\[
d\omega_{12}=-K\,dA
\]
（符号随定向约定变化；以下与公式保持一致）。

沿边界的每条光滑弧，设切向量与 \(e_1\) 的夹角为 \(\theta\)。测地曲率满足 Liouville 公式
\[
k_g\,ds=d\theta+\omega_{12}.
\]
沿所有边界弧积分并相加，角函数 \(\theta\) 在顶点处跳跃。若内角为 \(\alpha_j\)，外转角贡献为
\[
\pi-\alpha_j.
\]
因此
\[
\int_{\partial D}k_g\,ds
=\int_{\partial D}\omega_{12}
+2\pi-\sum_j(\pi-\alpha_j)
\]
对单连通区域边界总转角为 \(2\pi\)。由 Stokes 定理，
\[
\int_{\partial D}\omega_{12}
=\int_D d\omega_{12}
=-\int_DK\,dA
\]
按上述符号代入，整理得
\[
\int_{\partial D}k_g\,ds+\sum_j(\pi-\alpha_j)+\int_DK\,dA=2\pi.
\]
这就是 Gauss--Bonnet 公式。
\end{solution}

\subsection{2013 年团体赛：几何与拓扑}

原题要求：从以下 6 道题中选择 5 道作答。

\begin{problem}
设
\[
X=\{(x,y,0):x^2+y^2=1\}\cup\{(x,0,z):x^2+z^2=1\}.
\]
求
\[
\pi_1(\R^3\setminus X).
\]
\end{problem}

\begin{solution}
空间 \(X\) 是两条圆在两点 \((\pm1,0,0)\) 相交形成的 theta 型图。作为图，它有
\[
V=2,\qquad E=4,
\]
故
\[
\operatorname{rank}H_1(X)=E-V+1=3.
\]

对 \(\R^3\) 中嵌入图的补空间，可用 Wirtinger 表示或 Alexander 对偶先判断 Abel 化：
\[
H_1(\R^3\setminus X)\cong H^1(X)\cong\mathbb Z^3.
\]
该具体嵌入是平面图在两个互相垂直平面中的标准嵌入，没有额外结关系；取图的正则邻域 \(N(X)\)，其边界是 genus \(3\) 的 handlebody 边界，而补空间 \(\R^3\setminus \operatorname{int}N(X)\) 也是 genus \(3\) 的 handlebody。handlebody 的基本群是自由群，其秩为 genus。

因此
\[
\pi_1(\R^3\setminus X)\cong F_3,
\]
即秩 \(3\) 的自由群。
\end{solution}

\begin{problem}
设 \(M\) 是光滑连通流形，\(f:M\to M\) 是单射光滑映射，并满足
\[
f\circ f=f.
\]
证明像集 \(f(M)\) 是 \(M\) 中的光滑子流形。
\end{problem}

\begin{solution}
由
\[
f(f(x))=f(x)
\]
可知 \(f\) 在像集 \(f(M)\) 上逐点为恒等。又 \(f\) 单射，若 \(f(x)=y\in f(M)\)，则
\[
f(y)=y=f(x),
\]
单射性推出
\[
y=x.
\]
因此每个 \(x\in M\) 都满足 \(f(x)=x\)，即 \(f=\operatorname{id}_M\)。

于是
\[
f(M)=M,
\]
当然是 \(M\) 的光滑子流形。
\end{solution}

\begin{problem}
设
\[
T^2=\{(z,w)\in\C^2:|z|=1,\ |w|=1\}.
\]
定义
\[
f:T^2\to T^2,\qquad f(z,w)=(zw^3,w).
\]
证明 \(f\) 是微分同胚。
\end{problem}

\begin{solution}
映射显然光滑。构造逆映射
\[
g:T^2\to T^2,\qquad g(Z,W)=(ZW^{-3},W).
\]
因为 \(|W|=1\)，所以 \(W^{-3}\) 光滑且仍在单位圆上。直接计算：
\[
g(f(z,w))=g(zw^3,w)=(zw^3w^{-3},w)=(z,w),
\]
\[
f(g(Z,W))=f(ZW^{-3},W)=(ZW^{-3}W^3,W)=(Z,W).
\]
因此 \(g=f^{-1}\)，且 \(f,g\) 都光滑，所以 \(f\) 是微分同胚。
\end{solution}

\begin{problem}
证明：任意三维 Einstein 流形具有常截面曲率。
\end{problem}

\begin{solution}
设三维 Riemann 流形满足
\[
\operatorname{Ric}=\lambda g.
\]
由 Schur 引理，维数至少 \(3\) 时 \(\lambda\) 为常数。三维中 Weyl 张量恒为零，Riemann 曲率张量完全由 Ricci 张量和 scalar curvature 决定：
\[
R_{ijkl}=g_{ik}R_{jl}-g_{il}R_{jk}-g_{jk}R_{il}+g_{jl}R_{ik}
-\frac S2(g_{ik}g_{jl}-g_{il}g_{jk}).
\]
代入
\[
R_{ij}=\lambda g_{ij},\qquad S=3\lambda,
\]
得
\[
R_{ijkl}
=\frac{\lambda}{2}(g_{ik}g_{jl}-g_{il}g_{jk}).
\]
这正是常截面曲率张量形式。因此任意二维截面的截面曲率都等于
\[
K=\frac{\lambda}{2},
\]
为常数。
\end{solution}

\begin{problem}
陈述并证明完备 Riemann 流形的 Myers 定理。
\end{problem}

\begin{solution}
Myers 定理：若 \((M^n,g)\) 完备且
\[
\operatorname{Ric}\ge (n-1)k>0,
\]
则
\[
\operatorname{diam}(M)\le \frac{\pi}{\sqrt{k}}.
\]
特别地，\(M\) 紧且 \(\pi_1(M)\) 有限。

证明。设 \(\gamma:[0,l]\to M\) 是单位速度最短测地线。取沿 \(\gamma\) 的平行正交标架
\[
E_1,\ldots,E_{n-1}\perp\gamma'.
\]
对变分场
\[
V_i(t)=\sin\frac{\pi t}{l}\,E_i(t)
\]
使用 index form。由于 \(\gamma\) 最短，
\[
\sum_{i=1}^{n-1}I(V_i,V_i)\ge0.
\]
而
\[
\sum_i I(V_i,V_i)
=\int_0^l\left((n-1)\left(\frac{\pi}{l}\right)^2\cos^2\frac{\pi t}{l}
-\operatorname{Ric}(\gamma',\gamma')\sin^2\frac{\pi t}{l}\right)dt.
\]
由 Ricci 下界，
\[
0\le (n-1)\int_0^l\left[
\left(\frac{\pi}{l}\right)^2\cos^2\frac{\pi t}{l}
-k\sin^2\frac{\pi t}{l}
\right]dt.
\]
两个三角函数平方积分都为 \(l/2\)，故
\[
0\le \frac{(n-1)l}{2}\left(\frac{\pi^2}{l^2}-k\right).
\]
因此
\[
l\le\frac{\pi}{\sqrt{k}}.
\]
任意两点可由最短测地线连接，所以直径满足同样上界。完备且有界闭球紧，由 Hopf--Rinow 定理得 \(M\) 紧。其 universal cover 仍满足同样 Ricci 下界，也紧；因此覆盖层数有限，\(\pi_1(M)\) 有限。
\end{solution}

\begin{problem}
设 \(C\) 是 \(\R^3\) 中正则闭曲线，挠率为 \(\tau\)。称
\[
\frac1{2\pi}\int_C\tau\,ds
\]
为 \(C\) 的总挠率。证明：给定 \(\R^3\) 中光滑曲面 \(M\)，若 \(M\) 上任意正则闭曲线 \(C\) 的总挠率总为整数，则 \(M\) 是球面或平面的一部分。
\end{problem}

\begin{solution}
设曲线 \(C\subset M\) 的 Darboux 标架为
\[
T,\quad u=N_M\times T,\quad N_M,
\]
其中 \(N_M\) 为曲面单位法向。Frenet 主法向与曲面法向之间的夹角记为 \(\theta\)。曲线的 Frenet 挠率满足
\[
\tau=\tau_g+\frac{d\theta}{ds},
\]
其中 \(\tau_g\) 是曲线在曲面上的测地挠率。闭曲线一周后 \(\theta\) 的变化为 \(2\pi\) 的整数倍，因此
\[
\frac1{2\pi}\int_C\tau\,ds
\equiv
\frac1{2\pi}\int_C\tau_g\,ds
\pmod{\mathbb Z}.
\]
题设对任意闭曲线给出左边为整数，所以
\[
\int_C\tau_g\,ds\in2\pi\mathbb Z.
\]
对足够小的闭曲线，该积分可任意小，故必为 \(0\)。于是任意足够小闭曲线都有
\[
\int_C\tau_g\,ds=0.
\]

在局部正交主方向标架中，测地挠率为
\[
\tau_g=(k_2-k_1)\sin\phi\cos\phi,
\]
其中 \(k_1,k_2\) 是主曲率，\(\phi\) 是曲线切向与第一主方向夹角。若某点不是 umbilic，即 \(k_1\ne k_2\)，可取围绕该点的小闭曲线，使 \(\sin\phi\cos\phi\) 的加权积分非零，矛盾。因此每一点都为 umbilic：
\[
k_1=k_2.
\]
曲面全脐。

Euclidean 三维空间中的连通全脐曲面只能是平面或球面的一部分。证明如下：若形算子 \(S=\lambda I\)，Codazzi 方程给出 \(d\lambda=0\)，故 \(\lambda\) 常数。若 \(\lambda=0\)，法向常值，曲面为平面的一部分；若 \(\lambda\ne0\)，则
\[
d\left(x+\frac1\lambda N\right)=0,
\]
所以 \(x+\lambda^{-1}N\) 为常向量，曲面为球面的一部分。
\end{solution}

\subsection{2014 年团体赛：几何与拓扑}

原题要求：从以下 6 道题中选择 5 道作答。

\begin{problem}
计算圆 \(S^1\) 与环面 \(T=S^1\times S^1\) 的楔和
\[
S^1\vee T
\]
的基本群和同调群。
\end{problem}

\begin{solution}
由 Seifert--van Kampen 定理，
\[
\pi_1(S^1\vee T)
\cong \pi_1(S^1)*\pi_1(T)
\cong \mathbb Z*(\mathbb Z^2).
\]
可写为
\[
\langle a,b,c\mid [b,c]=1\rangle.
\]

楔和的约化同调为直和：
\[
\widetilde H_i(S^1\vee T)
\cong \widetilde H_i(S^1)\oplus\widetilde H_i(T).
\]
因此
\[
H_0\cong\mathbb Z,
\]
\[
H_1\cong \mathbb Z\oplus\mathbb Z^2\cong\mathbb Z^3,
\]
\[
H_2\cong H_2(T)\cong\mathbb Z,
\]
而
\[
H_i=0,\qquad i\ge3.
\]
\end{solution}

\begin{problem}
设 \(F:G\times M\to M\) 是群 \(G\) 在光滑流形 \(M\) 上的 proper discontinuous 作用。证明 \(M/G\) 可定向当且仅当 \(M\) 可定向且每个 \(F(g,\cdot)\) 保持 \(M\) 的定向。用此结论证明 Möbius 带不可定向，并证明 \(\mathbb{RP}^n\) 可定向当且仅当 \(n\) 为奇数。
\end{problem}

\begin{solution}
若 \(M/G\) 可定向，则覆盖映射
\[
\pi:M\to M/G
\]
把 \(M/G\) 的定向拉回为 \(M\) 的定向。deck transformation \(g:M\to M\) 满足
\[
\pi\circ g=\pi,
\]
所以保持拉回定向。

反过来，若 \(M\) 可定向且 \(G\) 中每个元素保定向，则可在商空间上定义定向：取 \(M/G\) 中一点的均匀覆盖邻域，用任一 lift 的定向定义商邻域定向。若换另一个 lift，二者相差某个 \(g\in G\)，而 \(g\) 保定向，所以定义无关。故 \(M/G\) 可定向。

Möbius 带可看作条带
\[
\R\times(-1,1)
\]
关于变换
\[
(x,y)\mapsto(x+1,-y)
\]
的商。该变换反定向，因此商不可定向。

\(\mathbb{RP}^n=S^n/\{x\sim -x\}\)。对径映射 \(A:S^n\to S^n\) 的 degree 为
\[
\deg A=(-1)^{n+1}.
\]
它保定向当且仅当
\[
(-1)^{n+1}=1,
\]
即 \(n\) 为奇数。因此 \(\mathbb{RP}^n\) 可定向当且仅当 \(n\) 为奇数。
\end{solution}

\begin{problem}
\begin{enumerate}[label=(\alph*)]
\item 设 \(Y\) 由 \(S^2\times[0,1]\) 经如下识别得到：
\[
(x,0)\sim(-x,0),\qquad (x,1)\sim(-x,1).
\]
证明 \(Y\) 同胚于连通和 \(\mathbb{RP}^3\#\mathbb{RP}^3\)。
\item 证明 \(S^2\times S^1\) 是 \(\mathbb{RP}^3\#\mathbb{RP}^3\) 的二重覆盖。
\end{enumerate}
\end{problem}

\begin{solution}
\begin{enumerate}[label=(\alph*)]
\item \(\mathbb{RP}^3\) 去掉一个开 \(3\)-球后，同胚于
\[
(S^2\times[0,1/2])/(x,0)\sim(-x,0).
\]
这是因为 \(\mathbb{RP}^3\) 可由 \(3\)-球边界对径识别得到，去掉内部小球后剩下一个带有一端对径识别的圆柱。两个这样的空间沿未识别的 \(S^2\) 边界粘合，正得到
\[
(S^2\times[0,1])/\{(x,0)\sim(-x,0),(x,1)\sim(-x,1)\}.
\]
这就是 \(\mathbb{RP}^3\#\mathbb{RP}^3\)。

\item 定义
\[
S^2\times[0,1]\to Y
\]
为商映射。把 \(S^2\times[0,1]\) 的两端按
\[
(x,0)\sim(x,1)
\]
识别得到 \(S^2\times S^1\)。映射
\[
S^2\times S^1\to Y
\]
由把端点处的对径识别作为 deck transformation 生成。每个点在 \(Y\) 中有两个 lift，deck transformation 可写为
\[
(x,t)\mapsto(-x,-t)
\]
在圆周参数下。故它是二重覆盖。
\end{enumerate}
\end{solution}

\begin{problem}
证明 Lie 群 \(G\) 上的双不变度量具有非负截面曲率。
\end{problem}

\begin{solution}
对双不变度量，左不变向量场 \(X,Y\) 满足
\[
\nabla_XY=\frac12[X,Y].
\]
于是曲率张量为
\[
R(X,Y)Y=-\frac14[[X,Y],Y]
\]
按常用约定。利用 \(\operatorname{ad}\)-不变性：
\[
\langle [A,B],C\rangle=\langle A,[B,C]\rangle,
\]
得
\[
\langle R(X,Y)Y,X\rangle
=\frac14\|[X,Y]\|^2.
\]
因此截面曲率
\[
K(X,Y)=
\frac{\langle R(X,Y)Y,X\rangle}
{\|X\|^2\|Y\|^2-\langle X,Y\rangle^2}
=\frac14
\frac{\|[X,Y]\|^2}{\|X\wedge Y\|^2}
\ge0.
\]
\end{solution}

\begin{problem}
设 \(M\) 是上半平面 \(\R^2_+\)，度量为
\[
ds^2=\frac{dx^2+dy^2}{y^k}.
\]
问 \(k\) 取何值时 \(M\) 完备？
\end{problem}

\begin{solution}
该度量共形于 Euclidean 度量，长度元为
\[
ds=y^{-k/2}\sqrt{dx^2+dy^2}.
\]
需要检查到边界 \(y=0\) 和到无穷远的距离。

沿竖直曲线 \(x=\text{常数}\)，从 \(y=1\) 到 \(y=0\) 的长度为
\[
\int_0^1 y^{-k/2}\,dy.
\]
该积分发散当且仅当
\[
\frac{k}{2}\ge1,\qquad k\ge2.
\]
若 \(k<2\)，边界 \(y=0\) 距离有限，不完备。

再看到 \(y=\infty\) 的长度：
\[
\int_1^\infty y^{-k/2}\,dy.
\]
该积分发散当且仅当
\[
\frac{k}{2}\le1,\qquad k\le2.
\]
若 \(k>2\)，向上无穷远距离有限，不完备。

同时满足两端距离无限只可能
\[
k=2.
\]
此时度量为双曲上半平面度量，完备。因此 \(M\) 完备当且仅当
\[
\boxed{k=2}.
\]
\end{solution}

\begin{problem}
给定任意不可定向流形 \(M\)，证明存在光滑可定向流形 \(\widetilde M\)，它是 \(M\) 的二重覆盖。并求当 \(M=\mathbb{RP}^2\) 或 \(M\) 为 Möbius 带时的 \(\widetilde M\)。
\end{problem}

\begin{solution}
构造定向双覆盖。令
\[
\widetilde M=\{(p,o_p):p\in M,\ o_p\text{ 是 }T_pM\text{ 的一个定向}\}.
\]
每个点 \(p\) 有两个可能定向，所以自然投影
\[
\widetilde M\to M
\]
是二重覆盖。局部坐标的 Jacobian 符号决定两个 sheet 如何粘合，该空间自然可定向：在 \((p,o_p)\) 处选择与 \(o_p\) 一致的局部坐标即可。

若 \(M=\mathbb{RP}^2\)，其定向双覆盖为
\[
S^2\to\mathbb{RP}^2.
\]
若 \(M\) 是 Möbius 带，其定向双覆盖为圆柱
\[
S^1\times[0,1]\to M.
\]
\end{solution}

\subsection{2015 年团体赛：几何与拓扑}

原题要求：从以下 6 道题中选择 5 道作答。

\begin{problem}
设 \(SO(3)\) 为所有满足 \(\det A=1\) 且 \(A^TA=I\) 的 \(3\times3\) 实矩阵集合。证明 \(SO(3)\) 是光滑流形，并求其基本群。
\end{problem}

\begin{solution}
考虑映射
\[
F:GL(3,\R)\to \operatorname{Sym}(3),\qquad F(A)=A^TA.
\]
有
\[
O(3)=F^{-1}(I).
\]
其微分为
\[
dF_A(B)=B^TA+A^TB.
\]
若 \(A\in O(3)\)，给定任意对称矩阵 \(S\)，取
\[
B=\frac12AS,
\]
则
\[
dF_A(B)=S.
\]
所以 \(I\) 是正则值，\(O(3)\) 是 \(GL(3,\R)\) 中维数 \(9-6=3\) 的光滑子流形。\(SO(3)\) 是 \(O(3)\) 中由方程 \(\det A=1\) 选出的连通分支，故也是 \(3\) 维光滑流形。

单位四元数群
\[
S^3=\{q\in\mathbb H:|q|=1\}
\]
通过共轭作用在纯虚四元数空间 \(\operatorname{Im}\mathbb H\cong\R^3\) 上：
\[
v\mapsto qvq^{-1}.
\]
这给出满同态
\[
S^3\to SO(3),
\]
核为 \(\{\pm1\}\)。因此
\[
SO(3)\cong S^3/\{\pm1\}=\mathbb{RP}^3.
\]
于是
\[
\pi_1(SO(3))\cong\pi_1(\mathbb{RP}^3)\cong\mathbb Z/2.
\]
\end{solution}

\begin{problem}
设 \(X\) 是拓扑空间。\(X\) 的 suspension \(S(X)\) 是把 \(X\times[0,1]\) 中 \(X\times\{0\}\) 压缩为一点、\(X\times\{1\}\) 压缩为另一点得到的空间。描述 \(X\) 与 \(S(X)\) 的同调群关系。
\end{problem}

\begin{solution}
若 \(X\) 非空，则 \(S(X)\) 连通，并且约化同调满足 suspension 同构：
\[
\widetilde H_k(SX)\cong \widetilde H_{k-1}(X),\qquad k\ge1.
\]
特别地，
\[
\widetilde H_0(SX)=0.
\]

证明如下。把 suspension 写成两个 cone 的并：
\[
SX=C_+X\cup C_-X.
\]
两个 cone 都可缩，交集同伦等价于 \(X\)。Mayer--Vietoris 长正合列中，\(\widetilde H_*(C_+X)=\widetilde H_*(C_-X)=0\)，故连接同态给出
\[
\widetilde H_k(SX)\cong \widetilde H_{k-1}(X).
\]
\end{solution}

\begin{problem}
设 \(F:M\to N\) 是光滑映射。设 \(X_1,X_2\) 是 \(M\) 上光滑向量场，\(Y_1,Y_2\) 是 \(N\) 上光滑向量场。证明若
\[
Y_1=F_*X_1,\qquad Y_2=F_*X_2,
\]
则
\[
F_*[X_1,X_2]=[Y_1,Y_2].
\]
\end{problem}

\begin{solution}
这里 \(Y_i=F_*X_i\) 的意思是 \(X_i\) 与 \(Y_i\) 为 \(F\)-相关：
\[
dF(X_i)=Y_i\circ F.
\]
对任意 \(N\) 上光滑函数 \(h\)，有
\[
X_i(h\circ F)=(Y_i h)\circ F.
\]
于是
\[
[X_1,X_2](h\circ F)
=X_1X_2(h\circ F)-X_2X_1(h\circ F)
\]
\[
=(Y_1Y_2h)\circ F-(Y_2Y_1h)\circ F
=([Y_1,Y_2]h)\circ F.
\]
这正说明
\[
dF([X_1,X_2])=[Y_1,Y_2]\circ F.
\]
因此
\[
F_*[X_1,X_2]=[Y_1,Y_2].
\]
\end{solution}

\begin{problem}
设 \(M_1,M_2\) 是 \(\R^3\) 中两个紧凸闭曲面，\(f:M_1\to M_2\) 是微分同胚。假设对应点处 \(M_1,M_2\) 有相同的内法向量和 Gaussian 曲率。证明 \(f\) 是一个平移。
\end{problem}

\begin{solution}
用 Gauss 映射参数化凸曲面。对严格凸闭曲面，内法向 Gauss 映射
\[
N_i:M_i\to S^2
\]
是微分同胚；一般凸情形可由严格凸曲面逼近后取极限。题设说对应点有相同内法向，故可把两曲面都写成同一个法向 \(u\in S^2\) 的支撑函数：
\[
X_i(u)=\nabla_{S^2}h_i(u)+h_i(u)u.
\]
凸曲面的曲率满足 Minkowski 公式
\[
\frac1{K_i(u)}
=\det(\nabla^2h_i+h_i g_{S^2}).
\]
题设 \(K_1=K_2\)，因此两曲面具有同一面积测度
\[
\frac1K\,d\omega.
\]
Minkowski 唯一性定理说明：具有同一面积测度的闭凸曲面只差一个平移。因此存在常向量 \(a\in\R^3\)，使
\[
M_2=M_1+a.
\]
在支撑函数上，这对应
\[
h_2(u)=h_1(u)+a\cdot u.
\]
由相同法向对应，\(f\) 正是
\[
f(x)=x+a.
\]
所以 \(f\) 是平移。
\end{solution}

\begin{problem}
证明第二 Bianchi 恒等式：
\[
R_{ijkl;h}+R_{ijlh;k}+R_{ijhk;l}=0.
\]
\end{problem}

\begin{solution}
设 \(\omega_{ij}\) 是 Levi-Civita 联络形式，曲率形式为
\[
\Omega_{ij}=d\omega_{ij}+\sum_k\omega_{ik}\wedge\omega_{kj}.
\]
结构方程的 Bianchi 恒等式为
\[
d\Omega_{ij}
\sum_k\omega_{ik}\wedge\Omega_{kj}
\sum_k\omega_{jk}\wedge\Omega_{ik}=0,
\]
这里利用了 \(\omega_{ij}=-\omega_{ji}\) 与 \(\Omega_{ij}=-\Omega_{ji}\) 来写成对称形式。把曲率形式写作
\[
\Omega_{ij}=\frac12\sum_{k,l}R_{ijkl}\theta_k\wedge\theta_l.
\]
对上式展开，并比较
\[
\theta_h\wedge\theta_k\wedge\theta_l
\]
的系数，即得
\[
R_{ijkl;h}+R_{ijlh;k}+R_{ijhk;l}=0.
\]
\end{solution}

\begin{problem}
设 \(M_1,M_2\) 是两个完备 \(n\) 维 Riemann 流形，\(\gamma_i:[0,a]\to M_i\) 为弧长参数化测地线。令 \(\rho_i\) 为到 \(\gamma_i(0)\) 的距离函数。假设 \(\gamma_i(a)\) 不在 \(\gamma_i(0)\) 的 cut locus 中，并且对任意 \(0\le t\le a\)，有截面曲率不等式
\[
K_1\left(X_1,\gamma_1'\right)
\ge
K_2\left(X_2,\gamma_2'\right),
\]
其中 \(X_i\in T_{\gamma_i(t)}M_i\) 为任意垂直于 \(\gamma_i'\) 的单位向量。证明
\[
\operatorname{Hess}(\rho_1)(\widetilde X_1,\widetilde X_1)
\le
\operatorname{Hess}(\rho_2)(\widetilde X_2,\widetilde X_2),
\]
其中 \(\widetilde X_i\in T_{\gamma_i(a)}M_i\) 为垂直于 \(\gamma_i'(a)\) 的单位向量。
\end{problem}

\begin{solution}
这是 Hessian 比较定理。沿 \(\gamma_i\) 取 Jacobi 场 \(J_i\)，满足
\[
J_i(0)=0,\qquad J_i(a)=\widetilde X_i.
\]
由于 \(\gamma_i(a)\) 不在 cut locus 中，距离函数在该点附近光滑，并且
\[
\operatorname{Hess}\rho_i(\widetilde X_i,\widetilde X_i)
=I_i(J_i,J_i),
\]
其中 index form 为
\[
I_i(V,V)=\int_0^a\left(|D_tV|^2-\langle R_i(V,\gamma_i')\gamma_i',V\rangle\right)dt.
\]
Jacobi 场在给定端点条件下使 index form 取最小值。

用平行标架把 \(M_1\) 中沿 \(\gamma_1\) 的向量场与 \(M_2\) 中沿 \(\gamma_2\) 的向量场对应，使长度函数与端点单位向量对应。曲率假设 \(K_1\ge K_2\) 意味着对任意对应比较场 \(V_1,V_2\)，有
\[
I_1(V_1,V_1)\le I_2(V_2,V_2).
\]
取 \(M_2\) 中的 Jacobi 场 \(J_2\)，把它对应到 \(M_1\) 中的场 \(\overline J_1\)。则
\[
I_1(J_1,J_1)\le I_1(\overline J_1,\overline J_1)\le I_2(J_2,J_2).
\]
于是
\[
\operatorname{Hess}\rho_1(\widetilde X_1,\widetilde X_1)
\le
\operatorname{Hess}\rho_2(\widetilde X_2,\widetilde X_2).
\]
\end{solution}

\subsection{2016 年团体赛：几何与拓扑}

原题要求：从以下 6 道题中选择 5 道作答。

\begin{problem}
证明 \(\mathbb{CP}^{2n}\) 不覆盖除自身以外的任何流形。
\end{problem}

\begin{solution}
设
\[
\pi:\mathbb{CP}^{2n}\to M
\]
是连通流形 \(M\) 的覆盖，覆盖度为 \(d\)。因为 \(\mathbb{CP}^{2n}\) 紧，\(d\) 有限。又 \(\mathbb{CP}^{2n}\) 单连通，所以这是 universal cover，deck transformation 群 \(G\) 的阶为 \(d\)，并自由作用在 \(\mathbb{CP}^{2n}\) 上。

若 \(d>1\)，取非恒等 deck transformation \(g\)。它无不动点。计算 Lefschetz 数。上同调环为
\[
H^*(\mathbb{CP}^{2n};\mathbb Q)=\mathbb Q[h]/(h^{2n+1}).
\]
任意自同胚诱导 \(g^*h=\pm h\)。但 \(g\) 作为 deck transformation 保持定向；即使只看 Lefschetz 数，若 \(g^*h=h\)，则
\[
L(g)=2n+1\ne0.
\]
若 \(g^*h=-h\)，则
\[
L(g)=1-1+\cdots+(-1)^{2n}=1\ne0.
\]
因此 Lefschetz 不动点定理给出 \(g\) 有不动点，矛盾。故 \(d=1\)，覆盖为平凡覆盖，\(M=\mathbb{CP}^{2n}\)。
\end{solution}

\begin{problem}
设 \(X\) 是拓扑空间，\(p\in X\)。约化 suspension \(\Sigma X\) 是把
\[
(X\times\{0,1\})\cup(\{p\}\times[0,1])
\]
压缩为一点所得空间。描述 \(X\) 与 \(\Sigma X\) 的同调群关系。
\end{problem}

\begin{solution}
约化 suspension 满足
\[
\widetilde H_k(\Sigma X)\cong \widetilde H_{k-1}(X),\qquad k\ge1,
\]
并且若 \(X\) 非空，则 \(\widetilde H_0(\Sigma X)=0\)。

证明与普通 suspension 相同。把 \(\Sigma X\) 写成两个约化 cone 的并，它们都可缩，交集同伦等价于 \(X\)。Mayer--Vietoris 长正合列给出上述自然同构。
\end{solution}

\begin{problem}
陈述并证明可微流形上的 Frobenius 定理。
\end{problem}

\begin{solution}
Frobenius 定理：设 \(D\subset TM\) 是秩为 \(k\) 的光滑分布。则 \(D\) 在每点附近可积，即存在局部坐标
\[
(x^1,\ldots,x^n)
\]
使
\[
D=\operatorname{span}\left\{\frac{\partial}{\partial x^1},\ldots,\frac{\partial}{\partial x^k}\right\},
\]
当且仅当 \(D\) 对 Lie bracket 封闭：
\[
X,Y\in\Gamma(D)\quad\Longrightarrow\quad [X,Y]\in\Gamma(D).
\]

必要性显然：若 \(D\) 由坐标向量场张成，则坐标向量场两两括号为零，一般截面的括号仍在同一张成空间中。

充分性证明。取 \(D\) 的局部基 \(X_1,\ldots,X_k\)。由 involutive 条件，
\[
[X_i,X_j]=\sum_l c_{ij}^lX_l.
\]
通过解一阶线性方程，可重新选取局部基，使其两两对易。对易向量场的局部流彼此交换。取过点 \(p\) 的横截子流形 \(S\)，定义
\[
F(t^1,\ldots,t^k,s)
=\Phi_1^{t^1}\circ\cdots\circ\Phi_k^{t^k}(s).
\]
反函数定理给出局部坐标。在这些坐标中，\(D\) 正由前 \(k\) 个坐标向量场张成，积分叶为固定后 \(n-k\) 个坐标的子流形。故 \(D\) 可积。
\end{solution}

\begin{problem}
证明球面 \(S^n\) 上所有测地线恰为大圆。
\end{problem}

\begin{solution}
把 \(S^n\) 看作 \(\R^{n+1}\) 中单位球面。设 \(\gamma\) 是 \(S^n\) 上单位速度测地线。作为 \(\R^{n+1}\) 中曲线，有
\[
\langle\gamma,\gamma\rangle=1.
\]
微分得
\[
\langle\gamma',\gamma\rangle=0,\qquad
\langle\gamma'',\gamma\rangle=-|\gamma'|^2=-1.
\]
测地线条件说明 \(\gamma''\) 的切向分量为零，所以 \(\gamma''\) 纯法向；球面单位法向为 \(\gamma\)，故
\[
\gamma''=-\gamma.
\]
解此常微分方程：
\[
\gamma(t)=\cos t\,\gamma(0)+\sin t\,\gamma'(0).
\]
因此 \(\gamma\) 位于经过原点的二维平面
\[
\operatorname{span}\{\gamma(0),\gamma'(0)\}
\]
与 \(S^n\) 的交中，即大圆。

反过来，任意大圆可写为上述形式，满足 \(\gamma''=-\gamma\)，其加速度纯法向，故为测地线。
\end{solution}

\begin{problem}
设 \(M\) 是 Riemann 流形，曲率张量为 \(R\)，截面曲率为 \(K_M\)。若在点 \(x\in M\) 处
\[
a\le K_M\le b,
\]
证明对任意正交四标架 \(\{e_1,e_2,e_3,e_4\}\subset T_xM\)，有 Berger 估计
\[
|R(e_1,e_2,e_3,e_4)|\le \frac23(b-a).
\]
\end{problem}

\begin{solution}
令
\[
R_{ijkl}=R(e_i,e_j,e_k,e_l).
\]
利用代数曲率张量的极化恒等式，可把混合分量表示为截面曲率的组合：
\[
6R_{1234}
=K(e_1+e_3,e_2+e_4)-K(e_1+e_3,e_2-e_4)
\]
\[
-K(e_1-e_3,e_2+e_4)+K(e_1-e_3,e_2-e_4),
\]
其中每个 \(K(u,v)\) 表示未归一化量
\[
R(u,v,v,u)
\]
除以 \(|u\wedge v|^2\)。对正交四标架，上述四个二平面的面积因子相同。把每个截面曲率写成
\[
K=\frac{a+b}{2}+\delta,\qquad |\delta|\le\frac{b-a}{2},
\]
常数部分在极化组合中抵消，只剩 \(\delta\) 项。估计得
\[
|R_{1234}|\le\frac23(b-a).
\]
这就是 Berger 不等式。
\end{solution}

\begin{problem}
设 \(M\) 是 \(S^{n+1}\) 中闭极小超曲面，且 scalar curvature 为常数。记 \(S\) 为 \(M\) 的第二基本形式长度平方。证明
\[
S=0\quad\text{或}\quad S\ge n.
\]
\end{problem}

\begin{solution}
对 \(S^{n+1}\) 中极小超曲面，Gauss 方程给出 scalar curvature
\[
\operatorname{Scal}_M=n(n-1)-S.
\]
题设 scalar curvature 为常数，因此 \(S\) 为常数。

Simons 恒等式给出
\[
\frac12\Delta S=|\nabla A|^2+S(n-S),
\]
其中 \(A\) 是第二基本形式。由于 \(S\) 常数，左边为零：
\[
0=|\nabla A|^2+S(n-S).
\]
若 \(S=0\)，即全测地情形。若 \(S>0\)，则
\[
S(n-S)\le0,
\]
所以
\[
S\ge n.
\]
命题得证。
\end{solution}

\subsection{2017 年团体赛：几何与拓扑}

原题要求：从以下 6 道题中选择 5 道作答。

\begin{problem}
设
\[
X=(S^2\vee S^2)\cup_{f}D^3,
\]
其中贴映射 \(f:S^2\to S^2\vee S^2\) 把连接南北极的大圆压缩。计算 \(X\) 的同调群。
\end{problem}

\begin{solution}
把 \(S^2\) 沿一条大圆压缩后，商空间同胚于
\[
S^2\vee S^2.
\]
该贴映射在
\[
H_2(S^2)\to H_2(S^2\vee S^2)\cong\mathbb Z^2
\]
上把生成元送到
\[
(1,1)
\]
（取定两个半球给出的两个球面方向）。

CW 链复形为
\[
0\to C_3\cong\mathbb Z
\xrightarrow{\partial_3}
C_2\cong\mathbb Z^2
\to C_1=0\to C_0\cong\mathbb Z\to0,
\]
其中
\[
\partial_3(1)=(1,1).
\]
因此
\[
H_3(X)=\ker\partial_3=0,
\]
\[
H_2(X)=\mathbb Z^2/\langle(1,1)\rangle\cong\mathbb Z,
\]
\[
H_1(X)=0,\qquad H_0(X)=\mathbb Z.
\]
更高维同调为零。
\end{solution}

\begin{problem}
\begin{enumerate}[label=(\alph*)]
\item 设 \(A\) 是 \(\R^3\) 中一个圆，计算
\[
\pi_1(\R^3\setminus A).
\]
\item 设 \(A,B\) 是 \(\R^3\) 中两个不交圆，分别位于上半空间和下半空间，计算
\[
\pi_1(\R^3\setminus(A\cup B)).
\]
\end{enumerate}
\end{problem}

\begin{solution}
\begin{enumerate}[label=(\alph*)]
\item 单个圆的补空间同伦等价于实心环面的补，也同伦等价于 \(S^1\)。因此
\[
\pi_1(\R^3\setminus A)\cong\mathbb Z.
\]

\item 两个圆分别位于上、下半空间，互不链环。补空间同伦等价于两个圆的非链环 unlink 的补。其基本群为两个 meridian 生成的自由群：
\[
\pi_1(\R^3\setminus(A\cup B))\cong F_2.
\]
可由 van Kampen 定理证明：用一个平面把空间分成包含 \(A\) 和包含 \(B\) 的两部分，交集单连通，两个部分各贡献一个 \(\mathbb Z\)，自由积为 \(F_2\)。
\end{enumerate}
\end{solution}

\begin{problem}
考虑 \(\R^3\) 中微分 \(1\)-形式
\[
\omega=x\,dy-y\,dx+dz.
\]
证明对任何处处非零函数 \(f:\R^3\to\R\)，\(f\omega\) 都不是闭形式。
\end{problem}

\begin{solution}
若 \(f\omega\) 闭，则
\[
d(f\omega)=df\wedge\omega+f\,d\omega=0.
\]
两边与 \(\omega\) 再外乘，得
\[
f\,\omega\wedge d\omega=0,
\]
因为 \(\omega\wedge df\wedge\omega=0\)。但
\[
d\omega=d(xdy-ydx+dz)=2\,dx\wedge dy.
\]
所以
\[
\omega\wedge d\omega
=(xdy-ydx+dz)\wedge2dx\wedge dy
=2\,dz\wedge dx\wedge dy,
\]
这是处处非零的体积形式。由于 \(f\) 处处非零，
\[
f\,\omega\wedge d\omega\ne0,
\]
矛盾。因此 \(f\omega\) 不可能闭。
\end{solution}

\begin{problem}
证明
\[
Q^n=\left\{(x^1,\ldots,x^{n+1})\in\R^{n+1}:\sum_{i=1}^{n+1}(x^i)^4=1\right\}
\]
是可微流形。
\end{problem}

\begin{solution}
定义光滑函数
\[
F:\R^{n+1}\to\R,\qquad F(x)=\sum_{i=1}^{n+1}(x^i)^4.
\]
则
\[
Q^n=F^{-1}(1).
\]
梯度为
\[
\nabla F=(4(x^1)^3,\ldots,4(x^{n+1})^3).
\]
若 \(x\in Q^n\)，则至少一个 \(x^i\ne0\)，所以 \(\nabla F(x)\ne0\)。因此 \(1\) 是正则值。由正则值定理，
\[
Q^n
\]
是 \(\R^{n+1}\) 中的光滑 \(n\) 维子流形。
\end{solution}

\begin{problem}
设 \(M\) 是 \(\R^3\) 中闭曲面。证明
\[
\int_M |K|\,d\sigma\ge4\pi(1+g),
\]
其中 \(K\) 是 Gaussian 曲率，\(g\) 是 genus，\(d\sigma\) 是面积元。
\end{problem}

\begin{solution}
由 Gauss--Bonnet 定理，
\[
\int_M K\,d\sigma=2\pi\chi(M)=4\pi(1-g).
\]
另一方面，Gauss 映射
\[
N:M\to S^2
\]
的 Jacobian 为 \(K\)，所以
\[
\int_M |K|\,d\sigma
\]
是 Gauss 映射按重数计的面积。对嵌入闭曲面，任意方向的支撑平面至少在一个最大点和一个最小点与曲面相切；对 genus \(g\) 的曲面，结合 Morse 理论，典型高度函数至少有
\[
2+2g
\]
个临界点。这意味着 Gauss 映射覆盖 \(S^2\) 的总重数平均至少为 \(1+g\)。因此
\[
\int_M |K|\,d\sigma\ge (1+g)\operatorname{Area}(S^2)=4\pi(1+g).
\]
这就是 Chern--Lashof 全绝对曲率不等式在曲面情形的表述。
\end{solution}

\begin{problem}
设 \(M\) 是 \(n\) 维紧单连通 Riemann 流形。若截面曲率满足
\[
\frac14<K_M\le1,
\]
证明 \(M\) 同胚于 \(S^n\)。
\end{problem}

\begin{solution}
这是严格 \(1/4\)-pinched sphere theorem。由 Brendle--Schoen 微分球面定理，若紧单连通 Riemann 流形的截面曲率满足
\[
\frac14<K\le1,
\]
则该流形微分同胚于球面 \(S^n\)。题目只要求同胚，因此结论直接成立：
\[
M\cong S^n.
\]

在经典表述中，严格 pinching 条件排除了复射影空间、四元数射影空间和 Cayley 平面等等距 pinched 到 \(1/4\) 的例外情形；单连通条件排除了球面空间形式商。因此 \(M\) 必为球面。
\end{solution}

\subsection{2018 年团体赛：几何与拓扑}

原题要求：从以下 6 道题中选择 5 道作答。

\begin{problem}
设 \(X=(S^2\vee S^2)\cup_fD^3\)，其中贴映射
\[
f:S^2\to S^2\vee S^2
\]
把连接南北极的一条大圆压缩。计算 \(X\) 的同调群。
\end{problem}

\begin{solution}
这与 2017 年团体赛第 1 题同型。把 \(S^2\) 沿一条大圆压缩后得到两个 \(2\)-球的楔和，贴映射在
\[
H_2(S^2)\to H_2(S^2\vee S^2)\cong\mathbb Z^2
\]
上为
\[
1\mapsto(1,1).
\]
因此胞腔链复形为
\[
0\to\mathbb Z\xrightarrow{(1,1)}\mathbb Z^2\to0\to\mathbb Z\to0.
\]
故
\[
H_0(X)=\mathbb Z,\qquad H_1(X)=0,\qquad H_2(X)=\mathbb Z,\qquad H_3(X)=0,
\]
其余同调群为零。
\end{solution}

\begin{problem}
\begin{enumerate}[label=(\alph*)]
\item 设 \(A\) 是 \(\R^3\) 中单个圆，计算 \(\pi_1(\R^3-A)\)。
\item 设 \(A,B\) 是 \(\R^3\) 中不交圆，分别位于上半空间和下半空间，计算 \(\pi_1(\R^3-(A\cup B))\)。
\end{enumerate}
\end{problem}

\begin{solution}
\begin{enumerate}[label=(\alph*)]
\item 单个圆的补空间同伦等价于 \(S^1\)，所以
\[
\pi_1(\R^3-A)\cong\mathbb Z.
\]
\item 两个圆位于不同半空间，构成平凡二分量 unlink。补空间基本群由两个 meridian 自由生成：
\[
\pi_1(\R^3-(A\cup B))\cong F_2.
\]
\end{enumerate}
\end{solution}

\begin{problem}
考虑 \(\R^3\) 中微分 \(1\)-形式
\[
\omega=x\,dy-y\,dx+dz.
\]
证明对任意处处非零函数 \(f:\R^3\to\R\)，\(f\omega\) 都不是闭形式。
\end{problem}

\begin{solution}
若 \(f\omega\) 闭，则
\[
0=d(f\omega)=df\wedge\omega+f\,d\omega.
\]
两边外乘 \(\omega\)，得
\[
0=f\,\omega\wedge d\omega.
\]
但
\[
d\omega=2dx\wedge dy,
\]
所以
\[
\omega\wedge d\omega
=2(xdy-ydx+dz)\wedge dx\wedge dy
=2dz\wedge dx\wedge dy,
\]
处处非零。因 \(f\) 处处非零，矛盾。因此 \(f\omega\) 不闭。
\end{solution}

\begin{problem}
证明
\[
Q^n=\left\{(x^1,\ldots,x^{n+1})\in\R^{n+1}:\sum_{i=1}^{n+1}(x^i)^4=1\right\}
\]
是可微流形。
\end{problem}

\begin{solution}
令
\[
F(x)=\sum_{i=1}^{n+1}(x^i)^4.
\]
则 \(Q^n=F^{-1}(1)\)。梯度为
\[
\nabla F=(4(x^1)^3,\ldots,4(x^{n+1})^3).
\]
在 \(F^{-1}(1)\) 上至少一个坐标非零，因此 \(\nabla F\ne0\)。故 \(1\) 是正则值，\(Q^n\) 是 \(\R^{n+1}\) 中光滑 \(n\) 维子流形。
\end{solution}

\begin{problem}
设 \(M\) 是 \(\R^3\) 中闭曲面。证明
\[
\int_M |K|\,d\sigma\ge4\pi(1+g),
\]
其中 \(K\) 是 Gaussian 曲率，\(g\) 是 genus，\(d\sigma\) 是面积元。
\end{problem}

\begin{solution}
这是曲面的 Chern--Lashof 全绝对曲率不等式。Gauss 映射
\[
N:M\to S^2
\]
的 Jacobian 为 \(K\)，所以
\[
\int_M|K|\,d\sigma
\]
等于 Gauss 映射的面积按重数积分。

对几乎所有方向 \(v\in S^2\)，高度函数
\[
h_v(x)=\langle x,v\rangle
\]
是 Morse 函数，其临界点正是 Gauss 映射取值为 \(\pm v\) 的点。Morse 不等式给出临界点数至少为 Betti 数之和：
\[
b_0+b_1+b_2=1+2g+1=2+2g.
\]
因此对几乎所有 \(v\)，\(N^{-1}(v)\) 与 \(N^{-1}(-v)\) 的总点数至少为 \(2+2g\)，平均到 \(S^2\) 上得到 Gauss 映射的总重数至少为 \(1+g\)。于是
\[
\int_M|K|\,d\sigma\ge(1+g)\operatorname{Area}(S^2)=4\pi(1+g).
\]
\end{solution}

\begin{problem}
设 \(M\) 是 \(n\) 维紧单连通 Riemann 流形。若
\[
\frac14<K_M\le1,
\]
证明 \(M\) 同胚于 \(S^n\)。
\end{problem}

\begin{solution}
由严格 \(1/4\)-pinched sphere theorem，紧单连通 Riemann 流形若截面曲率满足
\[
\frac14<K\le1,
\]
则 \(M\) 同胚于球面。更强的 Brendle--Schoen 微分球面定理给出 \(M\) 微分同胚于 spherical space form；在单连通条件下只能是 \(S^n\)。因此
\[
M\cong S^n.
\]
\end{solution}

\subsection{2019 年团体赛：几何与拓扑}

原题标注：团体赛共 5 题。

\begin{problem}
\(TS^2\) 是否微分同胚于 \(S^2\times\R^2\)？说明理由。这里 \(TS^2\) 是 \(S^2\) 的切丛总空间。
\end{problem}

\begin{solution}
不微分同胚。若
\[
TS^2\cong S^2\times\R^2
\]
作为向量丛总空间并保持零截面附近结构，则切丛将是平凡的。但
\[
e(TS^2)=\chi(S^2)=2\ne0
\]
在
\[
H^2(S^2;\mathbb Z)\cong\mathbb Z
\]
中非零，所以 \(TS^2\) 不是平凡向量丛。

更直接地，零截面 \(S^2\subset TS^2\) 的自交数等于 Euler 数
\[
\langle e(TS^2),[S^2]\rangle=2.
\]
而在 \(S^2\times\R^2\) 中，零截面 \(S^2\times\{0\}\) 的法丛平凡，自交数为 \(0\)。自交数为微分拓扑不变量，因此二者不微分同胚。
\end{solution}

\begin{problem}
求电影《美丽心灵》中 Russell Crowe 给学生的问题：
\[
V=\{F:\R^3\setminus X\to\R^3:\nabla\times F=0\},\qquad
W=\{F=\nabla g\},
\]
求
\[
\dim(V/W).
\]
先给出任意闭集 \(X\subset\R^3\) 的一般答案，再分别讨论
\[
(a)\ X=\{0\},\qquad (b)\ X=\{x=y=0\},\qquad (c)\ X=\{x=0\}.
\]
\end{problem}

\begin{solution}
把向量场 \(F\) 视为 \(1\)-形式
\[
\alpha=F_1dx+F_2dy+F_3dz.
\]
条件 \(\nabla\times F=0\) 等价于
\[
d\alpha=0,
\]
而 \(F=\nabla g\) 等价于
\[
\alpha=dg.
\]
因此
\[
V/W\cong H^1_{\mathrm{dR}}(\R^3\setminus X).
\]
一般答案为
\[
\dim(V/W)=b_1(\R^3\setminus X).
\]
若 \(X\) 足够好，例如局部可缩紧集，则由 Alexander 对偶
\[
\widetilde H^1(\R^3\setminus X)\cong \widetilde H_1(S^3\setminus X)\cong \widetilde H^{1}(X)
\]
可把答案转化为 \(X\) 的 Alexander 对偶数据。

(a) 若 \(X=\{0\}\)，则
\[
\R^3\setminus\{0\}\simeq S^2,
\]
故
\[
H^1=0,\qquad \dim(V/W)=0.
\]

(b) 若 \(X=\{x=y=0\}\) 是 \(z\)-轴，则
\[
\R^3\setminus X\simeq S^1,
\]
故
\[
H^1\cong\R,\qquad \dim(V/W)=1.
\]

(c) 若 \(X=\{x=0\}\) 是一个平面，则
\[
\R^3\setminus X
\]
有两个连通分支，每个都是半空间，均可缩。故每个分支上 \(H^1=0\)，整体
\[
H^1=0,\qquad \dim(V/W)=0.
\]
\end{solution}

\begin{problem}
设 \(T^2=S^1\times S^1\) 带标准定向，\(F:T^2\to T^2\) 是 degree 为 \(1\) 的光滑映射，满足
\[
F\circ F=\operatorname{Id}
\]
且 \(F\) 无不动点。证明诱导映射
\[
F^*:H^1(T^2)\to H^1(T^2)
\]
为恒等。
\end{problem}

\begin{solution}
设
\[
A=F^*:H^1(T^2;\mathbb Z)\to H^1(T^2;\mathbb Z).
\]
则 \(A\in GL(2,\mathbb Z)\)。由
\[
F^2=\operatorname{Id}
\]
得
\[
A^2=I.
\]
又 \(\deg F=1\)，所以 \(F\) 保定向，故
\[
\det A=1.
\]
\(2\times2\) 整矩阵满足 \(A^2=I\)、\(\det A=1\)，其特征值为 \(\pm1\)，乘积为 \(1\)，所以要么 \(A=I\)，要么 \(A=-I\)。

若 \(A=-I\)，Lefschetz 数为
\[
L(F)=\operatorname{tr}(F^*|H^0)-\operatorname{tr}(F^*|H^1)+\operatorname{tr}(F^*|H^2).
\]
这里 \(H^0\) 上迹为 \(1\)，\(H^1\) 上迹为 \(-2\)，\(H^2\) 上因 degree 为 \(1\) 迹为 \(1\)。故
\[
L(F)=1-(-2)+1=4\ne0.
\]
Lefschetz 不动点定理推出 \(F\) 有不动点，矛盾。因此
\[
A=I,
\]
即 \(F^*\) 在 \(H^1(T^2)\) 上为恒等。
\end{solution}

\begin{problem}
设 \(U(n)\) 为 \(n\times n\) 酉矩阵群，\(O(n)\) 为 \(n\times n\) 正交矩阵群，
\[
SU(n)=\{A\in U(n):\det A=1\},\qquad
SO(n)=\{A\in O(n):\det A=1\}.
\]
这些都是带自然流形结构的 Lie 群。
\begin{enumerate}[label=(\alph*)]
\item 计算 \(SU(n)\) 和 \(SO(n)\) 的维数。
\item 计算 \(SU(n)\) 和 \(SO(n)\) 的基本群，\(n\ge2\)。
\end{enumerate}
\end{problem}

\begin{solution}
\begin{enumerate}[label=(\alph*)]
\item Lie 代数为
\[
\mathfrak{su}(n)=\{X\in M_n(\C):X^*+X=0,\ \operatorname{tr}X=0\}.
\]
\(\mathfrak u(n)\) 的实维数为 \(n^2\)，迹为零去掉一个实维数，所以
\[
\dim SU(n)=n^2-1.
\]
而
\[
\mathfrak{so}(n)=\{X\in M_n(\R):X^T+X=0\},
\]
反对称实矩阵由上三角元素决定，因此
\[
\dim SO(n)=\frac{n(n-1)}2.
\]

\item \(SU(2)\cong S^3\)，故
\[
\pi_1(SU(2))=0.
\]
纤维丛
\[
SU(n-1)\to SU(n)\to S^{2n-1}
\]
给出归纳；因为 \(S^{2n-1}\) 对 \(n\ge3\) 单连通且 \(\pi_2(S^{2n-1})=0\)，得
\[
\pi_1(SU(n))=0,\qquad n\ge2.
\]

对 \(SO(n)\)，
\[
SO(2)\cong S^1,\qquad \pi_1(SO(2))\cong\mathbb Z.
\]
对 \(n\ge3\)，纤维丛
\[
SO(n)\to SO(n+1)\to S^n
\]
以及 \(SO(3)\cong\mathbb{RP}^3\) 给出
\[
\pi_1(SO(n))\cong\mathbb Z/2,\qquad n\ge3.
\]
\end{enumerate}
\end{solution}

\begin{problem}
设 \((M,g)\) 是紧 Riemann 流形，曲率张量为 \(R\)。称 \((M,g)\) weakly negative，如果对任意点 \(p\in M\) 和任意非零向量 \(X\in T_pM\)，存在非零 \(Y\in T_pM\)，使
\[
R(X,Y,Y,X)<0.
\]
\begin{enumerate}[label=(\alph*)]
\item 设 \(X\) 是 Killing 向量场，\(f=\frac12g(X,X)=\frac12|X|^2\)。证明对任意向量场 \(V\)，
\[
\operatorname{Hess}f(V,V)
=g(\nabla_VX,\nabla_VX)-R(V,X,X,V).
\]
\item 证明若 \((M,g)\) weakly negative，则不存在非平凡 Killing 向量场。
\end{enumerate}
\end{problem}

\begin{solution}
\begin{enumerate}[label=(\alph*)]
\item 因为
\[
f=\frac12\langle X,X\rangle,
\]
有
\[
Vf=\langle\nabla_VX,X\rangle.
\]
在一点取 \(\nabla_VV=0\) 的延拓，则
\[
\operatorname{Hess}f(V,V)
=V\langle\nabla_VX,X\rangle
=\langle\nabla_V\nabla_VX,X\rangle+\langle\nabla_VX,\nabla_VX\rangle.
\]
Killing 场满足恒等式
\[
\nabla_V\nabla_VX=-R(X,V)V
\]
在该点成立。因此
\[
\langle\nabla_V\nabla_VX,X\rangle
=-\langle R(X,V)V,X\rangle
=-R(X,V,V,X).
\]
按题目曲率排列，
\[
-R(X,V,V,X)=-R(V,X,X,V).
\]
所以
\[
\operatorname{Hess}f(V,V)
=|\nabla_VX|^2-R(V,X,X,V).
\]

\item 因 \(M\) 紧，\(f=\frac12|X|^2\) 取最大值于某点 \(p\)。若 \(X(p)\ne0\)，由 weak negativity，存在 \(Y\ne0\)，使
\[
R(X(p),Y,Y,X(p))<0.
\]
由 (a)，在最大点 Hessian 应半负定，但
\[
\operatorname{Hess}f(Y,Y)
=|\nabla_YX|^2-R(Y,X,X,Y)>0,
\]
矛盾。因此在最大点 \(X(p)=0\)，最大值为 \(0\)，故 \(X\equiv0\)。不存在非平凡 Killing 场。
\end{enumerate}
\end{solution}

\end{document}
